\documentclass[11pt]{ctexart}
\usepackage[a4paper,margin=1in]{geometry}
\usepackage{graphicx}
\usepackage{xcolor}
\usepackage{hyperref}
\definecolor{refgray}{gray}{0.45}
\title{\textbf{市场最新观点汇总}\\[0.4em]{\large 宏观利率分化、AI硬件从光学转向网络、大宗商品结构性缺口、中国消费与地产K型修复}}
\author{KC桌面 自动生成}
\date{260624}
\begin{document}
\maketitle
\tableofcontents
\newpage
\section{一页摘要}
\begin{itemize}
  \item 美联储新主席Warsh的“沉默策略”正重塑全球利率定价逻辑，市场波动率预期上升，新兴市场债券存在10-15个基点的定价错位。
  \item AI硬件投资正从光学板块的“需求发现”转向网络设备的“推理落地”，半导体板块拥挤度已达历史极值，但AI需求本身未见放缓。
  \item 中国新能源车6月订单脉冲由新车型和季末促销共同驱动，7月大概率出现“订单空窗期”，行业整合的三个条件目前一个都未满足。
  \item 中国化工行业呈现“传统品疲弱、氟化工结构性受益”的分化格局，AI算力正成为氟化工材料新的定价锚。
  \item 中国消费数据呈现“量稳价降”特征，人均旅游消费连续五个假期同比下滑，消费股的真正考验在于渠道碎片化时代的定价权。
  \item 日本股市结构性牛市进入第二幕，通胀稳定、政策加速、治理改革、居民入场四大驱动力同步强化。
  \item 黄金短期承压但长期逻辑未破，与2013年崩盘有本质区别；铝市场存在约200万吨的供给缺口，产能天花板担忧被夸大。
\end{itemize}
\section{宏观与利率：美联储“沉默”重塑全球定价，日本央行内部分歧加大}
\textbf{美联储主动退出预期管理，将定价权交还市场，导致利率波动率结构性上升；日本央行加息周期已过中点，但鹰派声音正在变清晰，财政政策成为下一个焦点。}

\begin{itemize}
  \item 美联储新主席Warsh有意减少前瞻指引，市场将其“不反驳”解读为默许，导致市场定价更紧密追随点阵图，投资者需从“猜美联储想什么”转向“猜市场如何猜美联储”。
  \item MS指出，美联储的沟通策略转变意味着利率波动率将更高，宏观市场的信号功能会更强，推荐SOFR曲线平坦化交易。
  \item 日本央行6月加息至1\%后，NOM认为加息周期可能已过中点，终点可能在2027年中后期，但两位委员已明确呼吁加速加息。
  \item GS分析日本央行6月会议意见摘要发现，支持加息的理由中“通胀超调风险”已取代“经济过热风险”成为核心变量，未来加息节奏将更依赖实际通胀数据。
  \item 日本市场关注重心正从央行转向政府，增长战略与消费税减税方案将成为决定日本资产定价的下一个变量，LDP提案计划2027年起将食品饮料税率降至1\%。
  \item NOM的人民币中间价模型显示，当前中间价设定存在系统性偏低，逆周期因子已显著收窄，央行干预力度减弱，市场化定价因子权重上升。
  \item 新兴市场债券收益率与油价出现10-15个基点的脱钩，MS认为这一定价错位将被修复，看好新兴市场货币对G3货币升值。
  \item JPM指出中国5月财政存款升至近十年同期第二高位，问题不在于财政空间不够，而在于部署节奏慢于预期，下半年政策空间反而更充足。
\end{itemize}
\begin{figure}[htbp]
\centering
\includegraphics[width=0.88\linewidth]{figures/fig_365.jpg}
\caption{Exhibit 1 - Exhibit 1: Yields have lagged the move in oil prices Exhibit 2: Even using a broader measure of energy prices, yields look 10-15bp too high There are a few reasons EM yields have not fallen further US yields have moved higher:}
\end{figure}
\begin{figure}[htbp]
\centering
\includegraphics[width=0.88\linewidth]{figures/fig_366.jpg}
\caption{Exhibit 1 - Exhibit 1: Yields have lagged the move in oil prices Exhibit 2: Even using a broader measure of energy prices, yields look 10-15bp too high There are a few reasons EM yields have not fallen further US yields have moved higher:}
\end{figure}
\begin{figure}[htbp]
\centering
\includegraphics[width=0.88\linewidth]{figures/fig_026.jpg}
\caption{Figure 1 - (a) With regard to the economy, the latest important development is that tensions in the Middle East have started to ease, with the US and Iran signing a memorandum of understanding on 18 June (US time). However, in order for global economic activity to normal}
\end{figure}
\begin{figure}[htbp]
\centering
\includegraphics[width=0.88\linewidth]{figures/fig_027.jpg}
\caption{Figure 2 - This suggests that it will take time for quantitative supply constraints, which have a direct bearing on the economy, to ease even as the situation in the Middle East eases. Looking ahead, we will need to keep a close eye on the extent of the recovery in petro}
\end{figure}
{\small\color{refgray}\textbf{References:} [R001] 摩根斯坦利：MS：新美联储的“沉默”本身就是最响亮的信号; [R005] NOM：日本加息周期已过中点，但市场真正该关注的是财政牌; [R037] GS：日本央行内部已有人要求加速加息，关键看通胀是否会超2\%; [R021] NOM：人民币中间价模型释放了一个被忽视的鸽派信号; [R044] 摩根斯坦利：新兴市场债券的“油价逻辑”尚未兑现; [R020] JPM：财政不是没子弹，是没扣扳机}

\section{AI/算力：从“光学狂欢”到“网络务实”，半导体拥挤度达历史极值}
\textbf{AI硬件投资逻辑正从光学转向网络与服务器，半导体板块估值和拥挤度均处于极端水平，但AI需求本身未见放缓，Agentic AI正在改变内存需求的量纲。}

\begin{itemize}
  \item MS指出，过去一个月网络设备板块上涨6.7\%而光学板块下跌10\%，市场关注重心正从光通信的“需求发现”转向网络设备的“推理落地”。
  \item BARC报告显示，AI硬件股的隐含PE倍数已从两个月前的20-40倍跃升至40-60倍，光学板块回调更多是技术性而非基本面恶化。
  \item Bernstein认为半导体板块唯一的游戏还在，但拥挤度已到历史极值，SOX指数年初至今上涨107\%，Forward EPS同步上涨75\%，估值扩张并非全部。
  \item BofA将2030年全球半导体TAM从2.3万亿美元上调至2.7万亿美元，AI产业核心矛盾已从“值不值得投”变成“能不能建得出来”。
  \item MS日本半导体设备报告指出，Agentic AI的工作负载模式将导致DRAM和NAND需求发生指数级跃迁，仅NVIDIA Vera CPU就可能使总DRAM bit需求增加约16\%。
  \item JEF供应链调研显示，AI服务器下一代背板架构Kyber大概率推迟到2028年甚至取消，2027年AI PCB市场规模将因此下调5\%-8\%。
  \item 中国大模型在Token消耗上已反超美国，JEF数据显示中国模型周消耗18.8万亿Token，美国模型仅5.8万亿，API成本仅为后者的几分之一。
  \item MS认为，连接性将成为下一个瓶颈，当芯片和内存密度足够高时，信号在铜缆上的传输距离本身就成了物理限制。
\end{itemize}
\begin{figure}[htbp]
\centering
\includegraphics[width=0.88\linewidth]{figures/fig_298.jpg}
\caption{Exhibit 27 - EXHIBIT 2: As AI has strengthened, “bottlenecks’ have taken off YTD Average Returns By Sub-Sector As of 6/22/2026 EXHIBIT 3: Semiconductor stocks have been outperforming the broader market meaningfully YTD As of 6/22/2026}
\end{figure}
\begin{figure}[htbp]
\centering
\includegraphics[width=0.88\linewidth]{figures/fig_299.jpg}
\caption{EXHIBIT 2 - EXHIBIT 4: SOX EPS is up \textbackslash{}\textasciitilde{}75\% YTD, accounting for the bulk of the stock performance over that period As of 6/22/2026}
\end{figure}
\begin{figure}[htbp]
\centering
\includegraphics[width=0.88\linewidth]{figures/fig_141.jpg}
\caption{Exhibit 8 - Exhibit 9: While each industry upcycle has differed in duration, unit vs. ASP contribution, and fundamental drivers, we observe that in most cases WFE intensity tends to rise during the period WFE performance during every semi indu}
\end{figure}
\begin{figure}[htbp]
\centering
\includegraphics[width=0.88\linewidth]{figures/fig_146.jpg}
\caption{Exhibit 22 - Exhibit 22: Industry HBM shipment could reach 16.190 CY30E, growing at +40\% CAGR between CY25-30E HBM Unit Shipment Forecast by Vendor (GB bn) 3 Exhibit 2 By Industry HBM sales could reach \textbackslash{}\$168bn by CY30E, growing at +37\% CAGR betw}
\end{figure}
\begin{figure}[htbp]
\centering
\includegraphics[width=0.88\linewidth]{figures/fig_241.jpg}
\caption{Exhibit 2 - Exhibit 1 - Weekly token consumption over the past 12 months as of 22 Jun (tn) Exhibit 2 - Silicon Data LLM Token Expenditure Index Exhibit 3 - Blended prices for China models vs US models (1Q26)}
\end{figure}
\begin{figure}[htbp]
\centering
\includegraphics[width=0.88\linewidth]{figures/fig_242.jpg}
\caption{Exhibit 1 - Exhibit 1 - Weekly token consumption over the past 12 months as of 22 Jun (tn) Exhibit 2 - Silicon Data LLM Token Expenditure Index Exhibit 3 - Blended prices for China models vs US models (1Q26) Exhibit 4 - Top 10 ranking in}
\end{figure}
{\small\color{refgray}\textbf{References:} [R042] 摩根斯坦利：MS：AI投资正从“光学狂欢”转向“网络务实”; [R025] BARC：AI硬件股的估值游戏正在从光学转向服务器; [R038] Bernstein：半导体唯一的游戏还在，但拥挤度已到历史极值; [R012] 美国银行：芯片行业下一个1万亿美元，只需五年; [R034] 摩根斯坦利：半导体设备的下一个瓶颈不是芯片，是连接; [R040] JEF：AI服务器PCB的“Kyber延迟”比想象中更值得警惕; [R029] JEF：中国大模型正在用“成本-智能”剪刀差改写全球AI竞争格局}

\section{能源与大宗：铝缺口被低估，黄金与2013年本质不同，航运逻辑切换}
\textbf{铝市场存在约200万吨供给缺口，产能天花板担忧被夸大；黄金短期承压但长期逻辑未破；航运市场的真正驱动力正从“绕航”切换为“抢运”。}

\begin{itemize}
  \item Citi认为市场对电解铝48百万吨产能上限的恐慌源于被误读的统计局数据，5月年化产量已回落至45.8百万吨，产能管控政策没有松动。
  \item Citi预计2026年全球铝市场存在约200万吨的供给缺口，2027年缺口收窄至27万吨，三季度铝价预计维持在4000美元/吨左右。
  \item UBS指出黄金与2013年崩盘有本质不同：当前通胀在上升而非下降，美联储政策方向仍偏宽松，央行购金趋势未变只是买家结构在变化。
  \item UBS认为黄金短期风险确实在上升，但长期逻辑未被打破，美国债务/GDP超137\%提供结构性支撑，亚洲占实物投资需求超70\%。
  \item UBS航运报告指出，VLCC运价飙升是事件性的，真正值得关注的是亚洲至美国航线的集装箱“前置抢运”正在加速，SCFI综合指数环比+5\%、同比+67\%。
  \item 霍尔木兹海峡恢复通航后，可运出的铝库存不到40万吨，远低于市场预期，印尼新增产能受电力建设制约，2026年海外实际产量同比减少160万吨。
  \item 中国MDI出口正从“以量换价”转向“以价换量”，5月出口均价2193美元/吨创近两年新高，欧洲正取代美国成为第一增量市场。
\end{itemize}
\begin{figure}[htbp]
\centering
\includegraphics[width=0.88\linewidth]{figures/fig_015.jpg}
\caption{Figure 1 - Figure 1: Middle East/US Gulf/West Africa-China VLCC TCE +122\%/+6\%/-20\% vs. pre-conflict level Figure 3: Average earnings of Capesize rallied by 11\% WoW last week}
\end{figure}
\begin{figure}[htbp]
\centering
\includegraphics[width=0.88\linewidth]{figures/fig_018.jpg}
\caption{Figure 3 - Figure 6: Container volume new in transit from Asia to US has been showing a front loading pattern}
\end{figure}
\begin{figure}[htbp]
\centering
\includegraphics[width=0.88\linewidth]{figures/fig_019.jpg}
\caption{Figure 4 - Figure 7: SCFI +5\% WoW and +67\% YoY}
\end{figure}
\begin{figure}[htbp]
\centering
\includegraphics[width=0.88\linewidth]{figures/fig_067.jpg}
\caption{Exhibit 4 - Exhibit 4: KEI Industries Q1FY27E expectations Exhibit 5: Copper on an average has increased 1-5\% MoM every month in Q1 Exhibit 6: Aluminium also increased in April and May MoM, but has corrected \textbackslash{}\textasciitilde{}1\% in June YTD Exhibit 7: On}
\end{figure}
\begin{figure}[htbp]
\centering
\includegraphics[width=0.88\linewidth]{figures/fig_068.jpg}
\caption{Exhibit 5 - Exhibit 5: Copper on an average has increased 1-5\% MoM every month in Q1 Exhibit 6: Aluminium also increased in April and May MoM, but has corrected \textbackslash{}\textasciitilde{}1\% in June YTD Exhibit 7: On a YoY basis, we are seeing the largest quantum}
\end{figure}
{\small\color{refgray}\textbf{References:} [R036] Citi：铝价缺口比市场想象的更深，中国铝业和中国宏桥的买入窗口已打开; [R006] UBS：黄金正处在“无人区”，但这次和2013年不一样; [R004] UBS：霍尔木兹海峡的VLCC运价已不是关键，集装箱的“前置抢运”才是; [R022] JPM：中国MDI出口量价齐升，欧洲市场正成为新的定价锚}

\section{中国新能源车：6月订单透支7月，行业整合远未到终局}
\textbf{6月新能源车订单脉冲由新车型和季末促销共同驱动，7月大概率出现“订单空窗期”；行业整合的三个条件目前一个都未满足，出口是唯一亮点。}

\begin{itemize}
  \item MS周度数据显示6月15-21日主要品牌订单环比飙升34\%，但警告这种订单前置正在透支7月需求，形成典型的“订单空气包”。
  \item 比亚迪周订单环比暴增61\%至10.45万辆，主要靠唐系列新车型预订单转化和季度末促销，GS数据显示新能源平均折扣7.48\%微幅收窄。
  \item GS提出行业整合三条件框架：多数车企EBITDA为负、管理层全面投降、资产负债表普遍净债务，目前14家样本车企中仍有10家保持正EBITDA。
  \item GS预计2026年NEV渗透率达60\%，但国内NEV零售仅增长1\%，出口同比增长70\%是唯一亮点，比亚迪、零跑、小鹏在海外布局上占优。
  \item MS认为欧盟将反补贴关税扩展至PHEV的担忧被过度定价，股价跌幅与南向资金持仓高度相关而非PHEV敞口，车企的适应能力比市场预期的强。
  \item GS数据显示6月前两周新能源零售渗透率63.9\%，批发渗透率67.9\%，均高于5月全月水平，但碳酸锂价格继续回落至16.75万元/吨。
  \item MS指出，真正依靠产品力实现“自然转化”的品牌与依赖促销“催熟”的品牌正在走向分化，新车型驱动的转化可持续，促销驱动的将在7月暴露真实需求底色。
\end{itemize}
\begin{figure}[htbp]
\centering
\includegraphics[width=0.88\linewidth]{figures/fig_153.jpg}
\caption{Exhibit 2 - Exhibit 2: NEV dealer discount trend By brand dealer discount is based on best-selling models. Exhibit 3: ICE dealer discount trend}
\end{figure}
\begin{figure}[htbp]
\centering
\includegraphics[width=0.88\linewidth]{figures/fig_154.jpg}
\caption{Exhibit 2 - Exhibit 2: NEV dealer discount trend By brand dealer discount is based on best-selling models. Exhibit 3: ICE dealer discount trend By brand dealer discount is based on best-selling models.}
\end{figure}
\begin{figure}[htbp]
\centering
\includegraphics[width=0.88\linewidth]{figures/fig_070.jpg}
\caption{Exhibit 1 - Exhibit 2: Aggregate 2026E/2027E EBITDA consensus estimates for China OEMs have undergone four consecutive rounds of cuts, and momentum has been accelerating \#\# Inflection monitor - 2025/1Q26 Bottom line: Based on our three wat}
\end{figure}
\begin{figure}[htbp]
\centering
\includegraphics[width=0.88\linewidth]{figures/fig_071.jpg}
\caption{Exhibit 2 - Exhibit 2: Aggregate 2026E/2027E EBITDA consensus estimates for China OEMs have undergone four consecutive rounds of cuts, and momentum has been accelerating \#\# Inflection monitor - 2025/1Q26 Bottom line: Based on our three wat}
\end{figure}
{\small\color{refgray}\textbf{References:} [R033] 摩根斯坦利：MS：6月订单透支7月，新能源车市的分水岭不在销量，在“催熟”与“真转化”; [R013] GS：新能源车周订单暴增34\%，但真正的胜负手不在销量; [R010] GS：中国车企的“淘汰赛”还远未到终局; [R031] 摩根斯坦利：MS：PHEV关税担忧被过度定价，中国车企的缓冲比市场想象的厚}

\section{中国化工：传统品疲弱，氟化工与MDI的结构性机会}
\textbf{传统化工品陷入库存与开工率双重压力，但AI算力驱动的含氟化学品和MDI出口正打开被低估的定价空间，市场尚未充分定价这一结构性变化。}

\begin{itemize}
  \item UBS数据显示中国地方炼厂开工率已跌至45.74\%，原油库存自6月初持续下降，但主营炼厂开工率相对稳定，问题在结构而非总量。
  \item 聚酯链压力最大，涤纶长丝开工率74\%同比下滑15个百分点，PTA开工率69\%同比下滑13个百分点，但聚氨酯有亮点，TDI开工率环比上升7个百分点至71\%。
  \item UBS明确指出，氟化工材料在AI领域的应用潜力正在快速增长，市场尚未完全定价AI带给氟化工材料的增长机会，当前估值相比电子化学品存在显著折价。
  \item AI服务器对散热效率和低介电常数提出远超传统服务器的要求，氟化液和PTFE成为液冷和PCB基板的理想材料，国产替代机会明确。
  \item JPM数据显示5月中国MDI出口量10.1万吨同比飙升49\%，出口均价2193美元/吨创近两年新高，欧洲正取代美国成为第一增量市场。
  \item 中国钛白粉出口呈现结构性变化，氯化法出口增速是硫酸法的6.5倍，前5个月氯化法出口同比增长39\%，中国企业正从中低端向中高端市场渗透。
  \item 中国锂贸易已发生结构性逆转，前5个月氢氧化锂从净出口35.4千吨转为净进口10.5千吨，碳酸锂进口依赖度加深，进口均价同比上涨102\%。
\end{itemize}
\begin{figure}[htbp]
\centering
\includegraphics[width=0.88\linewidth]{figures/fig_047.jpg}
\caption{Figure 1 - Figure 1: China crude oil inventory Figure 2: Utilisation of China SOE/teapot refiners Chemical utilisation Figure 3: China ethylene (naphtha cracking) capacity utilisation}
\end{figure}
\begin{figure}[htbp]
\centering
\includegraphics[width=0.88\linewidth]{figures/fig_048.jpg}
\caption{Figure 1 - Figure 1: China crude oil inventory Figure 2: Utilisation of China SOE/teapot refiners Chemical utilisation Figure 3: China ethylene (naphtha cracking) capacity utilisation Figure 4: China ethylene (MTO) capacity utilisation}
\end{figure}
\begin{figure}[htbp]
\centering
\includegraphics[width=0.88\linewidth]{figures/fig_177.jpg}
\caption{Figure 2 - Figure 3: China's electronic-grade hydrofluoric acid consumption and revenue}
\end{figure}
\begin{figure}[htbp]
\centering
\includegraphics[width=0.88\linewidth]{figures/fig_178.jpg}
\caption{Figure 3 - Figure 4: UPSSS-grade EG-HF prices are higher than those of other grades}
\end{figure}
{\small\color{refgray}\textbf{References:} [R008] UBS：当炼厂开工率跌至45\%，化工行业下一个机会在AI制冷剂; [R018] UBS：氟化工正被AI重新定价，但市场还没看懂; [R022] JPM：中国MDI出口量价齐升，欧洲市场正成为新的定价锚; [R024] JPM：中国钛白粉出口的“量价齐升”并非故事的全部; [R023] JPM：中国锂贸易已发生结构性逆转，净进口时代开启}

\section{中国消费与地产：K型分化加速，定价权成为核心考验}
\textbf{消费数据呈现“量稳价降”特征，人均旅游消费连续五个假期下滑；地产市场修复脆弱但一线城市出现积极信号，消费股的真正考验在于渠道碎片化时代的定价权。}

\begin{itemize}
  \item Citi数据显示端午假期国内旅游人次同比增长4.4\%，但人均消费同比下降0.4\%，且从元旦到端午连续五个假期人均消费同比增速均为负值。
  \item NOM指出端午消费数据是消费动能系统性回落的第三个信号，增速从春节的8.2\%降至五一的3.5\%再到端午的接近0\%，天气只是加速了已存在的放缓过程。
  \item MS对中国消费品公司2026/27年净利润预测大幅下调，中顺洁柔被砍20\%/27\%，上海家化被砍35\%/39\%，核心问题是线上渠道侵蚀平均售价。
  \item JPM认为地产股急跌是四重噪音叠加而非基本面逆转，一线城市二手挂牌量持续下降和房价环比连续第四个月转正这两个信号被市场低估。
  \item MS周度数据显示50城新房成交量同比下降16\%，端午节假期带来日历效应，但剔除后市场修复仍是“低位企稳”而非“趋势性回暖”。
  \item Citi认为OTA行业正经历“消费降级式繁荣”，人均消费持续下行但出游人次保持正增长，平台竞争逻辑从“卖更贵的”转向“卖更多的低价短途产品”。
  \item JPM指出香港楼市最大的风险不是加息而是恒指，历史数据显示房价滞后股市3-6个月，恒指从高点下跌约15\%的影响可能才刚刚开始被计入。
\end{itemize}
\begin{figure}[htbp]
\centering
\includegraphics[width=0.88\linewidth]{figures/fig_206.jpg}
\caption{Figure 1 - Brian Cho Figure 1. Golden Week Domestic Tourist Number \& Revenue as \% of 2019 © 2026 Citi Inc. No redistribution without Citi's written permission. Figure 2. Golden Week Domestic Tourist Numbers © 2026 Citi Inc. No redistribution without Citi's written permis}
\end{figure}
\begin{figure}[htbp]
\centering
\includegraphics[width=0.88\linewidth]{figures/fig_207.jpg}
\caption{Figure 1 - Figure 1. Golden Week Domestic Tourist Number \& Revenue as \% of 2019 © 2026 Citi Inc. No redistribution without Citi's written permission. Figure 2. Golden Week Domestic Tourist Numbers © 2026 Citi Inc. No redistribution without Citi's written permission. Note}
\end{figure}
\begin{figure}[htbp]
\centering
\includegraphics[width=0.88\linewidth]{figures/fig_159.jpg}
\caption{Figure 1 - Figure 2: Iceberg Index - tier-1 cities' secondary listing volume (冰山指数二手挂牌量) No. of secondary listings ('000 units) - Tier 1 cities}
\end{figure}
\begin{figure}[htbp]
\centering
\includegraphics[width=0.88\linewidth]{figures/fig_160.jpg}
\caption{Figure 1 - Figure 2: Iceberg Index - tier-1 cities' secondary listing volume (冰山指数二手挂牌量) No. of secondary listings ('000 units) - Tier 1 cities \#\# Official sales registrations (lag of a few weeks) Figure 3: 60-city weekly primary sales reg}
\end{figure}
\begin{figure}[htbp]
\centering
\includegraphics[width=0.88\linewidth]{figures/fig_211.jpg}
\caption{Figure 1 - Figure 1: Hong Kong secondary home price index (and major events) \textbackslash{}- Transaction volume: Remained strong year-to-date; the latest 12-month rolling Y/Y growth is tracking 40\%. Figure 2: HK monthly private residential transactions}
\end{figure}
\begin{figure}[htbp]
\centering
\includegraphics[width=0.88\linewidth]{figures/fig_212.jpg}
\caption{Figure 1 - Figure 1: Hong Kong secondary home price index (and major events) \textbackslash{}- Transaction volume: Remained strong year-to-date; the latest 12-month rolling Y/Y growth is tracking 40\%. Figure 2: HK monthly private residential transactions}
\end{figure}
{\small\color{refgray}\textbf{References:} [R026] Citi：端午出游数据背后，中国消费韧性比想象中更“挑”; [R027] Citi：端午节旅游数据背后的真正信号——OTA行业正在经历一场“消费降级式繁荣”; [R035] NOM：端午消费数据透露的信号，比表面数字更值得警惕; [R030] 摩根斯坦利：MS：中国消费股的真正考验不是需求，是定价权; [R014] JPM：地产股急跌不是基本面反转，而是噪音窗口; [R016] 摩根斯坦利：MS：中国房地产的修复，还差一个“端午效应”的量级; [R028] JPM：香港楼市最大的风险不是加息，而是恒指}

\section{日本市场：结构性牛市第二幕，四大驱动力同步强化}
\textbf{日本股市正在从“低通胀、低利润”的旧范式切换至新定价框架，通胀稳定、政策加速、治理改革、居民入场四大驱动力同步强化，日元不再是企业盈利的唯一驱动因素。}

\begin{itemize}
  \item MS对TOPIX基准目标价为2027年6月底4300点，核心逻辑是日本正进入一个“更稳定、更可持续”的通胀环境，企业定价能力增强。
  \item 日本企业自疫情后持续改善利润率，美元兑日元与日本股市的正相关性在过去十年中持续下降，日元不再是日本企业盈利的唯一驱动因素。
  \item 执政联盟大胜后政策可预测性提高，2022年东证市场改革后Prime市场PBR已涨约40\%，2026年预计有治理准则修订推动企业更有效使用现金。
  \item NISA免税账户推动日本家庭从储蓄转向股票，2027年还计划推出“Kids NISA”，未成年人每年可免税投资60万日元。
  \item NOM认为日本央行加息周期可能已过中点，但财政政策的不确定性比加息本身更值得关注，消费税减税方案的资金来源与推进节奏将深刻影响内需和国债定价。
  \item GS分析日本央行6月会议意见摘要发现，有两位委员明确主张加速加息，认为政策利率应尽快接近中性利率水平，但这一主张尚未扩散为央行共识。
  \item MS日本电子化学品报告指出，FC-BGA覆铜板正经历“非常强劲”的需求周期，Resonac可能启动产能扩张并进一步提价，材料供应商的议价权正在重塑。
\end{itemize}
\begin{figure}[htbp]
\centering
\includegraphics[width=0.88\linewidth]{figures/fig_026.jpg}
\caption{Figure 1 - (a) With regard to the economy, the latest important development is that tensions in the Middle East have started to ease, with the US and Iran signing a memorandum of understanding on 18 June (US time). However, in order for global economic activity to normal}
\end{figure}
\begin{figure}[htbp]
\centering
\includegraphics[width=0.88\linewidth]{figures/fig_027.jpg}
\caption{Figure 2 - This suggests that it will take time for quantitative supply constraints, which have a direct bearing on the economy, to ease even as the situation in the Middle East eases. Looking ahead, we will need to keep a close eye on the extent of the recovery in petro}
\end{figure}
\begin{figure}[htbp]
\centering
\includegraphics[width=0.88\linewidth]{figures/fig_028.jpg}
\caption{Figure 4 - Fig. 2: Sharp decline in imports of oil and petroleum products Fig. 3: Delays in delivery of raw materials and parts \#\# Prices: We expect CPI inflation to peak in 2027 Q1 The second factor that will need to be monitored to see whether the BOJ is likely to cont}
\end{figure}
{\small\color{refgray}\textbf{References:} [R003] 摩根斯坦利：MS：日本股市的结构性牛市才刚刚进入第二幕; [R005] NOM：日本加息周期已过中点，但市场真正该关注的是财政牌; [R037] GS：日本央行内部已有人要求加速加息，关键看通胀是否会超2\%; [R041] 摩根斯坦利：日本电子化学品正在经历一场“材料级”的供应链重构}

\section{区域市场：亚洲太空经济被低估，印度制造业成本驱动型修复}
\textbf{亚洲太空经济股票相对全球同行存在60\%的市盈率折价，市场定价仍停留在“概念主题”阶段；印度制造业正从“量价齐升”切换至“价升量平、利润承压”的新区间。}

\begin{itemize}
  \item GS报告指出亚洲太空经济股票市盈率折价60\%、市净率折价25\%，处于历史低位附近，而全球太空主题基金AUM已从10亿美元飙升至约250亿美元。
  \item 2025年全球发射卫星4400+颗同比增65\%，未来5年还有约7万颗卫星计划发射，LEO卫星市场预计到2035年增长7倍至1080亿美元。
  \item 亚洲是全球太空供应链的硬件脊梁，中国、日本、韩国、台湾、印度各有分工，且都有政府订单支撑，AI+太空正在融合推动“轨道AI计算”进入采购阶段。
  \item GS数据显示印度电缆和EMS公司Q1收入增长主要驱动力是铜铝价格同比涨幅超35\%带来的价格传导与渠道补库存，而非需求爆发。
  \item 印度电缆公司Polycab预计Q1营收同比+36\%，但EBITDA利润率从14.5\%降至13.4\%；EMS公司Dixon营收同比+8\%但利润率从3.8\%降至3.4\%。
  \item JPM指出可持续基金资金流出已停止但大规模流入前提尚未满足，全球可持续基金低配科技股的结构性弱点正在被AI主题放大。
  \item MS认为中国生物医药创新正从“工程能力验证”进入“生物学差异化”第二阶段，全球药企开始认可中国从“快速跟进”向差异化生物学转变。
\end{itemize}
\begin{figure}[htbp]
\centering
\includegraphics[width=0.88\linewidth]{figures/fig_005.jpg}
\caption{Exhibit 6 - Exhibit 6: Space-themed funds have seen meaningful inflows since 2015, with aggregate AUM reaching US\textbackslash{}\$24bn across 40+ funds globally Exhibit 7: Asia supply chain companies remain materially underrepresented in global thematic al}
\end{figure}
\begin{figure}[htbp]
\centering
\includegraphics[width=0.88\linewidth]{figures/fig_006.jpg}
\caption{Exhibit 6 - Exhibit 6: Space-themed funds have seen meaningful inflows since 2015, with aggregate AUM reaching US\textbackslash{}\$24bn across 40+ funds globally Exhibit 7: Asia supply chain companies remain materially underrepresented in global thematic al}
\end{figure}
\begin{figure}[htbp]
\centering
\includegraphics[width=0.88\linewidth]{figures/fig_007.jpg}
\caption{Exhibit 8 - Exhibit 8: Asia Space Economy stocks have outperformed the broader regional equity market since 2025, led by Korea and Taiwan suppliers, although performance has recently consolidated following a strong rally in late 2025 Exhibit}
\end{figure}
\begin{figure}[htbp]
\centering
\includegraphics[width=0.88\linewidth]{figures/fig_067.jpg}
\caption{Exhibit 4 - Exhibit 4: KEI Industries Q1FY27E expectations Exhibit 5: Copper on an average has increased 1-5\% MoM every month in Q1 Exhibit 6: Aluminium also increased in April and May MoM, but has corrected \textbackslash{}\textasciitilde{}1\% in June YTD Exhibit 7: On}
\end{figure}
\begin{figure}[htbp]
\centering
\includegraphics[width=0.88\linewidth]{figures/fig_068.jpg}
\caption{Exhibit 5 - Exhibit 5: Copper on an average has increased 1-5\% MoM every month in Q1 Exhibit 6: Aluminium also increased in April and May MoM, but has corrected \textbackslash{}\textasciitilde{}1\% in June YTD Exhibit 7: On a YoY basis, we are seeing the largest quantum}
\end{figure}
{\small\color{refgray}\textbf{References:} [R002] GS：亚洲太空经济正在从“故事”转向“订单”，但市场定价仍落后60\%; [R009] GS：印度电缆与EMS公司正经历一轮“成本驱动型”收入修复; [R011] JPM：可持续投资正在“喘口气”，但真正的风险来自科技; [R032] 摩根斯坦利：MS：中国生物医药的“全球叙事”正在换剧本}

\section{风险偏好：AI主题拥挤但需求无虞，可持续投资“喘口气”}
\textbf{AI主题已成为半导体板块唯一的驱动力，拥挤度达历史极值但需求未见放缓；可持续投资资金流出已止血，但科技股主导市场放大了其结构性弱点。}

\begin{itemize}
  \item Bernstein指出当前半导体行情的驱动力已从基本面改善切换至“只有这一个游戏可玩”的叙事惯性，SOX指数年初至今上涨107\%，是标普500涨幅的11倍。
  \item Bernstein的Tearsheet显示7个关键指标中6个已从三个月前的“中性”转为“负面”，估值、拥挤度、库存天数均在警示风险。
  \item JPM可持续投资报告显示全球可持续基金总资产3.8万亿美元，但市场份额从高点持续下滑至5.6\%，资金流从净流出转向中性。
  \item JPM认为可持续基金要实现大规模资金流入的前提是相对业绩持续改善，但目前欧洲可持续基金小幅跑赢、美国小幅跑输，不足以吸引增量资金。
  \item JPM指出可持续基金普遍低配科技股，而AI主题恰恰集中在美股，这一结构性弱点正在被放大，ESG基金的IT低配是最大的跟踪误差来源。
  \item Bernstein认为车企造人形机器人不是跨界而是“降维复用”，电机、减速器、传感器与新能源汽车电驱系统高度同源，现有生产线可直接迁移。
  \item MS将2026年中国人形机器人出货量预测从此前的2.8万台大幅上调至5万台，商业化验证、政策支持与供应链反馈正同步加速。
\end{itemize}
\begin{figure}[htbp]
\centering
\includegraphics[width=0.88\linewidth]{figures/fig_298.jpg}
\caption{Exhibit 27 - EXHIBIT 2: As AI has strengthened, “bottlenecks’ have taken off YTD Average Returns By Sub-Sector As of 6/22/2026 EXHIBIT 3: Semiconductor stocks have been outperforming the broader market meaningfully YTD As of 6/22/2026}
\end{figure}
\begin{figure}[htbp]
\centering
\includegraphics[width=0.88\linewidth]{figures/fig_299.jpg}
\caption{EXHIBIT 2 - EXHIBIT 4: SOX EPS is up \textbackslash{}\textasciitilde{}75\% YTD, accounting for the bulk of the stock performance over that period As of 6/22/2026}
\end{figure}
\begin{figure}[htbp]
\centering
\includegraphics[width=0.88\linewidth]{figures/fig_098.jpg}
\caption{Figure 1 - Figure 1: Sustainable AuM has stabilized YTD... Figure 2: Sustainable AuM market share Sustainable AuM market share by region (\%) Figure 3: European funds account for the vast majority of sustainable assets AuM of Sustainable}
\end{figure}
\begin{figure}[htbp]
\centering
\includegraphics[width=0.88\linewidth]{figures/fig_099.jpg}
\caption{Figure 1 - Figure 4: Equity funds account for roughly two-thirds of sustainable funds AuM of Sustainable Funds, by asset class and region (\%)}
\end{figure}
\begin{figure}[htbp]
\centering
\includegraphics[width=0.88\linewidth]{figures/fig_166.jpg}
\caption{Exhibit 2 - Exhibit 2: China humanoid development has seen rapid developments in 2024-26; we are in the early commercialization stage. \#\# China Humanoid Market Forecast Stronger push toward commercialization. As highlighted in our 2026 outlo}
\end{figure}
\begin{figure}[htbp]
\centering
\includegraphics[width=0.88\linewidth]{figures/fig_167.jpg}
\caption{Exhibit 3 - Exhibit 3: We raise our China 2026 humanoid forecast to 50k unit and expect full-size humanoids to see more adoption with improving abilities. China Humanoid Sales ('k unit) Exhibit 4: We expect the industry's blended ASP to decl}
\end{figure}
{\small\color{refgray}\textbf{References:} [R038] Bernstein：半导体唯一的游戏还在，但拥挤度已到历史极值; [R011] JPM：可持续投资正在“喘口气”，但真正的风险来自科技; [R039] Bernstein：车企造人形机器人，真正的护城河不在机器人本身; [R015] 摩根斯坦利：人形机器人2026年出货量上调近八成，关键不在技术而在“部署数据闭环”}

\section{结语}
综合来看，当前市场呈现高度分化的格局：宏观利率层面美联储的“沉默”正在重塑全球定价规则，AI硬件投资从光学转向网络与服务器，大宗商品中铝和氟化工存在结构性缺口，而中国消费与地产的修复仍显脆弱。投资者需要从“猜央行意图”转向“猜市场如何猜央行”，在AI主题拥挤度达历史极值的背景下，关注真正具备盈利兑现能力的细分领域。
\newpage
\section{报告覆盖清单}
\begin{itemize}
  \item [R001] 宏观与利率：美联储“沉默”重塑全球定价，日本央行内部分歧加大 - 摩根斯坦利：MS：新美联储的“沉默”本身就是最响亮的信号
  \item [R002] 区域市场：亚洲太空经济被低估，印度制造业成本驱动型修复 - GS：亚洲太空经济正在从“故事”转向“订单”，但市场定价仍落后60\%
  \item [R003] 日本市场：结构性牛市第二幕，四大驱动力同步强化 - 摩根斯坦利：MS：日本股市的结构性牛市才刚刚进入第二幕
  \item [R004] 能源与大宗：铝缺口被低估，黄金与2013年本质不同，航运逻辑切换 - UBS：霍尔木兹海峡的VLCC运价已不是关键，集装箱的“前置抢运”才是
  \item [R005] 宏观与利率：美联储“沉默”重塑全球定价，日本央行内部分歧加大 / 日本市场：结构性牛市第二幕，四大驱动力同步强化 - NOM：日本加息周期已过中点，但市场真正该关注的是财政牌
  \item [R006] 能源与大宗：铝缺口被低估，黄金与2013年本质不同，航运逻辑切换 - UBS：黄金正处在“无人区”，但这次和2013年不一样
  \item [R007] 未被主题摘要引用 - Bernstein：数据中心每兆瓦值多少钱，答案藏在商业模式里
  \item [R008] 中国化工：传统品疲弱，氟化工与MDI的结构性机会 - UBS：当炼厂开工率跌至45\%，化工行业下一个机会在AI制冷剂
  \item [R009] 区域市场：亚洲太空经济被低估，印度制造业成本驱动型修复 - GS：印度电缆与EMS公司正经历一轮“成本驱动型”收入修复
  \item [R010] 中国新能源车：6月订单透支7月，行业整合远未到终局 - GS：中国车企的“淘汰赛”还远未到终局
  \item [R011] 区域市场：亚洲太空经济被低估，印度制造业成本驱动型修复 / 风险偏好：AI主题拥挤但需求无虞，可持续投资“喘口气” - JPM：可持续投资正在“喘口气”，但真正的风险来自科技
  \item [R012] AI/算力：从“光学狂欢”到“网络务实”，半导体拥挤度达历史极值 - 美国银行：芯片行业下一个1万亿美元，只需五年
  \item [R013] 中国新能源车：6月订单透支7月，行业整合远未到终局 - GS：新能源车周订单暴增34\%，但真正的胜负手不在销量
  \item [R014] 中国消费与地产：K型分化加速，定价权成为核心考验 - JPM：地产股急跌不是基本面反转，而是噪音窗口
  \item [R015] 风险偏好：AI主题拥挤但需求无虞，可持续投资“喘口气” - 摩根斯坦利：人形机器人2026年出货量上调近八成，关键不在技术而在“部署数据闭环”
  \item [R016] 中国消费与地产：K型分化加速，定价权成为核心考验 - 摩根斯坦利：MS：中国房地产的修复，还差一个“端午效应”的量级
  \item [R017] 未被主题摘要引用 - UBS：中国券商股不只是反弹，是估值体系正在被政策重写
  \item [R018] 中国化工：传统品疲弱，氟化工与MDI的结构性机会 - UBS：氟化工正被AI重新定价，但市场还没看懂
  \item [R019] 未被主题摘要引用 - GS：中国经济5月边际回暖，但内需疲软与信用紧缩的“温差”才是真正挑战
  \item [R020] 宏观与利率：美联储“沉默”重塑全球定价，日本央行内部分歧加大 - JPM：财政不是没子弹，是没扣扳机
  \item [R021] 宏观与利率：美联储“沉默”重塑全球定价，日本央行内部分歧加大 - NOM：人民币中间价模型释放了一个被忽视的鸽派信号
  \item [R022] 能源与大宗：铝缺口被低估，黄金与2013年本质不同，航运逻辑切换 / 中国化工：传统品疲弱，氟化工与MDI的结构性机会 - JPM：中国MDI出口量价齐升，欧洲市场正成为新的定价锚
  \item [R023] 中国化工：传统品疲弱，氟化工与MDI的结构性机会 - JPM：中国锂贸易已发生结构性逆转，净进口时代开启
  \item [R024] 中国化工：传统品疲弱，氟化工与MDI的结构性机会 - JPM：中国钛白粉出口的“量价齐升”并非故事的全部
  \item [R025] AI/算力：从“光学狂欢”到“网络务实”，半导体拥挤度达历史极值 - BARC：AI硬件股的估值游戏正在从光学转向服务器
  \item [R026] 中国消费与地产：K型分化加速，定价权成为核心考验 - Citi：端午出游数据背后，中国消费韧性比想象中更“挑”
  \item [R027] 中国消费与地产：K型分化加速，定价权成为核心考验 - Citi：端午节旅游数据背后的真正信号——OTA行业正在经历一场“消费降级式繁荣”
  \item [R028] 中国消费与地产：K型分化加速，定价权成为核心考验 - JPM：香港楼市最大的风险不是加息，而是恒指
  \item [R029] AI/算力：从“光学狂欢”到“网络务实”，半导体拥挤度达历史极值 - JEF：中国大模型正在用“成本-智能”剪刀差改写全球AI竞争格局
  \item [R030] 中国消费与地产：K型分化加速，定价权成为核心考验 - 摩根斯坦利：MS：中国消费股的真正考验不是需求，是定价权
  \item [R031] 中国新能源车：6月订单透支7月，行业整合远未到终局 - 摩根斯坦利：MS：PHEV关税担忧被过度定价，中国车企的缓冲比市场想象的厚
  \item [R032] 区域市场：亚洲太空经济被低估，印度制造业成本驱动型修复 - 摩根斯坦利：MS：中国生物医药的“全球叙事”正在换剧本
  \item [R033] 中国新能源车：6月订单透支7月，行业整合远未到终局 - 摩根斯坦利：MS：6月订单透支7月，新能源车市的分水岭不在销量，在“催熟”与“真转化”
  \item [R034] AI/算力：从“光学狂欢”到“网络务实”，半导体拥挤度达历史极值 - 摩根斯坦利：半导体设备的下一个瓶颈不是芯片，是连接
  \item [R035] 中国消费与地产：K型分化加速，定价权成为核心考验 - NOM：端午消费数据透露的信号，比表面数字更值得警惕
  \item [R036] 能源与大宗：铝缺口被低估，黄金与2013年本质不同，航运逻辑切换 - Citi：铝价缺口比市场想象的更深，中国铝业和中国宏桥的买入窗口已打开
  \item [R037] 宏观与利率：美联储“沉默”重塑全球定价，日本央行内部分歧加大 / 日本市场：结构性牛市第二幕，四大驱动力同步强化 - GS：日本央行内部已有人要求加速加息，关键看通胀是否会超2\%
  \item [R038] AI/算力：从“光学狂欢”到“网络务实”，半导体拥挤度达历史极值 / 风险偏好：AI主题拥挤但需求无虞，可持续投资“喘口气” - Bernstein：半导体唯一的游戏还在，但拥挤度已到历史极值
  \item [R039] 风险偏好：AI主题拥挤但需求无虞，可持续投资“喘口气” - Bernstein：车企造人形机器人，真正的护城河不在机器人本身
  \item [R040] AI/算力：从“光学狂欢”到“网络务实”，半导体拥挤度达历史极值 - JEF：AI服务器PCB的“Kyber延迟”比想象中更值得警惕
  \item [R041] 日本市场：结构性牛市第二幕，四大驱动力同步强化 - 摩根斯坦利：日本电子化学品正在经历一场“材料级”的供应链重构
  \item [R042] AI/算力：从“光学狂欢”到“网络务实”，半导体拥挤度达历史极值 - 摩根斯坦利：MS：AI投资正从“光学狂欢”转向“网络务实”
  \item [R043] 未被主题摘要引用 - JPM：AI颠覆论被高估，数字银行龙头才是真正的护城河
  \item [R044] 宏观与利率：美联储“沉默”重塑全球定价，日本央行内部分歧加大 - 摩根斯坦利：新兴市场债券的“油价逻辑”尚未兑现
\end{itemize}
\section{逐篇报告摘录}
\subsection{[R001] 摩根斯坦利：MS：新美联储的“沉默”本身就是最响亮的信号}
{\small\color{refgray}归类：宏观与利率：美联储“沉默”重塑全球定价，日本央行内部分歧加大}

这份来自MS全球宏观论坛的报告，核心判断并非关于加息终点或降息时点，而是指向一个更深层的结构性变化：美联储正在主动退出“预期管理”的角色，将定价权交还给市场。对于习惯了“看美联储脸色”做交易的投资者而言，这意味着一套全新的游戏规则正在建立。

报告由MS首席固定收益策略师Vishwanath Tirupattur领衔，联合首席全球经济学家Seth Carpenter、全球宏观策略主管Matthew Hornbach、首席美股策略师Michael Wilson等核心团队共同撰写。它不只是一份市场展望，更像是一份“新美联储操作手册”的导论。

报告明确指出，新任美联储主席Warsh有意在沟通中减少“前瞻指引”。他认为，减少对市场的前瞻性引导是其政策思路的核心。这种做法与他的前任们形成了鲜明对比。过去十年，市场习惯了从每一次FOMC会议声明、每一次记者会中寻找关于未来利率路径的蛛丝马迹。现在，这种“确定性”正在被主动撤回。

然而，市场的反应却耐人寻味。尽管Warsh明确表示希望避免任何像前瞻指引的暗示，市场却将他“没有对当前定价进行反驳”的行为，解读为一种默许。结果就是，市场定价反而更加紧密地追随了点阵图。这构成了一种新的博弈：美联储不再告诉你它会怎么做，但它不纠正你，就意味着它认可你的判断。

KC评论： 这相当于美联储把方向盘交给了市场，但又没有完全撒手。投资者需要从“猜美联储想什么”转向“猜市场如何猜美联储”。这种二阶博弈的复杂度，远高于过去。完整报告里包含了对SOFR曲线平坦化交易的详细论证，以及为什么在当前环境下，曲线平坦化比直接做空利率更具性价比。

Warsh的“沉默策略”最直接的后果，是利率波动率的结构性抬升。报告警告，由于缺乏前瞻指引，波动率不仅不低，反而可能进一步上升。这对所有利率敏感的资产类别都构成了根本性挑战。

对于抵押贷款支持证券（MBS）而言，这意味着一个经典的“双重打击”：更高的长期波动率会压缩抵押贷款指数的相对价值，同时利率曲线熊市平坦化会削弱银行对MBS的需求。MS因此战术性建议做空FNCL 5.5s。这并非看空美国房地产，而是基于一个清晰的相对价值判断——在当前环境下，MBS相对于国债的利差补偿不足。

KC评论： 很多投资者习惯用历史波动率来定价，但报告暗示，我们需要把波动率中枢上移作为一个结构性假设。这会影响所有依赖低波动环境的策略，从MBS套利到风险平价。报告里有一张图展示了当前抵押贷款指数相对于量化宽松时期的水平，非常直观地说明了为什么现在的位置并不安全。

这是报告中一个容易被忽略但极其重要的判断。首席美股策略师Michael Wilson认为，Warsh减少前瞻指引、更加依赖市场信号的做法，实际上会提高美联储决策的质量和时机。但真正驱动这轮牛市的核心，从来不是美联储的鸽派姿态，而是企业盈利。

报告指出，盈利的广度和韧性已经抵消了今年以来的估值压缩，并且很可能继续缓冲未来来自更紧缩货币政策或不确定性带来的冲击。这意味着，市场正在从“流动性驱动”切换为“盈利驱动”。那些盈利增长强劲、能消化更高利率成本的公司，将继续跑赢；而那些依赖低利率环境进行估值扩张的“故事型”公司，将面临更严峻的考验。

KC评论： 这实际上是在说，不要因为美联储“不友好”就轻易看空美股。只要盈利还在增长，市场就有支撑。但这里的核心问题在于，盈利的韧性是否已经被充分定价？以及，如果利率波动率持续高企，它最终会不会反过来侵蚀企业融资成本和消费者信心？这些二阶效应，报告并没有完全展开，值得在完整报告里深入阅读。

4. 欧洲市场的“上海解封交易”正在失去动能，但结构性机会仍在

报告用大量篇幅分析了欧洲市场与美联储政策之间的联动。一个有趣的发现是，欧洲各板块对“上海解封”主题的敏感度，与其对美联储政策

[... middle omitted ...]

提振，又受制于美联储紧缩带来的估值压力。报告中的图表显示，欧洲斯托克600指数在触底后的走势，与过去几次市场低谷的表现非常相似，但节奏更慢。这种“慢修复”是否意味着后续有更大的上行空间，还是仅仅是阶段性疲软？报告没有给出确定的答案，这恰恰是值得继续跟踪讨论的地方。

5. 报告尚未完全回答的关键问题：当美联储失去“可信度”时会发生什么？

这是整份报告最具思辨价值的问题。Warsh的策略在逻辑上很清晰：减少干预，让市场自己发现价格。但这里隐含了一个关键假设——市场能够有效吸收和解读美联储的“沉默”。然而，如果市场在解读过程中出现系统性偏差，比如将“不反驳”误读为“认可”，进而导致定价过度偏离基本面，美联储是否有能力在不引发恐慌的前提下进行纠偏？

报告提到了“市场信号”的增强，但并没有深入探讨这种信号机制在极端情况下的可靠性。当市场因为缺乏前瞻指引而出现流动性枯竭或价格剧烈扭曲时，美联储是否还能保持“沉默”？这个问题的答案，可能决定了未来一到两年内最大的一次市场波动。

\subsection{[R002] GS：亚洲太空经济正在从“故事”转向“订单”，但市场定价仍落后60\%}
{\small\color{refgray}归类：区域市场：亚洲太空经济被低估，印度制造业成本驱动型修复}

GS：亚洲太空经济正在从“故事”转向“订单”，但市场定价仍落后60\%

这份GS最新发布的亚洲太空经济专题报告，核心判断并不复杂：全球太空经济正从概念验证进入大规模部署阶段，而亚洲供应链公司在这一轮部署中扮演着不可替代的硬件核心角色。但市场对这些公司的定价，仍然停留在“概念主题”阶段，而非“盈利交付”阶段。

报告最值得关注的信号，不是太空经济的长期增长数字——那些已经足够震撼，LEO卫星市场预计到2035年增长7倍至1080亿美元——而是三个更具体的、对投资者有直接含义的发现：

第一，亚洲太空经济股票相对于全球同行的市盈率折价高达60\%，市净率折价25\%，且这一折价水平处于历史区间的低位附近。第二，全球太空主题基金和ETF的资产管理规模已经从2025年初的约10亿美元飙升至峰值约250亿美元，但亚洲供应链公司在这些基金中的配置比例仍然极低。第三，亚洲太空经济股票自2025年以来已经跑赢更广泛的区域市场，但在过去三个月进入盘整区间，而全球同行仍在上涨。

这意味着什么？如果太空经济确实在进入一个由订单和收入支撑的部署周期，那么亚洲供应链公司目前被低估的程度，可能提供了一个罕见的、基本面与定价之间存在系统性错位的窗口。

这份报告由GS亚洲策略团队Alvin So、Timothy Moe、Kinger Lau等多位分析师联合撰写，于2026年6月发布。它不仅仅是一份行业扫描，更是一份带有明确投资主张的“买入框架”。以下是我们从报告中提炼出的五个核心洞察。

KC评论： 报告的核心主张是“亚洲太空经济股票被低估”，但支撑这一主张的关键证据——60\%的市盈率折价——需要结合报告中的具体拆解来理解。折价本身并不自动意味着机会，关键在于折价是否被基本面恶化所解释。GS认为不是，因为盈利势头仍然积极。这正是完整报告中最值得细读的部分：分国家、分环节的盈利修正数据和相对估值图表。

1. 亚洲不是太空经济的“配角”，而是全球硬件供应链的“骨架”

许多投资者仍然将太空经济等同于马斯克的SpaceX或美国的星座运营商。但GS的报告清晰地指出，亚洲的角色远不止于此。

报告将太空经济价值链划分为四个环节：上游（发射与推进）、卫星制造与组件、地面段与下游应用、太空级材料与电子。在这四个环节中，亚洲几乎占据了硬件制造的全部关键节点。

中国的角色是“全栈式”的，从卫星总线、有效载荷到可重复使用火箭发动机，已经形成工业级量产能力，且受到ITU频率申报时间表的强约束，订单可见性至少到2028年。日本在发射系统和在轨服务领域占据独特位置，其商业碎片清除能力是目前全球唯一在运营的。韩国正在从火箭发动机制造向集成化太空基础设施模式转型，其大规模卫星生产设施预计2025年底投产。台湾则凭借半导体产业基础，主导了射频前端模块、相控阵天线和GNSS芯片等地面电子环节。

这一分工格局意味着，无论全球哪个星座运营商最终胜出，亚洲供应链公司都大概率是硬件供应商。这种“卖铲子”的商业模式，在太空经济的部署阶段，往往比星座运营商本身拥有更确定的收入和更低的执行风险。

2. 主权项目提供了“非可自由裁量需求”，这是与其他主题最根本的区别

太空经济投资常常被质疑的一个问题是：终端用户需求是否真实？如果企业客户和消费者对卫星宽带或地球观测服务的付费意愿不足，那么整个产业链的订单可能只是“一次性建设”。

GS的报告提供了一个有力的反论据：主权星座项目已经在2028年之前获得资金并进入执行阶段。这些项目不是由商业需求驱动的，而是由频谱申报、国家安全和产业政策驱动的。这意味着，即使商业端需求低于预期，上游供应链仍然有确定的订单支撑。

日本的1万亿日元太空战略基金、韩国的2000亿韩元私人太空基金、印度的IN-SPACe风险投资设施、新加坡的NSAS计划——这

[... middle omitted ...]

6年的早期采购阶段，驱动力来自地面能源和审批限制。

逻辑链条是这样的：地面数据中心的能源消耗和碳排放限制日益严格，审批周期也在拉长。而太空数据中心的优势在于，它可以利用太阳能发电，且不需要占用地面土地资源。虽然这一领域目前仍处于早期阶段，但GS认为，它与AI基础设施资本开支的关联性，将显著拓宽太空经济的投资逻辑，并延长主题的投资久期。

对投资者而言，这意味着太空经济不再只是“卫星和火箭”的故事。它正在与全球最大的资本开支主题——AI基础设施——产生交集。日本、韩国和台湾是这一领域的主要硬件受益者。

报告提供了一组令人惊讶的数据：全球太空主题基金和ETF的数量已经超过40只，资产管理规模峰值达到约250亿美元，而2025年初这一数字仅有约10亿美元。资金流入的速度堪称爆炸性。

但问题在于，这些资金绝大多数流向了美国上市的太空公司。亚洲供应链公司在全球主题配置和指数基准中的代表性仍然极低。GS认为，这种“定位与基本面之间的差距”是亚洲受益者潜在超额收益的关键来源。

\subsection{[R003] 摩根斯坦利：MS：日本股市的结构性牛市才刚刚进入第二幕}
{\small\color{refgray}归类：日本市场：结构性牛市第二幕，四大驱动力同步强化}

当全球投资者还在争论日本经济能否走出“失去的三十年”时，MS用一份夏季策略会报告给出了一个更激进的判断：日本股市正在进入一个结构性转折点，而支撑这个判断的四大驱动力——通胀稳定、政策加速、治理改革、居民入场——都在同步强化。

这不是一个简单的周期性反弹故事。MS对TOPIX的基准目标价是2027年6月底达到4,300点，较当前水平约有8\%的上涨空间。但真正值得关注的不是这个数字，而是报告提出的一个核心逻辑：日本股市正在从“低通胀、低传导、低利润”的旧范式，切换到一个全新的定价框架。在这个新框架下，历史估值区间不再构成上限。

对于中国投资者而言，理解这个切换，远比纠结于短期汇率波动或中东油价冲击更重要。因为这意味着日本资产在全球配置中的角色正在发生根本性变化。

过去三十年，日本投资者最怕听到的两个字就是“通胀”。通缩思维深植于企业定价行为、居民消费习惯和资产配置逻辑中。但MS认为，日本正在进入一个“更稳定、更可持续”的通胀环境，而这一变化正在从根本上改变企业行为。

报告的核心论据是产出缺口正在逐步收窄。当经济不再有闲置产能，企业开始具备定价权，工资上涨开始传导至终端价格，形成一个正向循环。这听起来像是宏观教科书的标准叙述，但放在日本语境下，这是一个历史性的转变。

从数据上看，日本企业已经在行动。报告指出，自疫情后经历进口价格飙升以来，日本企业持续致力于改善利润率。这意味着即使日元升值，企业盈利的脆弱性也比过去大幅降低。一个关键证据是：美元兑日元与日本股市的正相关性在过去十年中持续下降，日元不再是日本企业盈利的唯一驱动因素。

KC评论： 这个判断对投资者的含义很直接——如果你还在用“日元升值=日本股市利空”的老框架来交易，你可能已经错过了结构性的变化。企业定价能力的提升意味着，即使日元从160回到140，日本企业的盈利可能也不会像过去那样大幅缩水。完整报告中有一张进口价格与国内企业商品价格的对比图，清晰地展示了这一传导链条的改善，值得仔细研究。

石破茂领导的自民党在众议院选举中赢得316席，单独超过三分之二多数。这个结果对资本市场的影响，远不止于“政策可预测性提高”这么简单。

MS从三个维度拆解了政治稳定的意义：更高的政策可预测性、更快的决策速度、以及增长战略更高的执行确定性。这三者叠加，直接作用于两个关键变量：中期盈利展望和市场要求的风险溢价。

这里有一个容易被忽略的细节。报告提到，政治稳定带来的不仅仅是政策连贯性，更重要的是加速了“增长导向型政策”的推进，包括国家安全相关的产业政策。这意味着，日本正在从一个“防御型经济体”转向一个“战略型经济体”，政府开始主动引导资源配置。

对于资产定价而言，风险溢价的下降意味着同样的盈利预期可以支撑更高的估值。这正是MS敢于给出17.5倍远期市盈率目标的原因——这个估值水平在历史上的日本市场并不常见，但在新范式下却是合理的。

3. 公司治理改革从“口号”进入“现金”阶段，ROE提升还有很大空间

东京证券交易所2022年4月的市场重组改革已经带来了显著效果：Prime市场的市净率平均上涨了约40\%。但MS认为，真正的重头戏还在后面。

报告预期，2026年6月左右，日本将进行公司治理准则的修订，核心方向是推动企业更有效地使用现金。这意味着，过去几年日本企业“口头承诺提升股东回报”的阶段即将结束，接下来将进入“真金白银回馈股东”的阶段。

从数据上看，日本企业的现金持有比例在全球主要市场中仍然偏高。即使经过几年的改善，ROE水平与欧美同行相比仍有明显差距。治理准则的修订一旦落地，将直接推动企业通过回购、增派股息、业务重组等方式释放价值，进而通过市净率和市盈率的双重扩张支撑股价。

KC评论： 这是报告中一个容易被低估的催化剂。很多投资者关注日本公司治理改革

[... middle omitted ...]

日本家庭金融资产中，现金和存款的比例仍然高达50\%以上，远高于美国的10\%和欧洲的30\%。即使只有一小部分资金从存款转向股票，其边际影响也足以改变市场的供需格局。

更重要的是，这个趋势是自我强化的。当通胀稳定、工资上涨、股价上升形成正向循环，居民从“储蓄型”转向“投资型”的意愿会持续增强。NISA只是加速了这个过程。

5. MS没有完全回答的一个关键问题：中东风险的真实尾部有多大

报告坦诚地展示了宽幅的情景分析。基准情景假设中东局势缓和，TOPIX目标4,300点。但熊市情景下，如果中东冲突升级导致油价飙升并引发全球衰退，TOPIX可能跌至2,600点，较当前水平下跌约35\%。牛市情景则看到4,900点。

这个宽幅区间（2,600-4,900点）本身就说明了一个事实：当前市场的核心不确定性不在日本国内，而在外部。MS对宏观情景的假设——中东能源供给冲击在2026年温和拖累全球增长，2027年复苏——是一个“有序”的情景。但地缘政治的尾部风险本质上是不可预测的。

\subsection{[R004] UBS：霍尔木兹海峡的VLCC运价已不是关键，集装箱的“前置抢运”才是}
{\small\color{refgray}归类：能源与大宗：铝缺口被低估，黄金与2013年本质不同，航运逻辑切换}

UBS：霍尔木兹海峡的VLCC运价已不是关键，集装箱的“前置抢运”才是

全球航运市场的目光正集中在霍尔木兹海峡的恢复通行上。UBS最新发布的周度航运供应链报告给出了一个反直觉的判断：VLCC运价的短期飙升是事件性的，真正值得关注的、且已经在高频数据中显现的结构性变化，是亚洲至美国航线的集装箱“前置抢运”正在加速。这轮运价的上涨，其驱动力正在从“绕航”切换为“抢运”，其影响深度和持续性可能远超市场预期。

这份报告的价值不在于复述海峡通航的新闻，而在于它利用UBS Evidence Lab的高频数据，提供了一套区分“噪音”与“信号”的观测框架。对于关注全球贸易、航运资产、以及供应链韧性的决策者而言，理解这个信号切换，比猜测下一周运价的涨跌更重要。

1. 霍尔木兹的恢复是“假摔”信号，真正的运价支撑来自大西洋的“抢油”

报告中最直接的冲击来自VLCC。数据显示，上周中东至中国的VLCC平均收益飙升了87\%，达到约19.5万美元/天。这个数字看起来很惊人，但UBS的解读非常克制：这更多是“预期”的体现，即市场预期波斯湾重新开放后，会有大量油轮涌入装载货物，导致短期运力紧张。

然而，报告中的另一组数据揭示了更底层的逻辑。VLCC一年期期租费率保持稳定。这意味着，真正的产业资本（需要长期锁定运力的炼油厂或贸易商）并未恐慌。现货市场的飙升，更多是投机性运力与短期补库需求的一次性冲高。一旦首批货物装运完毕，运价将迅速回归。

KC评论： 这提醒我们，不能把事件驱动的脉冲式上涨，误判为周期的反转。UBS报告里没有明说但值得追问的是：如果中东恢复，那批从大西洋绕道好望角的VLCC会如何？它们会形成新的运力冗余，对运价构成下行压力。要理解这背后的二阶效应，需要完整阅读报告中对“全球船队动态”的拆解。

真正需要警惕的，不是VLCC的脉冲，而是大西洋盆地的“抢油”行为可能正在发生。报告提到，从霍尔木兹海峡向东驶出的油轮数量自上周五以来显著增加。这暗示，部分原本依赖中东的买家，可能正在全球范围内抢购替代油源，从而拉升了西非、美国湾等地的运价，这才是VLCC运价获得支撑的长期逻辑。

集装箱市场的数据更具说服力。中国主要港口的集装箱吞吐量周环比增长3\%，同比增长2\%。这本身并不意外，但结合上海出口集装箱运价指数（SCFI）周环比上涨5\%、同比暴涨67\%来看，供需缺口仍在扩大。

UBS报告中最值得关注的一个判断，是“亚洲至美国航线的集装箱运输量自5月以来显示出强劲的同比增长”。这不是季节性因素可以解释的。结合图6“从亚洲到美国的集装箱新在途量”显示出的前置装载模式，一个清晰的逻辑链条浮现出来：面对持续的地缘政治不确定性（红海绕航仍在高位，较去年同期高出4-5\%），以及潜在的关税或贸易政策风险，进口商正在主动将下半年的货物提前到上半年发运。

KC评论： 这是一个典型的“抢跑”行为。它意味着，当前的高运价和高运量，透支的是未来的需求。一旦“抢运”结束，市场可能面临需求真空。报告中的“全球海运目的地监测”数据集，能帮我们量化这种前置装载的规模到底有多大，以及它何时会达到拐点。这是决定航运股和货代企业下半年业绩的关键变量，完整报告里有更细致的区域分解。

3. 散货市场的“结构性”上涨与造船价格的高位盘整，指向不同的产业周期

与集装箱和油轮的热闹不同，干散货市场上周有所软化，波罗的海干散货运价指数（BDI）周环比下跌4\%。但报告同时指出，好望角型散货船的平均收益周环比上涨了11\%。这揭示了干散货市场的内部结构分化：好望角型船（主要运输铁矿石和煤炭）的运价韧性，远强于巴拿马型船（主要运输谷物和小宗散货）。这背后是不同大宗商品的需求周期错位。

造船市场则给出了另一个信号。新造船价格指数周环比持平，但需求依然强劲，特别是VLCC和液化气

[... middle omitted ...]

跌。

但UBS的数据似乎在暗示另一种可能性：即便红海恢复通航，当前的贸易格局和运力部署已经发生了永久性改变。进口商为规避风险而建立的多元化供应链、以及为应对不确定性而增加的库存，都不会轻易逆转。这意味着，即使绕航结束，运价中枢也可能高于疫情前的水平。

这是一个巨大的不确定性。报告没有给出明确的结论，但它提供了跟踪这一问题的关键数据集——“全球海运贸易监测”和“每日海运贸易中断监测”。通过这些数据，我们可以实时观察船舶的实际航线和通行量，从而比市场更早地判断“新常态”是否到来。

基于UBS这份报告，我们建议读者建立一个三层观察框架，而不是简单跟踪运价：

第一层，看“量”而非“价”。运价是结果，运量是原因。重点关注中国主要港口吞吐量、亚洲至美国航线在途量等高频数据，判断“前置抢运”何时见顶。

第二层，看“船型”而非“指数”。VLCC、集装箱、散货船面临完全不同的供需格局。将三者混为一谈，会得出错误的结论。当前，集装箱的确定性最强，油轮看事件驱动，散货看结构分化。

\subsection{[R005] NOM：日本加息周期已过中点，但市场真正该关注的是财政牌}
{\small\color{refgray}归类：宏观与利率：美联储“沉默”重塑全球定价，日本央行内部分歧加大 / 日本市场：结构性牛市第二幕，四大驱动力同步强化}

当全球投资者的目光仍聚焦于日本央行何时结束负利率、加息路径是否被打断时，NOM在最新一期《日本经济周报》中给出了一个值得细品的判断：日本央行的加息周期可能已经走过了一半。更关键的是，随着6月货币政策会议尘埃落定，市场关注重心正从央行转向政府——增长战略与消费税减税将成为决定日本资产定价的下一个变量。

这份报告的核心信号并不复杂：日本经济、物价和金融条件三者目前都支持继续加息。但NOM同时指出，真正需要投资者警惕的，不是加息本身，而是财政政策的不确定性，尤其是消费税减税方案的资金来源与推进节奏。后者可能比加息更深刻地影响日本内需、通胀路径和国债定价。

NOM在报告中明确写道：“我们认为，日本央行可能已经过了本轮加息周期的中点。” 这一判断基于三个关键信号：一是政策委员会成员Asada投下反对票，暗示鸽派力量正在集结；二是央行对经济风险的评估出现微妙变化，认为“经济显著放缓的风险相比此前有所下降”；三是央行在措辞上删除了“实际利率处于极低水平”的表述，代之以“金融条件一直宽松”。

KC评论： 措辞调整看似微小，但信号意义很强。央行主动降低对宽松程度的评估，意味着未来加息的理由将更依赖数据而非“补涨”逻辑。这对市场意味着，加息节奏可能从“追赶式”转向“确认式”，每一次加息都需要更强的通胀或增长数据支撑。

NOM维持其基准情景：2026年12月和2027年6月各加息一次，共两次。风险情景则是在此基础上再加一次，时间点落在2026年10月、2027年3-4月和9-10月。无论哪种情景，加息周期的终点都指向2027年中后期，而非市场此前预期的2028年或更晚。

一个容易被忽视的细节是，NOM对中东局势的评估并非“问题解决”，而是“问题开始缓解”。美国与伊朗于6月18日签署谅解备忘录，但霍尔木兹海峡的油轮通行量已从冲突前的每天60-70艘骤降至几乎为零。即使邻近的曼德海峡通行量在3月下旬后略有回升，但替代航线远未建立。

这意味着，即便地缘政治风险溢价下降，量化的供应约束——即原油和石油产品的进口量——仍将在一段时间内压制日本经济的复苏弹性。NOM提醒投资者关注两个关键指标：石油进口量的恢复程度，以及主要经济体制造业面临的原材料采购延迟幅度。

KC评论： 市场往往对“停火”或“谅解”这类消息反应过度乐观，但供应链的物理恢复需要时间。日本作为能源进口依赖度极高的经济体，其通胀与增长对原油供应的敏感度远高于美国或欧洲。这份报告提醒我们，不要因为油价短期回落就忽视供应端风险。

NOM预计，日本核心CPI通胀可能在2027年一季度见顶。这一判断基于价格传导的时滞：日元计价的进口价格已大幅上涨，这些成本将从上游（PPI）向下游（CPI）逐步传导，整个过程大约需要6-9个月。

但报告同时指出，日本央行在6月声明中明确警告“核心CPI存在向上偏离2\%目标的风险”。这意味着，即便通胀如预期在2027年一季度见顶，其峰值水平可能高于央行和市场的当前预测。如果这一风险兑现，加息的终点利率也将相应上调。

一个支持继续加息的硬数据是贷款增长。日本银行和信用金库的贷款余额同比增长6\%，超过CPI涨幅。城市银行的贷款增长尤为显著。住房贷款也在持续增长，尽管加息推高了房贷利率，但房价上涨抵消了部分利率压力，贷款需求依然旺盛。

NOM认为，持续的贷款增长是“政策利率低于中性利率”的直接证据。只要贷款增速保持高位，央行就有理由继续加息，直到金融条件从“宽松”转向“中性”或“紧缩”。

5. 财政牌才是下半年的最大变量：370万亿投资与消费税减税

如果说加息是“已知的未知”，那么财政政策就是“未知的未知”。NOM明确指出，随着6月货币政策会议结束，市场焦点将转向政府的增长战略和消费税减税方案。

增长战略方面，日本增长战略委员会正在考虑在

[... middle omitted ...]

027年4月起将食品饮料消费税降至1\%，并在2027年和2028年秋季实施与收入挂钩的现金发放，总额相当于消费税收入的1\%（约6000亿日元）。但反对党强烈反对这一方案，尤其是资金缺口问题——每年约5万亿日元的成本尚未明确来源。

6. 报告尚未完全回答的关键问题：消费税减税的资金来源与政治博弈

NOM在报告中坦承，消费税减税的资金来源“仍未明确”。一个可能的选项是利用外汇基金特别账户（FEFSA）的盈余资金，但该账户每年的盈余规模有限，且部分已被指定用于外汇操作。另一个选项是发行专项国债，但这与自民党“不依赖赤字融资债券”的承诺相矛盾。

更复杂的是政治博弈。自民党在众议院拥有三分之二多数，理论上可以强行通过减税法案。但若该法案未能在超党派的国家社会保障委员会中获得充分讨论，强行通过将损害执政党的公信力。国民民主党（DPP）提出了一个相对接近的方案，但要求现金发放时间提前至2026年底，而非2027年秋季。两党能否在6月底前达成一致，将直接决定减税方案的推进节奏。

\subsection{[R006] UBS：黄金正处在“无人区”，但这次和2013年不一样}
{\small\color{refgray}归类：能源与大宗：铝缺口被低估，黄金与2013年本质不同，航运逻辑切换}

一方面是实际利率飙升、市场持续定价加息，金价承压明显；另一方面，央行购金、私人部门多元化配置、以及不断膨胀的美国债务，又在底部提供着结构性支撑。UBS在最新发布的贵金属研报中，用一个精准的比喻描述了这种状态：“precious metals navigate no-man's land”——贵金属正穿行于一片无人区。

这份研报的核心判断是：黄金的短期风险确实在上升，但长期逻辑并未被打破。真正值得关注的，不是金价会跌多少，而是这一轮调整结束之后，谁会率先重新入场。

许多投资者看到金价回调，第一反应是翻出2013年4月黄金崩盘的记忆。当时金价单年下跌29\%，背后是美联储货币政策正常化的预期转向、塞浦路斯央行抛售黄金的恐慌蔓延，以及实际利率由负转正的重大拐点。

UBS在报告中专门做了一组对比分析，结论很清晰：2026年的宏观环境与2013年存在根本性差异。

2013年黄金面临的是“货币政策体制转换”——美联储从极度宽松转向退出QE，实际利率从负值区域翻正，这对零息资产的打击是致命性的。而当前，美联储才刚刚完成175个基点的降息周期，处于“暂停观察”状态，政策方向仍偏向宽松。实际利率虽然也在上升，但起点和幅度都远不及2013年。

更重要的是，2013年市场担心的核心问题是“央行会不会集体抛售黄金”，而今天市场的担忧恰恰相反——央行购金仍在持续，只是买家结构在发生变化。2013年全球通胀处于下行通道（从2.0\%降至1.0\%），而当前通胀从2.7\%升至4.2\%，这本身就是黄金的长期利好。

KC评论： UBS这张对比表格非常值得仔细看。它不是在说黄金不会跌，而是在说“即使跌，驱动逻辑也不同”。投资者最大的风险不是金价下跌，而是用2013年的框架去理解2026年的市场。完整报告中包含了这张详细的宏观对比表，覆盖了美元、通胀、就业、财政赤字、金融条件等15个维度，是判断当前黄金所处阶段的核心参考。

2. 美国经济韧性背后藏着AI的结构性依赖，这才是黄金最大的短期风险

黄金多头当前最头疼的问题，是美国经济数据持续超预期。UBS在报告中指出，美国投资和消费的韧性，很大程度上是由人工智能和科技投资拉动的。

AI相关投资不仅直接支撑了企业资本开支，还通过财富效应支撑了高收入家庭的消费。这创造了一个自我强化的循环：AI驱动增长，增长支撑就业和收入，消费维持经济热度，经济热度推迟降息预期，降息推迟又压制金价。

但UBS也提醒，这种依赖AI的增长模式本身存在脆弱性。如果AI投资出现预期调整，或者生产率提升的兑现速度慢于市场定价，美国经济的下行风险将迅速暴露。从另一个角度看，AI带来的生产率提升，如果最终转化为更低的通胀和更高的实际增长，反而会成为黄金的长期逆风。

KC评论： 报告在这里留下了一个关键伏笔——AI对美国经济的支撑到底有多可持续？UBS自己的宏观预测显示，美联储要到2027年才会重启降息，这意味着黄金在未来12-18个月可能都要面对一个“不友好”的宏观环境。完整报告中包含了对美国经济分项数据的详细拆解，以及UBS内部对GDP、失业率、通胀和利率的完整预测表，这些是判断黄金拐点时间窗口的基础。

市场对央行购金的叙事往往过于简化——“央行在买黄金，所以金价有底”。UBS的研报揭示了一个更复杂的现实：央行购金的参与者结构正在发生重大转变。

2022年的前五大黄金买家是土耳其、中国、乌兹别克斯坦、印度和卡塔尔。到2026年，这个名单已经变成了波兰、印度、土耳其、中国和捷克。中国和土耳其的购金节奏在放缓，但波兰、捷克等东欧国家以及部分亚洲央行正在接棒。

更值得关注的是，UBS测算，其他官方机构（主权基金、养老基金等）的资产管理规模仍在增长，而这些机构目前对黄金的配置比例极低。如果这些机构将黄金配置比例提高1\%，对

[... middle omitted ...]

S的数据表明，央行的决策框架根本不是“高抛低吸”。真正需要关注的，是哪些央行还有空间继续增持，以及主权基金这个更大的买家群体何时入场。完整报告中有2022年与2026年购金国对比图、官方机构资产增长图，以及1\%配置调整对金价的模拟测算，这些图表是理解央行购金长期逻辑的关键。

金价处于历史高位，实物需求却出人意料地坚挺。UBS的数据显示，全球实物黄金（金条金币）需求总量并未因价格创新高而崩塌。

更关键的是区域结构的变化。亚洲目前贡献了全球超过70\%的实物黄金投资需求，且几乎所有区域的实物需求都录得正增长。这与2013年形成了鲜明对比——当时西方投资者在金价暴跌中大规模抛售，而东方投资者被动接盘。

这一次，亚洲投资者在价格高位仍在持续买入。这种结构性变化意味着，黄金定价权的重心正在从COMEX期货和西方ETF转向亚洲实物市场。如果亚洲经济体的居民财富持续增长，黄金作为储蓄和保值工具的底层需求将更加稳固。

5. 白银和铂族金属的处境更复杂，但中国期货市场正在给出信号

\subsection{[R007] Bernstein：数据中心每兆瓦值多少钱，答案藏在商业模式里}
{\small\color{refgray}归类：未被主题摘要引用}

Bernstein：数据中心每兆瓦值多少钱，答案藏在商业模式里

AI基础设施的军备竞赛中，一个核心问题始终困扰着投资者：当一家数据中心或GPU云公司签下一份合同，这到底是一笔好交易吗？市场往往被“数十亿美元”“数万兆瓦”这类宏大数字吸引，却忽略了单位经济学的根本差异。

Bernstein最新发布的研究报告，提供了一个极为实用的分析框架：每兆瓦（MW）能产生多少收入。这个看似简单的指标，在拆解不同商业模式后，揭示出惊人的差异——从不到100万美元到超过5500万美元。差异本身不是答案，理解差异背后的结构性原因，才是判断哪些公司真正拥有护城河的关键。

这份报告的核心判断是：对于规模尚未达到临界点的提供商，Bernstein继续偏好托管模式（colo model）而非GPU云模式（neocloud model）。但报告同时指出，当GPU云公司拥有自有房地产时，部分合同同样具有吸引力。

这不是一个简单的“谁更好”的结论，而是一个关于商业模式选择、资本配置能力和风险回报权衡的深度分析。

Bernstein将数据中心相关公司划分为三类，每类在每兆瓦收入上呈现清晰的阶梯状分布。

第一类是传统托管REITs。Equinix以约500万美元/兆瓦位居首位，这得益于其零售导向模式——不仅出租机柜空间，还叠加了网络互联服务和高价值收入。Digital Realty和CoreSite等更偏向批发模式的公司，收入密度在300万至400万美元/兆瓦之间。新兴AI基础设施提供商（如WULF、CIFR）则更低，主要因为它们在偏远地区预建未完成的项目。

第二类是GPU云公司。CoreWeave的收入密度达到2080万美元/兆瓦，Nebius为840万美元/兆瓦。它们之所以能获得更高收入，是因为在同一兆瓦电力上叠加了硬件（GPU）和平台服务，实现了垂直整合。

第三类是比特币矿工转型的AI基础设施公司。IREN是一个特例，其收入密度达到1040万美元/兆瓦，接近GPU云公司水平，因为它也运营垂直整合的GPU云模式。而CIFR、WULF、CORZ等采用托管模式，收入密度仅为175万美元/兆瓦，反映的是纯粹的“电力房东”定位。

KC评论： 收入密度差异本身不构成投资判断。高收入密度的公司往往也承担更高的运营成本和资本开支。关键在于，这份收入在扣除所有成本后，每兆瓦能留下多少利润。Bernstein在完整报告中提供了详细的利润率和资本回报率拆解，这是判断商业模式可持续性的关键。

托管模式和GPU云模式在运营成本结构上的差异，比收入端更为根本。

对于托管REITs，电力成本和设施运营费用大部分通过合同结构转嫁给客户。Digital Realty的运营利润率仅为8.2\%，而Equinix达到66.2\%。这一巨大差异并非因为Equinix效率更高，而是其收入结构中包含了大量高利润的互联服务收入，且其开发支出在利润表中占比更低。

GPU云公司则截然不同。CoreWeave和Nebius的每兆瓦运营支出分别高达1860万美元和3330万美元——它们自己采购电力、运营IT堆栈、承担GPU折旧、支付编排软件和支持工程成本。CoreWeave的运营利润率约为10.8\%，与Digital Realty相当，但收入基数完全不同。

新兴AI基础设施公司则走向另一个极端。它们签订的长期合同通常采用“总净”（gross net）或“三净”（triple net）结构，将大部分或全部电力及运营成本转嫁给租户，资产层面的EBITDA利润率可达80\%-95\%。

这一对比揭示了一个关键洞察：每兆瓦收入的数字本身会说话，但真正决定投资回报的是“谁在承担什么风险”。

KC评论： 投资者容易陷入“高收入密度=好生意”的思维陷阱。Bernstein的分析提醒我们，需

[... middle omitted ...]

为未来增长储备了空间。

GPU云公司则处于快速扩张期。CoreWeave在AI需求激增的背景下建立了集中的GPU部署，实现了强劲的营收增长。但高收入密度背后，是密集的基础设施和网络投入。这些公司的单位经济学仍在演变中，尚未达到稳定状态。

IREN提供了一个有趣的中间案例。它从北美市场扩展为全球公司，通过收购Nostrum进入欧洲，并进入澳大利亚市场，目前拥有5.8GW的电力容量。报告特别指出，IREN与NVIDIA的新合同条款较之前与微软的合同有明显改善——每IT兆瓦的收入从950万美元提升至1400万美元。这表明，随着规模扩Daiwa执行能力得到验证，单位经济学可以持续优化。

对于新兴AI基础设施公司，WULF和CIFR则展示了另一种路径。它们以REIT-like模式运营，通过长期合同锁定利用率，获得高可见性的现金流。CIFR获得AWS作为回头客，证明了其执行能力获得了超大规模客户的信任。这些公司的未来上行空间，取决于能否将拥有的电力站点转化为创收的托管合同。

\subsection{[R008] UBS：当炼厂开工率跌至45\%，化工行业下一个机会在AI制冷剂}
{\small\color{refgray}归类：中国化工：传统品疲弱，氟化工与MDI的结构性机会}

UBS：当炼厂开工率跌至45\%，化工行业下一个机会在AI制冷剂

中国化工行业正在经历一轮罕见的“双低”震荡：原油库存自6月初持续下降，地方炼厂开工率已跌至45.74\%，而聚酯、聚丙烯等主要化工品的产能利用率同比下滑幅度普遍在10个百分点以上。市场弥漫着对需求疲软的担忧，但UBS最新发布的化工周报给出了一条清晰的差异化路径——当传统化工品陷入库存与开工率的双重压力时，AI算力驱动的含氟化学品需求正在打开一个被低估的定价空间。

这份报告的核心判断是：市场尚未充分定价AI对氟化学材料带来的结构性增长机遇，部分氟化工企业存在明确的估值重估潜力。而这一判断的逻辑，需要从当前化工行业最真实的运营数据说起。

1. 炼厂开工率持续下滑，原油库存降至年内低位，但这不是系统性危机

UBS周报数据显示，截至上周，中国主营炼厂开工率小幅回升0.23个百分点至67.41\%，但地方炼厂开工率环比大幅下降2.9个百分点至45.74\%。中国原油总加工量环比下降12.8万桶/日，至1232.1万桶/日。与此同时，根据Kpler数据，中国原油库存自6月初以来持续下降，截至6月22日已降至12.68亿桶。

这两个数据放在一起看，容易让人联想到需求崩塌。但UBS的解读更克制：开工率分化本身说明问题不在总量，而在结构。主营炼厂开工率相对稳定，地方炼厂则面临更为严峻的利润压力和原料获取能力差异。原油库存下降叠加开工率走低，更多反映的是炼厂主动去库存和调节开工节奏，而非终端需求的全面失速。

KC评论： 这个地方最容易被误读。很多人看到开工率45\%就会联想到“行业萧条”，但UBS把主营炼厂和地方炼厂分开看，就是为了说明这轮压力更多集中在成本传导能力弱、议价权低的环节。对于下游化工品而言，原料端的价格压力正在通过炼厂减产释放，这反而可能为后续某些化工品的价差修复创造条件。完整报告里有一张炼厂开工率的分化图表，值得仔细看。

2. 化工品开工率普遍承压，但库存数据透露出需求并非全线溃败

报告覆盖的主要化工品中，除了TDI开工率环比上升7个百分点至71\%、PTA开工率环比上升4个百分点至69\%外，其余品种的产能利用率普遍低于去年同期。聚酯长丝开工率同比下降15个百分点至74\%，PP开工率同比下降15个百分点至63\%，PVC开工率同比下降9个百分点至70\%。

库存数据则呈现分化格局。PP库存同比下降26\%，PE库存同比下降5\%，PVC库存同比下降21\%，二氧化钛库存同比下降43\%。但聚酯长丝库存同比飙升112\%，硅酮DMC库存同比下降21\%。

这里的关键洞察在于：库存下降快的品种，并非需求旺盛，而是供给收缩更猛烈。PP和PVC的库存去化主要来自开工率大幅下降，而非下游采购回暖。聚酯长丝库存高企则直接反映出纺织终端需求的疲弱，即使开工率已经大幅下调，库存仍在累积。

KC评论： UBS的库存数据来自样本工厂，不是全行业统计，但趋势指向明确。当前化工品市场正在经历一轮“供给主动出清”的阶段，而不是需求自然修复。这意味着，接下来谁能在供给出清中保持开工率、同时控制库存，谁就能在下一轮价格反弹中占据主动。这份周报里对每个品种的开工率和库存做了连续追踪，如果你在关注具体标的，这些数据是判断拐点的基础。

3. 氟化工是报告中唯一的“结构性增长叙事”，AI算力正在改写估值坐标

在通篇以“下降”和“承压”为主调的周报中，UBS专门发布了一份氟化工行业报告，并明确提出：氟化学品在半导体和数据中心的应用正在快速增长，AI带来的算力需求提升和材料体系迭代，将推动氟化材料的增长。

目前，氟化工材料供应商的估值相较电子化学品企业存在显著折价。UBS认为，市场尚未充分定价AI给氟化学材料带来的增长机会，部分氟化工企业具备进一步重估的潜力。

这个判断的逻辑链条很清晰：A

[... middle omitted ...]

的定价锚正在向AI资本开支和半导体设备投资迁移。估值折价的存在意味着，如果这个逻辑被市场接受，重估空间可能相当可观。完整报告里对氟化工的细分应用场景和具体受益标的都有拆解。

4. 报告尚未完全回答的关键问题：AI对氟化学品的需求何时进入利润表

UBS的氟化工逻辑在方向上是清晰的，但报告没有完全回答几个关键问题。

第一，AI对氟化学品的需求增长，在时间维度上有多快？是未来12个月就能体现在订单中，还是需要更长的孵化周期？报告提到了“fast growth”和“AI tailwinds”，但没有给出具体的需求增速预测或渗透率假设。

第二，当前氟化工企业的产能布局和客户结构，是否已经为AI需求做好了准备？哪些企业已经进入了数据中心或半导体客户的供应链，哪些还在导入阶段？估值折价的存在，部分原因可能是市场对这些企业能否真正兑现AI红利存在疑虑。

第三，如果AI需求确实快速增长，氟化工行业的供给端是否存在瓶颈？产能扩张需要多长时间？如果供给跟不上需求，价格弹性有多大？

\subsection{[R009] GS：印度电缆与EMS公司正经历一轮“成本驱动型”收入修复}
{\small\color{refgray}归类：区域市场：亚洲太空经济被低估，印度制造业成本驱动型修复}

印度资本市场正在迎来一轮财报季预热。GS在最新发布的印度电缆与电子制造服务（EMS）行业前瞻报告中，给出了一个清晰且反直觉的判断：即将到来的第一季度（Q1FY27）收入增长，主要驱动力不是需求爆发，而是上游原材料价格上涨带来的价格传导与渠道补库存。对于关注印度制造业、供应链转移以及新兴市场资产定价的读者来说，这份报告的价值不在于预测具体数字，而在于它揭示了一个关键的结构性信号——当成本通胀成为收入增长的核心引擎时，企业的盈利质量、估值逻辑和投资风险都会发生微妙但重要的变化。

报告覆盖四家公司：电缆领域的Polycab India（中性评级）和KEI Industries（买入评级），以及EMS领域的Amber Enterprises（中性评级）和Dixon Technologies（卖出评级）。GS调高了前三家的盈利预测和目标价，唯一调低展望的是Dixon——这家印度最大的手机EMS厂商正面临记忆体成本通胀挤压利润率、PLI补贴退坡以及合资项目审批延迟的多重压力。

这份报告的核心主张可以浓缩为一句话：印度制造业正在从一个“量价齐升”的阶段，切换到一个“价升量平、利润承压”的新区间。对于投资者而言，真正需要追问的不是Q1收入增长多少，而是这种增长能否持续、利润能否同步改善、以及不同公司在成本压力下的竞争地位将如何分化。

1. 铜铝价格同比涨幅创16个季度新高，收入增长主要靠“涨价”而非“走量”

GS的数据显示，2026年第一季度（印度财年Q1FY27），铜和铝的LME价格同比涨幅分别超过35\%，创下过去16个季度以来的最大涨幅。这一数字本身并不令人意外，但GS将其与公司层面的收入预期联系起来后，结论变得耐人寻味。

对于Polycab和KEI这样的电缆电线企业，GS预计Q1的销量增速仅为中到高个位数（mid-to-high single digit），但收入增速将分别达到36\%和27\%的同比水平。这意味着，收入增长的主要贡献来自价格传导，而非需求扩张。GS在报告中明确写道：“revenue growth for these companies will be largely pricing-driven during the quarter。”

KC评论： 这不是一个“需求旺盛、供不应求”的故事，而是一个“成本上涨、被迫提价”的故事。两者在报表上都会体现为收入增长，但对利润率和可持续性的含义完全不同。完整报告中包含了铜铝价格月度环比变化图表（Exhibit 5-7），这些数据能帮助读者判断价格压力的持续性以及何时可能出现逆转。

2. 电缆企业：量价分离的格局下，KEI的产能扩张成为关键差异化因素

在电缆电线板块，GS维持了对KEI的买入评级，对Polycab则给出了中性。为什么同样是受益于原材料价格上涨，评级却不同？答案藏在产能和市场份额的动态里。

KEI是印度第二大电缆电线企业，市场份额约8\%。GS认为其具备几个结构性优势：拥有超高压（EHV）电缆生产能力，包括220kV及以上等级，最高可生产400kV电缆；在出口市场取得进展，近期获得美国市场准入；正在推进的产能扩张将在12-24个月内逐步释放；经销商网络计划每年增长5\%-6\%。

相比之下，Polycab虽然体量更大、品牌更强，但GS认为其风险回报已经趋于平衡（well-balanced），暗示当前股价已经反映了大部分利好。报告中对Polycab的FMEG（风扇、照明等消费品）业务和EPC（工程总承包）业务也给出了相对保守的预期。

KC评论： 在成本通胀环境下，能够把价格涨上去、同时不丢失市场份额的公司，才是真正的赢家。KEI的EHV电缆壁垒和出口扩张，意味着它在价格谈判中拥有更强的议价权。而Polycab的多元化业务中，FMEG面临更激烈

[... middle omitted ...]

利润率（以百分比计）会显著下降。GS预计Q1手机业务利润率将从3.6\%降至3.0\%。

第二，PLI（生产挂钩激励）将从FY27开始逐步退出。过去几年，PLI是印度EMS厂商利润的重要补充。Dixon的Mobile \& EMS部门在FY26的利润中，PLI贡献了相当比例。一旦退出，即使收入增长，净利润也可能承压。

第三，与Vivo的合资项目仍在等待政府批准，导致FY27的产能爬坡预期放缓。GS因此将FY27的EPS下调3\%，同时将部分预期量推迟到FY28。

KC评论： Dixon的故事是对“印度制造”光环的一次现实检验。市场往往把EMS简单理解为“受益于供应链转移”，但Dixon的案例说明，当成本结构、补贴政策和客户集中度三大变量同时出现不利变化时，高增长预期可能难以兑现。完整报告中对Dixon的季度预期表（Exhibit 9）还包含了一个值得关注的细节：市场共识预期（consensus）比GS高5\%，这意味着如果Q1实际数据低于预期，股价可能面临进一步下调压力。

\subsection{[R010] GS：中国车企的“淘汰赛”还远未到终局}
{\small\color{refgray}归类：中国新能源车：6月订单透支7月，行业整合远未到终局}

当市场普遍认为中国新能源汽车行业的“淘汰赛”已经进入白热化阶段，甚至即将迎来终局时，GS在最新发布的研报中给出了一个反直觉的判断：行业整合，短期内不会发生。

这份由 Tina Hou 和 Jenny Du 主笔的报告，基于对14家主要中国车企的财务与经营数据的系统分析，提出了一个清晰的“行业触底框架”。结论是，尽管国内市场销量暴跌19\%，但强劲的出口增长和内部成本削减，使得行业整体的现金利润仍在增长。绝大多数车企的现金流仍然健康，管理层并未放弃扩张，而资产负债表上的净债务状态也远未普遍出现。

这意味着，我们可能正在经历一场“低烈度、长周期”的消耗战，而非一场速战速决的决战。对于投资者和产业决策者而言，理解这个框架，比猜测下一个出局者是谁更为重要。

KC评论： GS这份报告的核心价值不在于预测股价，而在于提供了一个可验证的“行业整合三条件”框架。它告诉我们，不要被短期的价格战和销量排名迷惑，真正决定行业格局的是现金流、管理层心态和资产负债表。这三点，才是判断行业何时见底的硬指标。

GS的报告提出了一个简洁的分析框架，认为只有当以下三个条件同时满足时，行业才会出现真正意义上的整合与出清：

*第一，大部分车企的EBITDA（税息折旧及摊销前利润）为负，即处于“现金成本”水平。** 这意味着企业连维持基本运营的现金都赚不回来，更无力进行新的投资或产能扩张。截至2026年第一季度，14家样本车企中，仍有9家保持着正的EBITDA。尽管相比2025年全年的11家有所减少，但多数企业依然“有血可流”。

*第二，管理层在利润率和扩张计划上全面“投降”。** 这指的是管理层放弃对盈利能力的追求，并停止或大幅缩减产能扩张计划。报告指出，当前管理层的态度是“混合”的。一方面，大家都承认国内需求疲软；另一方面，几乎所有车企都对海外增长寄予厚望，并仍在积极规划海外产能。这种“东方不亮西方亮”的心态，使得管理层远未到“投降”的时刻。

*第三，车企的资产负债表转为净负债。** 这是最硬性的财务指标。截至2026年第一季度，14家车企中仅有1家处于净负债状态。这意味着一方面，大多数车企仍有家底可以消耗；另一方面，即使是有亏损的企业，也可以通过股权融资等方式补充弹药，从而延长生存时间。

KC评论： GS这个框架的价值在于它把“行业出清”从一个模糊的概念，变成了三个可以定期跟踪的硬指标。目前来看，前两个条件（现金成本和净负债）在2027年之前都很难满足。这意味着，未来一到两年，我们大概率会看到的是“边打边谈”的持久战，而不是一两家巨头迅速吃掉所有市场的闪电战。

报告中的另一个关键信号是，市场对车企的盈利预期正在加速下调。2026年和2027年的EBITDA共识预期在过去一年中经历了四轮连续下调，最近一次下调幅度为4\%。

然而，GS同时指出，尽管预期在下调，但大多数车企的盈利水平仍然高于现金成本线。这形成了一个有趣的“剪刀差”：市场预期在快速变差，但企业的实际现金流状况还没有恶化到触发行业整合的临界点。

这种局面意味着什么？它意味着，当前的股价可能已经反映了部分悲观预期，但并未完全反映“持久战”带来的漫长煎熬。对于投资者而言，盈利预期的加速下调本身就是一个需要警惕的信号，它可能预示着未来几个季度，超预期的盈利下滑仍会不断出现。

在国内市场销量预计同比下降9\%的背景下，GS将2026年的出口增长预期大幅上调至30\%。出口已经成为支撑整个行业产量和批发量的唯一增长引擎。

但这里存在一个悖论。出口的增长，尤其是向利润率更高的海外市场出口，正在帮助那些在国内市场表现不佳的车企“续命”。它使得许多车企即使在国内亏损，也能通过海外利润实现整体正的EBITDA。

这恰恰是GS判断行业整合不会很快到来的核心原因之一。出口为整个行业提供了

[... middle omitted ...]

告对此没有给出明确的答案，这恰恰是完整报告中需要深入探讨的部分。

4. 报告尚未完全回答的关键问题：海外扩张的“成本”是多少？

GS的报告清晰地展示了海外扩张对销量的拉动，但一个尚未被完全量化的风险是：海外扩张的“成本”究竟有多高？

报告提到了“更多的车企正在考虑海外产能扩张”。这意味着，从出口整车到海外建厂，车企的商业模式正在发生根本性转变。海外建厂需要巨大的前期资本开支，这会显著增加企业的财务负担，并拉长投资回报周期。

目前，市场对这部分海外投资的资本开支和折旧摊销影响，似乎还没有形成清晰的共识。如果海外建厂的资本开支超出预期，或者海外市场的回报低于预期，那么出口带来的现金红利可能会被迅速侵蚀，反而加速部分企业的现金耗尽。

这个“第二层效应”，是判断行业整合时间表的关键变量，但它在公开的摘要中并未被充分展开。

基于GS的分析框架，对于希望跟踪行业动态的决策者而言，一个更务实的做法是，放弃对“谁最终会赢”的宏大叙事，转而建立一个更微观的“血条监测”框架。

\subsection{[R011] JPM：可持续投资正在“喘口气”，但真正的风险来自科技}
{\small\color{refgray}归类：区域市场：亚洲太空经济被低估，印度制造业成本驱动型修复 / 风险偏好：AI主题拥挤但需求无虞，可持续投资“喘口气”}

全球可持续基金总资产在今年4月达到了3.8万亿美元，同比增长3\%。这个数字看起来不错，但市场份额却从高点持续下滑，目前仅占全球基金总资产的5.6\%。JPM在最新发布的《可持续投资2026年中展望》中，给出了一个克制的判断：可持续基金的资金流出已经停止，但大规模资金流入的前提尚未满足。

过去两年，可持续投资经历了从狂热到退潮的完整周期。ESG主题基金遭遇了前所未有的信任危机，美国市场的“反ESG”政治浪潮、欧洲监管的不确定性、以及全球利率环境的变化，都让这个赛道从“必选项”变成了“可选项”。但JPM的这份报告试图告诉市场：最糟糕的阶段可能已经过去，新的结构性变量正在形成。

真正值得关注的不是资金流出的止跌，而是在科技股持续主导市场的背景下，可持续基金的结构性弱点正在被放大。这个问题的答案，将决定可持续投资的下一个十年。

JPM的数据显示，可持续股票基金的资金流已经从净流出转向中性。欧洲市场的流入抵消了亚洲和美国的流出，而固定收益类可持续基金则持续保持净流入。这是一个积极信号，但远未到乐观的时候。

报告明确指出，可持续基金要实现实质性的大规模资金流入，前提是相对业绩能够持续改善。而目前的情况是，可持续基金的整体表现仅仅是与大盘“基本持平”。在欧洲，可持续基金小幅跑赢欧洲大盘基金；在美国，可持续基金则小幅跑输。这种表现不足以吸引增量资金。

固定收益端的表现相对更好，可持续债券基金实现了正收益，且与整体市场基本一致。这解释了为何固定收益类可持续产品能持续获得资金流入。但债券市场的故事与股票市场不同——债券投资者的决策更多受到评级、久期和信用利差驱动，ESG标签的边际吸引力在下降。

KC评论： 资金流从“流出”到“中性”是一个重要的拐点信号，但不要误解为复苏的开始。可持续基金面临的根本问题是：它能否在科技主导的市场中证明自己的alpha生成能力。如果不能，资金流入的中性状态可能只是暂时的喘息，而非趋势反转。

这是整份报告中最值得深挖的部分。JPM的分析师指出，全球可持续基金普遍存在对科技板块的低配。原因很简单：可持续基金的地域配置偏向欧洲，而欧洲市场的科技股权重远低于美国。而恰恰是科技股，尤其是AI相关主题，正在成为全球表现最好的板块。

问题在于：这种“低配科技”的结构并非主动选择，而是可持续基金固有特征的结果。可持续基金的低配科技，不是因为它不看好科技，而是因为它的欧洲偏好在客观上导致了低配。当JPM的策略师明确看好AI主题（这更多是美国故事）时，可持续基金就陷入了一个两难境地。

报告进一步指出，这种风险比上一次展望时更加显著。JPM的新兴市场和日本策略师也加入了超配科技的行列，这意味着可持续基金在全球范围内都面临科技低配带来的相对业绩压力。

这里有一个更深的逻辑：可持续基金的传统优势板块——公用事业、工业、能源——恰恰是JPM策略师也看好的领域。但如果去掉科技板块，其他行业的观点大致相互抵消。也就是说，可持续基金的命运实际上被一个它天然低配的板块所决定。

KC评论： 可持续投资正在经历一个“身份困境”。它的投资框架天然倾向于低配科技，但市场却在奖励科技。这不是一个短期波动问题，而是可持续投资框架是否需要重新审视的问题。完整报告中对ESGQ因子的分析，提供了一个有趣的视角——它可能成为可持续基金的“隐形AI对冲”。

欧洲仍然是可持续投资最坚定的支持者。欧盟批准了到2040年较1990年减排90\%的目标，但报告敏锐地捕捉到一个关键的转变：政策工具正在从“大棒”转向“胡萝卜+大棒”。碳定价和监管仍在强化，但补贴、国家援助、公共采购、有针对性的保护主义和基础设施投资正在成为更重要的工具。这种“务实化”的转向意味着欧洲的可持续投资正在从理念驱动转向产业政策驱动。

美国的情况则更加复杂。报告区分了“噪音

[... middle omitted ...]

，中国的“十五五”规划设定了非化石能源占比25\%的目标，电网投资可能超过5万亿元人民币。更值得关注的是，中国的资产管理协会发布了A股市场的首个“尽责管理准则”，这被报告解读为“公司治理海啸”的延续。

报告的核心贡献之一是提出了六个可投资的可持续主题，覆盖全球、EMEA和美国市场。这些主题的平均潜在涨幅在20\%到51\%之间。

全球层面，三个主题分别是电网、电池供应链和气候变化适应。电网主题的平均涨幅预期为20\%，电池供应链为39\%，气候变化适应为30\%。EMEA市场则聚焦于“电主权”（Electro-Sovereignty）和股东积极主义，前者涵盖电气化、核能、汽车和电网，后者则受益于股东积极主义推动的价值释放，平均涨幅预期高达51\%。美国市场的主题是“低碳赢家”，预期涨幅为28\%。

这些主题的共同特征是：它们不再依赖ESG标签，而是基于产业趋势和政策驱动的确定性。电网投资、电池供应链、气候变化适应，这些都是全球资本正在涌入的方向，无论是否被贴上“可持续”的标签。

\subsection{[R012] 美国银行：芯片行业下一个1万亿美元，只需五年}
{\small\color{refgray}归类：AI/算力：从“光学狂欢”到“网络务实”，半导体拥挤度达历史极值}

半导体行业用了大约50年才实现第一个1万亿美元的累计销售额。美国银行（BofA）最新发布的研报给出了一个大胆判断：在AI的驱动下，第二个1万亿美元将在未来五年内完成。

这不是一个线性外推，而是一个结构性跃迁。该机构将2030年全球半导体行业总可寻址市场（TAM）从此前的2.3万亿美元大幅上调至2.7万亿美元，对应的年复合增长率从23\%提升至28\%。上调的核心驱动力并非消费电子复苏，而是AI基础设施从“论证投资回报率”阶段，全面转向“解决物理约束”阶段——芯片产能、电力供应、封装能力，正在成为新的稀缺资源。

这份报告发布于2026年6月，正值市场对AI资本开支持续性存在分歧的时刻。美国银行选择在这个时点大幅上调行业预测，其背后逻辑值得每一位产业决策者和投资者仔细拆解。

KC评论： 2.7万亿美元的TAM意味着什么？作为参照，2025年全球半导体销售额预计约为7870亿美元。即使按最乐观的预测，到2030年市场规模也要翻三倍以上。这不再是“成长”，而是“爆发”。报告的标题已经点明核心驱动力：AI将行业可见性延伸至2028年。如果你还在用传统半导体周期的框架来思考，很可能错过这轮结构性机会。完整报告中的情景分析和敏感性测试，能帮助你理解这个数字在不同假设下的弹性空间。

1. 从“论证ROI”到“解决物理约束”：AI产业的底层逻辑已变

过去两年，市场对AI的最大质疑是：巨额资本开支能否带来足够的回报。美国银行认为，这个阶段已经过去了。当前AI产业的核心矛盾，已经从“值不值得投”变成了“能不能建得出来”。

报告明确指出，AI行业正在从“必须证明投资回报率”转向“应对结构和物理约束”——芯片短缺、电力供应紧张、先进封装产能不足。这是一个根本性的叙事转变。当需求不再是问题，供给能力就成了定价权的来源。

这一判断也解释了为何美国银行大幅上调了晶圆厂设备（WFE）的预测。该机构将2028年WFE预测从此前的2030亿美元大幅上调至2500亿美元，上调幅度高达23\%。2027年WFE预测也被上调至1900亿美元，同比增长31\%。到2030年，WFE市场规模有望达到2920亿美元。

驱动WFE增长的因素包括：2028年洁净室可用性的大幅提升、存储器长期协议提供的2-3年供需可见性，以及先进制程和封装技术的复杂性持续推高单晶圆设备支出。

KC评论： 设备是芯片的“卖铲人”。美国银行对WFE的激进上调，本质上是在押注AI芯片的需求会迫使晶圆厂持续扩产，而扩产就意味着设备订单。但报告也隐含了一个重要假设：洁净室和产能建设能够按计划推进。如果地缘政治或供应链扰动导致建设延迟，WFE的实际增速可能低于预期。完整报告中对英特尔、三星等主要晶圆厂的产能进度分析，是验证这一假设的关键。

美国银行预测，到2030年，存储器市场规模将从2025年的2200亿美元飙升至1.65万亿美元，年复合增长率高达49.6\%。相比之下，核心半导体（非存储器）的同期年复合增长率仅为13.9\%。

这一差异意味着，未来五年半导体行业的新增销售额中，存储器将贡献绝大部分。驱动存储器增长的核心逻辑有两个：

第一，AI系统中存储器的单位价值增速超过了计算芯片。HBM（高带宽存储器）的渗透率持续提升，单颗AI芯片需要搭配的HBM容量和带宽都在快速增长。美国银行在报告中特别指出，美光（MU）的目标价从950美元大幅上调至1500美元，核心依据就是HBM业务的估值提升。

第二，存储器供给弹性结构性下降。与过去几年不同，存储器厂商现在更倾向于通过长期协议锁定需求和价格，而非盲目扩产。美光与Anthropic的合作就是典型案例。这种供给纪律意味着存储器价格周期的波动性可能降低，头部厂商的盈利稳定性将显著增强。

美国银行对2026年的终端市场预测清晰揭示了“分化

[... middle omitted ...]

计将在2026年进入补库周期。汽车市场虽然单位销量增长缓慢，但单车芯片含量的提升（尤其是智能驾驶和电动化）正在为相关芯片公司提供结构性增长。

4. 报告尚未完全回答的关键问题：AI资本开支的持续性如何验证？

美国银行将行业可见性延伸至2028年，但报告本身并没有提供一套严格的框架来验证AI资本开支的持续性。这是当前市场最大的分歧点，也是这份报告最值得追问的地方。

目前支撑高资本开支的核心假设是：AI计算需求“永不满足”。但这一假设面临几个潜在挑战：第一，AI应用能否在消费端产生足够的经济效益，以支撑云服务商的持续投入？第二，推理效率的提升是否会降低单位计算需求？第三，地缘政治风险是否会导致关键供应链中断？

美国银行在报告中提到了“物理AI”和“Agentic CPU”等长期驱动力，但并未对这些场景下的需求弹性进行深入拆解。对于希望独立验证AI资本开支持续性的读者，完整报告中对云服务商资本开支趋势、芯片利用率以及HBM定价模型的详细分析，提供了更底层的数据支撑。

\subsection{[R013] GS：新能源车周订单暴增34\%，但真正的胜负手不在销量}
{\small\color{refgray}归类：中国新能源车：6月订单透支7月，行业整合远未到终局}

6月第三周，中国新能源车市场交出了一份让多数人意外的高分答卷。GS最新发布的周度追踪数据显示，在6月15日至21日这一周，主要新能源品牌合计订单同比大增31\%，环比更是飙升34\%。比亚迪、小米、零跑成为环比增长最快的三家车企，增幅分别达到134\%、44\%和22\%。

这份数据发布的时间点值得注意。市场此前普遍担忧价格战持续侵蚀利润，叠加传统燃油车在6月加大折扣力度，新能源车渗透率的上升势头可能遭遇阻力。但GS的周度数据给出了相反的信号：新能源车零售渗透率在6月前两周已达到63.9\%，批发渗透率更是攀升至67.9\%，均高于5月全月的水平。

订单的爆发式增长，最直接的驱动力是新车型集中上市。比亚迪“大唐”于6月17日开启交付，零跑C10/C11/C16改款于6月16日上市。这些新车型在上市首周即贡献了显著的增量订单。但问题在于：这种由新品脉冲驱动的增长，能否持续？订单转化率如何？价格战是否真的在缓和？

GS这份周度图表报告提供了大量高频数据，但真正值得产业决策者和投资者关注的，不是某一家公司单周订单的起伏，而是这些数据背后隐藏的结构性变化。

KC评论： GS这份周报的核心价值不在于告诉你“哪家卖得好”，而在于它提供了从订单、折扣、渗透率到上游电池价格的完整高频拼图。对于跟踪竞争格局的人来说，单周数据是噪音，但连续数周的趋势变化，往往比季报更能提前暴露拐点。完整报告中的逐周对比图表和品牌级折扣数据，是判断“涨价是否可持续”的关键素材。

1. 这轮订单脉冲的含金量，取决于新车型能否跨越“尝鲜期”的鸿沟

比亚迪“大唐”上市首周贡献了惊人的订单增量，但仔细看GS的数据，比亚迪王朝与海洋系列YTD订单同比仍下降40\%。这意味着，新车型带来的增量，目前更多是在弥补老车型的订单流失，而非整体市场的净扩张。

零跑的情况类似。其改款车型上市后周订单环比增长22\%，但YTD订单数据缺失，难以判断其增长是否可持续。真正值得关注的是蔚来和小米。蔚来YTD订单同比增长96\%，在所有品牌中表现最为“抗跌”；小米虽然YTD订单同比下降33\%，但6月第三周订单环比增长44\%，显示出SU7交付爬坡后需求仍在恢复。

这里的关键问题在于：新车型带来的订单脉冲，通常会在上市后2-4周达到峰值，随后逐步回落。GS报告并未给出新车型订单的转化率数据，也没有区分“新增订单”与“老车型流失订单”的净效果。对于投资者而言，需要关注的是7月初各品牌公布的6月交付量，届时才能判断这轮脉冲有多少转化为了实际交付。

KC评论： GS这份周报最宝贵的地方，在于它提供了品牌级别的周度订单数据，这在其他公开渠道几乎无法获得。但报告没有回答一个关键问题：这些订单有多少是“真实新增”，多少是“老车型订单转移”？对于比亚迪这样的巨头，如果新车型只是分流了自家老车型的客户，那对整体利润率的提升就非常有限。完整报告中的月度订单对比表和YTD同比数据，是判断这一问题的起点。

2. 新能源车折扣在收窄，但燃油车折扣在扩大——这个分化比销量数字更重要

GS的数据显示了一个容易被忽略但意义深远的变化：截至6月20日，新能源车经销商平均折扣率为7.48\%，低于6月13日的7.54\%，更低于去年同期的8.15\%。而同期燃油车经销商平均折扣率从19.56\%扩大至19.62\%，且远高于去年同期的22.95\%（折扣率越高，意味着终端降价幅度越大）。

这个分化意味着什么？新能源车行业的“价格战”可能正在进入一个新阶段：头部品牌开始有能力收窄折扣，而燃油车为了守住市场份额，不得不进一步加大优惠力度。比亚迪的折扣率稳定在4.27\%，远低于行业平均，说明其定价权在增强。

但这并不意味着新能源车行业已经走出了价格战的泥潭。折扣率的收窄，更多是结构性因素导致：新车型上市初期通常折扣较小，而新车型在订单中的占比提升，拉低了整体折扣率。一旦新车型进入“走量期”，折扣是否还会重新放大，仍需观察。

KC评论： 折扣数据是判断行业利润拐点的先行指标。GS报告中的品牌级折扣趋势图，比整体平均数更有价值——它揭示了谁在主动涨价、谁在被动跟降。完整报告中的比亚迪、特斯拉、蔚来等品牌的逐周折扣曲线，是判断各品牌定价权强弱的关键证据。

3. 上游锂价再次承压，但电池成本传递到整车价格存在明显时滞

GS数据显示，电池级碳酸锂价格跌至16.75万元/吨，周环比下降4.0\%。而方形LFP电芯和方形NCM电芯价格周环比持平。

锂价下跌对整车企业来说本应是利好，但现实是：上游原材料价格波动传递到下游整车成本，通常需要2-3个季度。这意味着，2025年下半年的锂价上涨压力，正在体现在当前的新能源车成本中；而当前锂价的下跌，要到2026年下半年才能对整车成本产生显著影响。

更重要的是，电芯价格已经连续多周持平，并未跟随锂价下跌。这说明电池厂商在通过锁定利润来对冲上游波动，整车企业短期内难以享受到原材料降价的红利。对于毛利率本就承压的新能源车企来说，这意味着2026年下半年的成本压力可能比市场预期的更大。

KC评论： GS报告中的上游电池定价数据，表面上只是一个价格跟踪，但结合折扣率和订单数据，可以构建一个完整的“成本-定价-销量”分析框架。完整报告中的历史价格对比表和季度均价数据，是判断电池成本趋势是否出现拐点的必要素材。

\subsection{[R014] JPM：地产股急跌不是基本面反转，而是噪音窗口}
{\small\color{refgray}归类：中国消费与地产：K型分化加速，定价权成为核心考验}

过去四个交易日，中国地产板块跑输恒生指数9\%。对于持有国企开发商仓位的投资者来说，这轮急跌像是突如其来的冷水。但仔细拆解JPM这份最新报告可以发现，市场关心的核心问题——房价是否企稳、销售是否可持续——并没有被高频数据否定。

相反，报告给出的判断值得认真对待：这轮下跌主要是四重噪音的叠加，而非基本面趋势的逆转。真正关键的变化藏在两个被市场低估的信号里——一线城市二手挂牌量的持续下降，以及房价环比连续第四个月转正。这些信号指向的结论是：市场正在从“全面下跌”切换到“K型分化”，而分化的赢家已经清晰。

*KC评论：** 这份报告最核心的判断是，过去几天的急跌不是基本面的“第二只脚”，而是噪音。但噪音本身也值得认真对待——它告诉我们市场对什么信号敏感、什么变量正在被错误定价。完整报告里的高频数据拆解和个股估值表，能帮你判断哪些资产在这轮噪音中反而值得加仓。

JPM把过去四天的下跌归结为四个原因，每一个都指向“情绪和结构因素”，而非销售或价格数据的恶化。

第一，消费数据的疲弱产生了负面传导。当居民消费信心走弱，市场自然怀疑房地产复苏的可持续性。第二，地产板块的高贝塔属性在市场整体回调中被放大——同期恒生指数本身就在下跌，地产只是跌得更快。第三，短期缺乏催化剂。高频数据虽然“intact”（完好），但没有出现“明显更好”的边际改善，这让市场缺少一个买入的理由。第四，部分投资者过度关注了二手成交的环比下滑，而忽视了季节性因素——3月和4月通常是成交旺季，环比回落是正常现象。

报告特别指出，即使是华润置地和中海地产这样此前相对抗跌的国企龙头，也在过去两个交易日急跌了11\%，而同期恒生指数仅跌2\%。但即便如此，这两只股票年初至今依然分别跑赢恒生指数22和14个百分点。

KC评论： 这四重噪音中，最容易被误读的是第四点——环比数据。JPM明确说，用环比判断市场健康度是“不准确的”，同比才是更合理的指标。但市场上大量交易型资金和量化模型恰恰在用环比做决策。完整报告里关于季节性调整后的同比数据对比，能帮你建立更可靠的观测框架。

报告中最值得关注的图表之一，是一线城市二手挂牌量的冰山指数变化。数据显示，自3月峰值以来，一线城市二手挂牌量已经持续下降了2.5\%。

这个信号之所以重要，是因为供给端的收缩是房价企稳的必要条件。过去两年，中国房地产市场的核心矛盾一直是“卖盘过多、买盘不足”。挂牌量的持续下降意味着，一线城市正在从“买方市场”向“均衡市场”过渡。如果这个趋势能够持续，二手房价的环比转正将从“脉冲式修复”升级为“趋势性企稳”。

JPM明确表示，这个因素“likely under-appreciated by the market”——很可能被市场低估了。

*KC评论：** 挂牌量下降2.5\%听起来不大，但放在一线城市动辄数万套的存量池里，边际影响不可忽视。更重要的是，这个趋势是否可持续，取决于两个变量：一是价格预期是否已经触底，二是持有成本是否稳定。完整报告里对冰山指数的跟踪是周度的，能帮你判断这个趋势是在加速还是在放缓。

高频数据中，市场最担心的是“上周同比增速从+12\%骤降到+1\%”。但JPM的解读非常直接：这个下降主要来自端午节假期造成的基数效应。如果对比去年同样包含端午节的那一周，同比增速仍然是+15\%，与此前数周10-20\%的区间一致。

官方网签数据也呈现同样的模式。60城一手房网签同比下降23\%，但同样是因为假期因素；如果只对比端午节前三天的数据，同比增速是+60\%。12城二手房网签同比下降13\%，结束了连续9周的同比正增长，但JPM预期下周就会看到明显的环比和同比改善——因为网签数据本身有“季末冲量”的季节性特征。

报告强调：同比增速是评估市场趋势的更准确指标。而同比数据依然在10-20\%的

[... middle omitted ...]

地产和金茂这三家国企开发商，因为它们同时具备“销售增长跑赢同行”和“强一线城市敞口”两个特征。

而对于万科、融创等非国企开发商，报告维持谨慎态度，认为它们“可能无法从K型稳定中受益”。

KC评论： 这里有一个关键问题报告没有完全展开：K型分化的“K”会持续多宽？如果一线城市房价继续企稳但二三四线仍在探底，国企开发商的估值溢价能维持到什么水平？完整报告中的个股估值表和P/E对比，能帮你量化这个溢价是否已经price in。

JPM把“消费数据弱化”列为下跌原因之一，但报告并没有深入展开这个逻辑链。消费信心走弱如何传导到地产销售？是通过影响首付能力、还是通过影响购房意愿？如果消费数据持续恶化，一线城市已经企稳的房价会不会再次承压？

这可能是当前地产板块最大的不确定性。过去三个月的房价企稳，很大程度上建立在“政策刺激+积压需求释放”的基础上。如果居民收入预期和消费信心进一步走弱，释放的需求可能难以持续。报告对此没有给出明确判断，而是选择聚焦于高频数据的“事实层面”。

\subsection{[R015] 摩根斯坦利：人形机器人2026年出货量上调近八成，关键不在技术而在“部署数据闭环”}
{\small\color{refgray}归类：风险偏好：AI主题拥挤但需求无虞，可持续投资“喘口气”}

摩根斯坦利：人形机器人2026年出货量上调近八成，关键不在技术而在“部署数据闭环”

2026年6月，摩根斯坦利发布了一份关于中国人形机器人产业的深度研报，核心判断清晰且直接：商业化验证、政策支持与供应链反馈正同步加速，该机构将2026年中国人形机器人出货量预测从此前的2.8万台大幅上调至5万台，2030年预期达到44.6万台。

这不是一份关于“机器人能否造出来”的报告。报告的核心主张是：产业正在从“展示能力”进入“创造商业价值”的阶段，而决定这一轮成败的关键，不是单台机器人的技术参数，而是谁能最快通过规模化部署，形成“部署-数据采集-算法迭代-再部署”的正向循环。

对于关注中国制造业升级、供应链重构以及AI硬件落地的读者，这份报告提供了一个从宏观政策到微观订单、从整机厂到核心零部件供应商的完整观察框架。它既指出了明确的投资主线，也留下了需要持续跟踪的关键变量。

KC评论： 摩根斯坦利上调出货量预测的直接驱动因素，不是实验室的技术突破，而是国家电网一笔68亿元订单、顺丰和邮政的物流场景落地、以及MIIT与SASAC联合推动的“万级部署能力”目标。这意味着，人形机器人已经从“概念验证”进入了“政策+订单”双轮驱动的阶段。完整报告中对这些订单的细节、供应链验证的现场反馈、以及各环节公司的产能规划都有更细致的拆解，值得细读。

1. 这轮加速的真正推手不是技术，而是“用起来”的政策与订单

报告指出，2026年上半年，人形机器人的商业化进程明显提速，标志性事件密集出现。

政策层面的信号最为明确。2026年3月，中央政府首次在“十五五”规划中将机器人列为战略性新兴产业。6月，工信部和国资委联合发布新倡议，目标是在2026年底前形成“万级部署能力”，并落地超过100个高价值应用场景。具体要求是：每个省份至少确定20个、每个中央国企至少确定10个真实运营场地，将其改造为人形机器人训练环境。11月，两部委将对进展进行评估。

订单层面，国家电网一笔68亿元人民币的采购订单成为行业里程碑——采购内容包括500台人形机器人、3000台双臂机器人和5000台四足机器人。顺丰和邮政已经在物流中心部署Robotera的人形机器人。这些不是小规模试点，而是带有明确商业目标的规模化采购。

报告认为，商业化验证通常需要数月时间，2026年上半年启动测试的项目，预计将从下半年开始转化为实际部署。这意味着5万台的出货量预测并非激进假设，而是对已有订单和验证进度的合理外推。

2. 供应链已经在为“万台级”做准备：产能扩张与设备订单先行

供应链的反馈是验证产业趋势最扎实的“硬数据”。摩根斯坦利在近期调研中发现，上游零部件供应商的产能规划已经远远超过当前出货量，说明产业链对需求爆发的预期是真实且一致的。

谐波减速器龙头绿的谐波（Leaderdrive）的月产能已从2026年一季度的5万台提升至当前的约7万台，年底目标为10-12万台。其产能规划对标的是日本哈默纳科（Harmonic Drive）明年约200万台的年产能计划。行星滚柱丝杠供应商恒立液压（Hengli Hydraulic）的墨西哥工厂目标是在年底前具备支撑约10万台机器人的产能。信捷电气（Xinje）计划今年形成约200万套无框力矩电机产能，明年扩至约300万套。

更值得关注的信号是，产能扩张已经开始向上游设备端传导。先导智能（Wuxi Lead）已与北京人形机器人创新中心签署战略合作协议，涉及未来装配线设备。拓斯达（Topstar）表示，人形机器人已占其数控机床下游需求的30-40\%。这意味着，设备供应商已经收到了明确的扩产订单，产业链的投资正在从“零部件”向“生产零部件的工具”扩散。

KC评论： 供应链的产能规划往往比终端出货量领先6-12个月。当一家减速器企业将

[... middle omitted ...]

据采集、政府项目和互动型商业服务。

但报告预期，随着验证推进和全尺寸机型在操作能力上的提升，其占比将在2027年升至约50\%，2028年达到约70\%。这一结构性变化对供应链的影响是直接的：全尺寸机型需要更大扭矩的关节模组、更高精度的减速器、更复杂的灵巧手，其ASP（平均销售价格）显著高于半尺寸机型。

报告预计，2026年行业综合ASP将下降15\%，但2027年和2028年将分别回升5\%和2\%——这看似反直觉，实则反映了产品结构从“低价半尺寸”向“高价全尺寸”的切换。这意味着，即使出货量增速放缓，市场规模的扩张速度可能更快。

4. 优选标的的逻辑：谁在“部署-数据闭环”中卡住了关键位置

摩根斯坦利在这份报告中更新了其中国人形机器人价值链条，新增科达利（Kedali）和蓝思科技（Lens），剔除国茂股份（Guomao）、旭升集团（Xusheng）和中坚科技（Zhongjian），目前共覆盖45家公司，包括3家“大脑”公司、32家“身体”组件公司和10家整机集成商。

\subsection{[R016] 摩根斯坦利：MS：中国房地产的修复，还差一个“端午效应”的量级}
{\small\color{refgray}归类：中国消费与地产：K型分化加速，定价权成为核心考验}

摩根斯坦利：MS：中国房地产的修复，还差一个“端午效应”的量级

这份MS最新发布的每周数据库追踪报告，披露了一个关键信号：在刚刚过去的一周（截至6月21日），中国50城新房周度成交量同比下降16\%，较前一周的-3\%明显走弱。这份报告的核心判断不是“转冷”，而是“被节日效应掩盖的真实节奏正在显现”。

真正值得关注的，是数据背后所揭示的市场修复的结构性特征——不是全面回暖，而是分化的加速。这份研报提供了一组精细到城市层级、一二手市场、挂牌价与去化率的交叉数据，足以让产业决策者重新审视“修复”二字的真正含义。

如果只看周度同比数据，新房市场似乎出现了一个明显的回调。50城新房成交量同比下降16\%，而前一周还是-3\%。但MS明确指出了这背后的关键扰动因素——端午节假期带来的日历效应。

这并非简单的技术性解释。端午节期间，各地推盘节奏明显放缓，网签系统暂停或延迟，这些都会导致当周成交数据出现系统性偏低。更值得追问的是：如果剔除节日效应，真实的市场动能是什么？

答案是：修复仍在进行，但强度远不足以支撑全面反转。前一周的-3\%已经是一个偏弱的读数，而本周的-16\%即使进行日历调整，也很难回到正增长区间。这意味着，当前市场的修复更多是“低位企稳”而非“趋势性回暖”。

KC评论： 对于投资者而言，不要因为单周数据的大幅波动就改变对行业基本面的判断。节日效应在房地产周度数据中是一个常见的噪声源，但噪声背后往往隐藏着更重要的信号——即市场在没有政策刺激和假期干扰时，其自然成交动能到底有多强。完整报告中的四期移动平均线图表，才是判断趋势的真正工具。

分城市层级来看，数据呈现出一个令人不安的分化格局。一线城市新房成交同比下降15\%，较前一周的+2\%明显转弱；二线城市同比下降12\%，前一周为-7\%；三线城市更是骤降至-28\%，而前一周还是+9\%。

表面上看，一线城市似乎跌幅最大，但考虑到其前期基数较高（前一周还是正增长），这种回落更多是均值回归。真正值得警惕的是三线城市的剧烈波动——从+9\%到-28\%的跳变，说明低线城市的市场情绪极其不稳定，任何短期因素都能引发成交量的剧烈震荡。

而二线城市的表现则揭示了一个更本质的问题：当一线城市的热度无法有效传导至二线城市时，整个市场的修复链条就断了。二线城市周度12\%的同比降幅，意味着承接一线城市外溢需求的逻辑并未兑现。

KC评论： 这份报告最有价值的洞察之一，就是城市层级之间的传导失灵。过去市场修复的逻辑是“一线带动二线，二线辐射三线”，但现在一线城市的修复更多是内生性的（高净值人群的改善需求），其对周边城市的拉动效应正在减弱。完整报告中的城市层级细分数据，能够帮助读者判断哪些城市真正受益于这轮修复，哪些只是被动跟随。

二手房市场的数据同样不容乐观。10城二手房周度成交量同比下降15\%，较前一周的+4\%明显下滑。其中一线城市同比下降10\%，前一周为+3\%；二线城市同比下降19\%，前一周为+3\%。

这里有一个容易被忽视的细节：二手房市场的年初至今累计成交量仍为+5\%，这意味着虽然周度数据在走弱，但整体趋势并未逆转。更关键的是，二手房挂牌价的压力正在累积——中央指数六城二手房挂牌价跟踪指数为17.3\%，较前一周的17.4\%略有下降，但一线城市房产中介指数从53.2骤降至51.0。

中介指数的下降，往往意味着议价空间在扩大、成交周期在拉长。这对于急于出手的房东来说，是一个明确的信号：挂牌价的“韧性”正在被成交量下滑所侵蚀。

KC评论： 二手房市场是判断房地产市场真实温度的最佳指标。新房成交受推盘节奏、网签政策、开发商促销等因素影响较大，而二手房成交更多反映的是买卖双方的真实意愿。当前二手房成交量的下滑与挂牌价压力的上升，说明市场仍未形成“量价齐稳”的正循环。完整报告中的二手房

[... middle omitted ...]

的吸引力，而二线城市的低去化率则意味着，非核心区域的项目正在被市场抛弃。

这种分化已经不仅仅是城市层级之间的差异，而是同一城市内部不同板块、不同产品之间的分化。在一个整体去化率只有47\%的市场中，能够实现59\%去化率的项目，必然具备某种稀缺性——地段、产品力、品牌溢价或价格优势。

KC评论： 去化率的分化是这轮修复最核心的特征之一。它告诉投资者：不是所有的房地产股票都值得投资，不是所有的项目都能卖得动。那些能够持续实现高去化率的开发商，才是真正的赢家。完整报告中的城市层级去化率数据，以及MS对重点开发商的评级，能够帮助读者构建一个更精准的投资框架。

MS在这份周度数据追踪报告中，提供了非常精细的成交与价格数据，但有一个问题它没有直接回答：当前的市场修复，在多大程度上依赖于政策支持？

从历史经验来看，中国房地产市场的每一次修复都离不开政策的推动。但政策的边际效用正在递减——过去一次降息或放松限购就能带动市场反弹，现在需要更密集、更精准的政策组合才能维持市场稳定。

\subsection{[R017] UBS：中国券商股不只是反弹，是估值体系正在被政策重写}
{\small\color{refgray}归类：未被主题摘要引用}

6月22日，中国券商H股板块单日大涨6.1\%，同期恒生指数下跌0.65\%。这不是一次简单的情绪修复。UBS在当天发布的研报中给出了一个值得认真对待的判断：券商板块正在进入一个由政策组合拳和基本面共振驱动的重新定价窗口。

这份报告的核心主张是，市场低估了两件事——第二季度业绩的强劲程度，以及政策层面正在发生的结构性变化。前者是短期催化剂，后者才是估值中枢上移的真正来源。

如果你还停留在“券商股是牛市放大器、熊市拖累者”的旧框架里，这篇研报可能会让你重新审视这个行业的定价逻辑。

UBS的判断起点是业绩。这不是一个模糊的“有望改善”，而是有数据支撑的强信号。

报告指出，A股二季度至今的股票日均成交额已跃升至2.8万亿元人民币，同比增长1.2倍，环比增长9\%。融资融券余额升至3.0万亿元，同比增长62\%，环比增长14\%。A股IPO募资额同比增长65\%，环比增长91\%。新发混合型基金规模达到960亿元，同比增长1.8倍。4至5月新开股票账户数达530万户，同比增长51\%。

这些数字合在一起意味着什么？券商的核心收入来源——经纪、两融、投行、资管——在第二季度几乎全线走强。而市场目前对券商股的定价，显然还没有充分反映这一现实。

以CITICH股为例，其年内涨幅仅为5.4\%，CICCH股年内涨幅11\%。对于一个营收增速可能在30\%以上的行业，这样的估值水平并不算贵。

KC评论： 这里的关键不是“业绩好”，而是“业绩好但股价还没充分反映”。UBS的测算框架显示，上市券商一季度经常性净利润已同比增长39\%，二季度大概率会更高。但板块年内H股整体下跌0.4\%，A股下跌6.8\%，与恒生指数下跌7.3\%相比只是小幅跑赢。这意味着，如果业绩兑现，估值修复的空间是真实的。完整报告中对各业务线的拆解和盈利预测，值得仔细看。

报告特别强调了上周陆家嘴论坛的政策信号。这不是一次简单的“安抚市场”，而是一个系统性的制度框架调整。

UBS提炼了三个层次：第一，市场稳定基调明确，吸引长期资金、扩大金融工具（包括活跃ETF）等举措，本质上是为市场结构加固；第二，对科创板和创业板的上市标准做出更包容的调整，重点聚焦AI和深科技领域，这意味着投行业务的增量空间被打开；第三，跨境层面，推进人民币外汇期货试点，支持香港推出5年期国债期货，同时通过沪深港通、QDII、跨境理财通等渠道引导资金流动。

这三个层次叠加在一起，指向一个结论：监管层正在用制度设计推动券商从通道型业务向买方投顾模式转型。这不是一个短期的政策刺激，而是一个中长期的结构性利好。

KC评论： 市场往往把政策解读为“利好出尽”或“短期情绪”，但UBS看到的是一套组合拳。比如，外资机构获得国债期货准入和基金投顾牌照，短期内可能加剧竞争，但长期看会加速行业向买方投顾模式转型——这对头部券商是利好，对中小券商是压力。完整报告中对政策影响的量化估算，是理解这个判断的关键。

UBS在报告中明确表达了偏好：CICCH股、CITICA/H股、HTSCA/H股、GTJAA股。这不是随机的选股，而是基于一个清晰的逻辑——行业集中度正在加速提升，头部券商的竞争优势正在被制度性放大。

报告提到，领先券商在科创板业务中具有显著优势，通过跟投承诺和私募股权持仓，能够实现投行、研究、投资的一体化协同。这种能力不是所有券商都具备的。同时，监管对并购重组、再融资的支持，以及对符合条件的港股上市公司回归A股上市的政策，都将使头部券商受益更多。

UBS的估值框架显示，其覆盖的券商平均交易在2026年预测市净率的1.2倍（H股）和0.8倍（A股）。对于一个盈利前景改善、结构性增长动力明确的行业，这个估值水平并不高。

但这里有一个关键问题：头部券商的优势能否持续转化为定价权和市场份额？这需要更细致的拆解。


[... middle omitted ...]

股市场出现持续调整，券商的收入结构会迅速承压。UBS报告中的一个细节值得注意：虽然IPO募资额大幅增长，但再融资承销额同比下降45\%。这意味着投行业务的结构正在发生变化，对头部券商的综合能力提出了更高要求。

KC评论： 报告没有完全回答的问题是，这轮政策驱动的增长周期，是否真的能改变券商行业的周期性特征？如果市场再次进入低迷期，头部券商能否凭借买方投顾和资本中介业务维持盈利韧性？这些问题的答案，需要在完整报告中对各家券商的业务结构和盈利质量做更深入的拆解。

基于UBS的分析框架，我们可以提炼出一个观察券商股的三个维度：

第一，市场活跃度能否维持。日均成交额、两融余额、新开户数这些领先指标，是判断业绩能否兑现的最直接窗口。如果这些指标出现拐点，券商的业绩预期需要同步调整。

第二，政策落地速度。陆家嘴论坛释放的信号需要转化为具体的规则和实施细则。外资机构获得国债期货准入和基金投顾牌照后，行业竞争格局会如何演变？头部券商能否在买方投顾转型中占据先机？这些都需要跟踪。

\subsection{[R018] UBS：氟化工正被AI重新定价，但市场还没看懂}
{\small\color{refgray}归类：中国化工：传统品疲弱，氟化工与MDI的结构性机会}

当AI算力竞赛从芯片蔓延到材料，一个被低估的板块正在浮出水面。

UBS最新发布的氟化工行业报告，给出了一个值得产业决策者和投资者认真审视的判断：氟化工材料在AI领域的应用潜力正在快速增长，但市场对这一结构性变化的定价仍然不足。报告明确指出，“氟化工材料供应商目前相对于电子化工材料企业存在显著折价，我们认为市场尚未完全定价AI带给氟化工材料的增长机会。”

这不是一个关于“氟化工涨价”的短期交易故事。这是一个关于材料体系迭代、供应链重构和国产替代交汇的长逻辑。

这份报告发布于2026年6月，覆盖了从电子级氢氟酸到氟化液、从PTFE到氟化气体的完整产业链。但它的核心洞察其实非常集中：AI对算力的需求正在倒逼材料体系的升级，而氟化工材料在性能上的不可替代性，使其成为这一进程中的关键受益者。

如果你只是把氟化工看作一个周期性化工子行业，你可能正在错过一个结构性重估的窗口。

1. AI对氟化工的需求不是“蹭概念”，而是材料体系的必然迭代

UBS报告开篇就点明了这一轮氟化工指数跑赢化工指数的两大驱动力：一是美国推迟部分第三代制冷剂退出时间表，二是市场对氟化工材料AI应用潜力的预期升温。但真正值得关注的，是后者。

AI服务器对散热效率和低介电常数提出了远超传统服务器的要求。当芯片功耗从几百瓦跃升至上千瓦，空气冷却已经达到物理极限，液冷成为必然选择。而在液冷方案中，氟化液凭借其化学惰性、不燃性、优异的电绝缘性和低表面张力，成为目前最理想的冷却介质之一。

同时，AI服务器对PCB基板材料的介电性能提出了更高要求。PTFE（聚四氟乙烯）因其极低的介电常数和介电损耗，被视为下一代PCB系统的候选材料。

这不是“概念炒作”，这是算力密度提升带来的物理约束倒逼材料升级。UBS的报告把这一点讲得很清楚：AI发展带来的更高算力需求和材料体系迭代，是氟化工材料需求增长的底层驱动力。

KC评论： 很多人看到“氟化工+AI”的第一反应是“又在蹭概念”。但这份报告给出的逻辑链是：算力密度提升 -> 散热方案从风冷到液冷 -> 氟化液成为关键材料。这不是想象，是物理规律。完整报告里有两张产业链图，清晰地展示了氟化工与AI从芯片制造到数据中心的全链条连接，值得细看。

如果说AI需求是需求侧的“拉力”，那么3M退出PFAS生产则是供给侧的“推力”。两个力量叠加，构成了中国氟化工企业的结构性机遇。

UBS报告指出，3M作为全球最大的半导体级氟化液供应商，于2025年底退出了PFAS生产。这一决定源于全球范围内对PFAS（全氟和多氟烷基物质）日益严格的监管。PFAS因其高持久性、不可降解性和潜在环境健康风险，正从单一产品限制转向更广泛的类别管控。

问题在于，氟化液在半导体制造中几乎不可替代。尤其是HFE（氢氟醚）和PFPE（全氟聚醚），它们用于精密电子元件的清洗和散热，在芯片制造过程中直接影响工艺稳定性和良率。海外监管机构正在考虑对电子和半导体等关键应用给予更长的豁免期，但这并不能解决供给短缺的问题。

这就给中国供应商创造了替代窗口。UBS报告明确提到，Capchem（新宙邦）在半导体级氟化液领域占据了有利位置，已与国内主要晶圆厂建立了稳定合作关系，在国内半导体冷却液市场排名第一。

KC评论： 3M退出PFAS这件事，很多人知道，但很少有人算清楚它对产业格局意味着什么。UBS这份报告给出了一个关键数据：半导体和汽车业务贡献了3M PFAS营收的35-40\%。这部分市场突然空出来，谁能接得住？完整报告里有一张中国主要上市公司氟化液产能布局表，可以看到谁在真金白银地扩产，谁还在“准备中”。

3. G5级电子级氢氟酸的国产替代，正在从“能做”走向“能打”

在半导体制造环节，电子级氢氟酸（EG-HF）是湿电子化学品中用量最大的品类之一

[... middle omitted ...]

9万吨，营收6.58亿元。更重要的是，国内G5级供应商的CR5高达98\%，意味着头部集中度极高。

多氟多、巨化股份旗下的凯圣氟化、滨化集团是主要的国内供应商。多氟多已向台积电、三星批量供应半导体级氢氟酸；凯圣氟化获得了SK海力士12英寸1Xnm工艺的批量供应确认。

但UBS报告也指出了关键变量：中东冲突导致全球硫磺供应趋紧，推动无水氢氟酸价格上涨，进而推高了EG-HF的提价预期。中国作为全球最大的无水氢氟酸供应国，国内EG-HF生产商在原材料供应稳定性上具有优势，有望进一步扩大全球市场份额。

UBS报告给出了明确的投资结论：偏好能够受益于AI驱动需求增长的氟化工材料生产商，如东岳集团、新宙邦、广钢气体等，并大幅上调了目标价。

但有一个问题，报告没有完全展开：当前氟化工材料企业的估值折价，何时会收敛？

UBS报告提到，氟化工材料供应商目前相对于电子化工材料企业存在显著折价。这是一个有趣的现象——如果市场已经认可AI对氟化工材料的拉动逻辑，为什么估值还没有跟上？

\subsection{[R019] GS：中国经济5月边际回暖，但内需疲软与信用紧缩的“温差”才是真正挑战}
{\small\color{refgray}归类：未被主题摘要引用}

GS：中国经济5月边际回暖，但内需疲软与信用紧缩的“温差”才是真正挑战

中国经济在5月出现了一丝令人宽慰的边际改善，但这份改善的底色，远非“复苏”二字可以概括。GS最新发布的6月中国自营经济指标报告揭示了一个更复杂的图景：制造业产出带动了整体活动指标的微幅回升，然而，隐藏在数据背后的内需疲软、信用紧缩和财政节奏滞后，构成了当前经济运行的真正挑战。对于投资者和决策者而言，理解这组“温差”信号，比关注单一数据点的好坏更为关键。

这份报告的核心价值不在于给出了一个乐观或悲观的结论，而在于它通过一系列自营的高频和结构性指标，拆解了经济增长的动力来源与阻力所在。它提醒我们，当前的经济修复更像是一次局部反弹，而非全面的趋势反转。

GS的中国当前活动指标在5月环比年化增长4.6\%，高于4月的2.9\%。这一改善的主要驱动力来自制造业部门。报告中的制造业增长代理指标在5月出现回升，这与官方PMI数据所显示的产出扩张趋势是一致的。

然而，我们必须追问一个“所以呢”的问题：制造业的韧性能够持续多久，并且能否有效传导至消费和服务业？目前来看，答案并不乐观。制造业的改善更多是基于供给端的修复和部分外需的支撑，而非内需的全面启动。如果消费和房地产等终端需求持续疲弱，制造业的反弹很可能只是库存周期的一次短暂补库，缺乏持续扩张的动力。这就像一场接力赛，第一棒跑得不错，但第二棒迟迟未能接棒。

KC评论： GS的制造业增长代理指标（基于机床、汽车、发电设备等产出）是值得持续跟踪的先行指标。但读者必须意识到，这一指标的改善与终端需求的关联度正在减弱。完整报告中关于消费和“其他”（包含房地产和劳动力市场）领域的CAI分项数据，才是判断这一反弹能否延续的关键。这部分数据在公开摘要中被一笔带过，但在完整报告中有更细致的拆解。

报告中最值得警惕的信号，来自其自营的进口隐含国内需求代理指标。该指标显示，二季度中国国内需求增长明显走弱。这与我们通过其他渠道感知到的消费降级、房地产财富效应消退等现象高度吻合。

GS通过将中国进口商品按其最终用途，并利用投入产出表进行归因，构建了这一指标。其核心逻辑是：进口是内需的“映射”。当企业进口减少，通常意味着它们对未来的终端销售预期变得谨慎。这一指标的走弱，直接回答了“为什么制造业反弹可能不可持续”的问题——因为下游没有足够的需求来消化上游的产出。

KC评论： 许多市场参与者更关注出口数据，因为外需表现尚可。但GS的这个进口隐含内需指标，实际上是在提醒我们：不要被月度出口数据的波动所迷惑。内需的疲软是结构性的，它决定了经济的长期均衡水平。完整报告中对这一指标的方法论和与国民账户内需数据的交叉验证，提供了更严谨的分析框架，值得深入阅读。

GS的中国金融条件指数在5月小幅收紧。更关键的是，其自营的信用脉冲指标预计将在下半年继续为负。信用脉冲衡量的是新增信贷占GDP比重的变化，其转为负值，意味着信贷对经济增长的边际贡献正在下降。

这一信号的重要性在于，它直接关联到资产定价。历史上，信用脉冲的走势对中国股市和房地产市场都有较强的领先意义。当信用脉冲为负时，意味着金融体系向实体经济输送的“血液”在减少，企业的融资环境在恶化，这将对投资和经营活动构成制约。5月份金融条件的收紧，部分由人民币对贸易加权篮子升值以及股市表现疲弱导致，这进一步压缩了宽松政策的空间。

报告尚未完全回答的一个关键问题是：这种信用紧缩是暂时的季节性波动，还是政策主动调控的结果？如果是后者，那么下半年经济面临的下行压力可能会超出市场预期。

GS的报告显示，其自营的中国宏观政策代理指数在5月小幅收紧，主要驱动力来自更紧的财政政策和更慢的信贷增长。同时，其广义财政赤字指标在5月进一步收窄。

这似乎与市场对“积极财政政策”的普遍预期相悖。GS的分

[... middle omitted ...]

目、形成实物工作量的效率仍有待提高。这才是决定下半年经济动能的关键变量。

报告中关于库存的追踪器提供了一个有趣的视角。GS的库存追踪器显示，二季度库存水平可能大致持平，而库存变化对GDP增长的贡献可能下降。

这看似是一个中性甚至偏正面的信号，因为它意味着企业没有在被动累库。但我们需要思考其“二阶效应”：库存水平的稳定，可能恰恰是企业主动去库的结果。当企业对未来需求预期悲观时，它们会主动压缩库存。这种主动去库行为，本身就会压制对上游原材料和中间品的采购，从而拖累工业生产和投资。因此，库存持平背后，可能隐藏着企业信心的低迷。报告中的库存追踪器是基于大宗商品、PMI分项、产成品库存等多个指标合成的，其细微变化值得持续跟踪。

KC评论： GS的库存分析框架非常精细。它告诉我们，不能只看库存的绝对水平，更要看库存变化的驱动因素。完整报告中对库存追踪器六个底层指标的拆解，可以帮助我们判断当前的库存行为究竟是主动还是被动。这对于判断大宗商品价格和工业品利润的短期走势至关重要。

\subsection{[R020] JPM：财政不是没子弹，是没扣扳机}
{\small\color{refgray}归类：宏观与利率：美联储“沉默”重塑全球定价，日本央行内部分歧加大}

5月财政数据公布后，市场对“政策刺激力度不足”的担忧再次升温。但JPM这份最新研报给出了一个与市场直觉相反的判断：问题不在于财政空间不够，而在于部署节奏慢于预期。5月财政存款升至近十年同期第二高位，恰恰说明子弹还在枪膛里。

这不是一个关于“有没有钱”的故事，而是一个关于“什么时候、以什么方式花钱”的故事。对于所有关注中国资产定价的投资者来说，理解这个时间差，比争论总量数字更有意义。

5月一般公共预算收入同比增长6.6\%，延续了4月以来的强劲势头。证券交易印花税同比暴增145.9\%，直观反映了市场活跃度。税收收入增长6.8\%，增值税贡献最大。

但支出端却连续第二个月收缩，同比下降1.6\%。这与一季度2.6\%的增速形成鲜明对比，更远低于全年4.4\%的预算目标增速。

*收入好、支出慢**，这个组合本身就是一个信号：财政当局没有选择在二季度加速发力。原因可能不是缺乏意愿，而是对一季度GDP超预期增长后的“政策耐心”。

KC评论： 很多投资者看到财政支出收缩就认为“政策转向紧缩”，但JPM的数据表明，问题出在基建支出这个具体环节，而非整体财政收缩。5月基建类支出同比下降12\%，直接拖累了基础设施固定资产投资同比下降9.4\%。这意味着，只要基建支出恢复，整体财政数据就会有明显改善。完整报告中关于支出分项的细拆图表，值得仔细对照阅读。

如果说一般公共预算还有韧性，政府性基金账本则面临更大的压力。5月政府性基金收入同比下降20.3\%，其中土地出让收入加速下滑至-35.8\%，比2025年全年-14.7\%的降幅进一步恶化。

土地收入的持续萎缩，直接拖累政府性基金支出同比下降11.5\%。由于政府性基金是地方基建投资的重要资金来源，这一链条的收缩对整体投资形成直接压制。

JPM在报告中特别指出，政府性基金目前已成为政府基金账目中最大的融资来源，但其高度依赖土地收入的模式，在当前房地产下行周期中面临结构性挑战。

KC评论： 市场习惯把“土地收入下降”和“地方财政困难”划等号，但JPM的分析提醒我们：关键变量不是土地收入本身，而是地方政府能否找到替代性融资渠道。报告提到的专项债发行节奏和政策性金融工具，正是观察这个替代路径的核心指标。这些图表在完整报告中都有详细呈现。

这是整份报告最值得关注的信号之一。5月财政存款进一步上升，达到近十年同期第二高水平，仅次于2025年。

财政存款高企，意味着已批准的预算资金尚未被充分部署到实际项目中。JPM认为，这背后至少有三个原因：一是合格项目储备不足，二是地方政府优先用于偿债而非新投资，三是政策性银行工具作为项目启动资金的落地速度偏慢。

*这不是一个“缺钱”的故事，而是一个“钱还没花出去”的故事。**

对于下半年的经济走势而言，这恰恰提供了一个缓冲垫：财政存款的释放空间，可以在需要时转化为实际的支出和投资。

好消息是，进入6月后，地方政府专项债发行节奏明显加快。截至6月第三周，月内发行规模已接近5000亿元，较4-5月月均1680亿元的水平大幅提升。

JPM将此定性为“反应式加速”：即面对二季度经济放缓，政策层面开始加快利用3月全国两会已经批准的财政资源。

但这里有一个关键的时间差：从债券发行到资金实际拨付到项目，再到形成实物工作量，通常需要数周甚至数月。因此，6月发行的加速，大概率不会在当月投资数据中完全体现。

KC评论： 这是投资者最容易误判的地方。看到专项债发行加速就认为“经济马上会有反应”，往往会在短期数据验证时失望。JPM给出的框架是：先看发行，再看拨款，最后看投资。这个传导链条的每个环节都有滞后。完整报告中关于债券发行节奏与投资数据的对比图表，能帮助读者更直观地理解这个时间差。

第一阶段是“执行期”：加快部署全国两会已批准的预算，加速项目审批，推动政府债

[... middle omitted ...]

标框架：

第一，财政存款的月度变化。如果财政存款开始明显下降，意味着资金正在从国库流向实体经济，这是最直接的“钱花出去了”的信号。

第二，专项债资金拨付与项目开工的匹配度。仅看发行量是不够的，需要跟踪资金是否真正进入了施工环节。基础设施投资数据将是最终的验证。

第三，土地收入降幅是否触底。土地收入是地方政府财力的核心变量，如果降幅不再扩大，地方财政的约束将边际放松，从而为基建投资提供更多空间。

这三个指标结合起来，比单纯盯着月度GDP数据，能更早地判断财政实际发力的拐点。

*顺着这些未解问题继续往下读：** JPM这份报告里还有大量关于支出结构、债券发行节奏、土地收入趋势的详细图表和细分数据，对于判断上述三个先行指标的拐点至关重要。我们每天会由AI agent结合人工review，生成国际投行中文摘要与KC评论，daily更新约10-40页，整理当天最新数据图表合集，方便喂给AI，也方便人工快速把握market dynamics。欢迎来星球微信群里继续讨论。

\subsection{[R021] NOM：人民币中间价模型释放了一个被忽视的鸽派信号}
{\small\color{refgray}归类：宏观与利率：美联储“沉默”重塑全球定价，日本央行内部分歧加大}

市场对人民币汇率的讨论，大多集中在央行意图、中美利差和出口表现上。但一份来自NOM亚洲外汇策略团队的模型报告，提供了一个更底层、更量化的观察视角：人民币每日中间价的定价机制本身，正在发出值得关注的信号。

这份报告的核心判断是：NOM的人民币中间价模型显示，当前中间价的设定存在系统性偏低，且逆周期因子已显著收窄。 这并不意味着央行在主动引导贬值，而是意味着在当前的定价框架下，市场力量正在发挥更大作用，而政策干预的边际成本正在上升。

对于关注宏观对冲、跨境资产配置和汇率风险管理的读者来说，理解这个“模型信号”比猜测“央行意图”更具实际意义。因为前者是可量化、可追踪的，后者则往往被噪音覆盖。

NOM模型在6月24日的最新投影为6.7910，较前值6.8171低了261个基点。这一数字本身并非关键，关键在于它与官方实际中间价之间的差距，以及这种差距的演变趋势。

模型投影是NOM基于其预设的定价因子（包括前日收盘价、一篮子货币变动、央行政策倾向等）计算出的“理论中间价”。当实际中间价显著低于模型投影时，通常意味着央行在主动引导汇率走弱；反之，则意味着央行在支撑汇率。

但报告提供的图2（模型误差走势图）显示，近期的模型误差正在收窄，甚至出现方向性逆转。这指向一个更值得关注的结论：央行对中间价的干预力度在减弱，市场化的定价因子权重在上升。

KC评论： 很多市场参与者过度关注“央行是否在引导贬值”这个模糊问题。NOM模型提供了一个更精确的观察框架：不是看央行“想”做什么，而是看“实际”中间价与“模型”中间价的偏差。偏差缩小，意味着市场定价权在回归。完整报告中包含了模型误差的历史走势图，能清晰看到这种结构性变化是从何时开始的，以及背后的驱动因子是什么。

报告特别指出，当加入逆周期因子后，模型投影变为6.8025，仅比前值低146个基点。这意味着，逆周期因子本身起到了显著的“缓冲”作用，但其调整幅度本身也在收窄。

逆周期因子是央行用来对冲市场顺周期行为、稳定汇率预期的重要工具。当该因子数值较大时，说明央行在主动干预；当它收窄时，说明央行对市场波动的容忍度在提高。

这份报告传递的信号是：央行正在从“强力干预”模式转向“有限干预、更多依靠市场”的模式。 这背后的逻辑可能包括：外部压力（如中美关系）阶段性缓和、内部经济复苏需要更灵活的汇率环境，以及央行对人民币汇率弹性的信心增强。

对于投资者而言，这意味着未来人民币汇率的波动性可能加大，但方向性趋势将更多由经济基本面而非政策指令决定。

3. 报告没有完全回答的关键问题：这种“干预退坡”是主动选择还是被动接受？

这是整份报告最值得追问的地方，也是NOM模型无法直接回答的问题。

一种可能是，央行主动选择减少干预，以换取更灵活的货币政策空间，或者为即将到来的政治议程（如7月政治局会议、APEC峰会）营造更平稳的宏观环境。另一种可能是，央行正在被动接受市场力量，因为干预的成本（如外汇储备消耗、利率传导效率损失）已经超过了收益。

模型本身只能描述“发生了什么”，无法解释“为什么发生”。但这个问题的答案，直接决定了未来人民币汇率的走势路径。如果是前者，汇率可能在可控范围内双向波动；如果是后者，一旦市场形成单边预期，央行可能被迫重新加大干预。

KC评论： 这是所有量化模型的天花板——它告诉你“是什么”，但无法告诉你“为什么”。要回答这个问题，需要结合中国的内部经济数据（如信贷脉冲、通胀压力）和外部地缘政治环境（如中美会谈进展）。报告中的日历（图4）列出了7月政治局会议、APEC峰会等关键事件，这些时点可能是判断央行意图的重要窗口。完整报告对这些事件与汇率政策联动性的分析，值得仔细研读。

4. 对资产定价的隐含含义：汇率弹性上升，可能改变所有人民币资产的定价逻辑

如果N

[... middle omitted ...]

 一个值得持续跟踪的观察框架：如何利用模型误差判断市场情绪

如果模型误差持续为正且扩大，说明央行在主动引导汇率走弱，市场情绪可能偏悲观。 如果模型误差持续为负且扩大，说明央行在支撑汇率，市场情绪可能过度悲观。 如果模型误差收窄或转正，说明央行干预减弱，市场定价权回归。

当前，模型误差的收窄暗示了央行对市场力量的认可。但这一趋势能否持续，取决于后续的经济数据和外部事件。投资者可以将这个框架纳入自己的宏观监测体系，每周或每月追踪一次，比单纯猜测“央行意图”更具操作性。

6. 报告尚未展开的更深层问题：人民币汇率定价机制改革的长期方向

NOM的这份模型报告聚焦于短期定价，但它触及了一个更根本的问题：人民币汇率形成机制的市场化改革，到底走到了哪一步？

从2015年“8·11汇改”至今，人民币汇率定价机制经历了多次调整。每一次调整都伴随着市场波动和政策博弈。当前，中间价定价机制中“逆周期因子”的权重下降，是否意味着改革正在重启？还是说，这只是特定外部环境下的权宜之计？

\subsection{[R022] JPM：中国MDI出口量价齐升，欧洲市场正成为新的定价锚}
{\small\color{refgray}归类：能源与大宗：铝缺口被低估，黄金与2013年本质不同，航运逻辑切换 / 中国化工：传统品疲弱，氟化工与MDI的结构性机会}

中国化工品出口正在经历一场静悄悄的结构性转变。JPM最新发布的MDI（二苯基甲烷二异氰酸酯）贸易数据显示，2026年5月中国MDI出口量达到10.1万吨，同比飙升49\%，出口均价环比上涨13\%、同比上涨25\%。这不是一次简单的周期性反弹——出口流向的重新布局和定价能力的恢复，正在改写全球MDI的竞争规则。

这份由Jeffrey J. Zekauskas领衔的研报揭示了一个关键信号：在中国出口总量尚未恢复到2024年峰值水平的情况下，出口价格已经率先突破。5月MDI出口均价达到2193美元/吨，不仅是近两年来的最高水平，更意味着中国供应商正在从“以量换价”转向“以价换量”的新阶段。

对于关注化工板块的投资者而言，这份数据最重要的判断是：全球MDI的供需再平衡正在加速，而定价权正在向供给端倾斜。2025年全年中国MDI净出口仅为40万吨，较2024年大幅下降49.6万吨，但进入2026年，前五个月净出口已达到25.8万吨，同比增长5.5万吨。这个拐点意味着什么？中国MDI出口的“量价双杀”阶段已经结束，取而代之的是“量价齐升”的新周期。

2025年，中国对美国的MDI出口同比下降了83\%，这是一个几乎断崖式的下滑。但与此同时，对欧盟27国的出口却在快速回升。2026年前五个月，中国对欧盟的MDI出口量达到8.2万吨，占出口总量的21\%，较2025年同期的5万吨增长了63\%。

这一结构性变化并非偶然。2025年中国MDI出口总量仅为80.5万吨，较2024年的120.4万吨下降了33\%，其中对美国和荷兰的出口降幅最为显著。但进入2026年，出口恢复的路径清晰指向了欧洲市场。除欧盟外，对土耳其的出口也从2.8万吨增长至4.7万吨，增幅达68\%。

这意味着什么？全球MDI的贸易流向正在被地缘政治和供应链重组重新定义。美国市场的大门在收窄，但欧洲市场正在为中国供应商打开一扇更大的窗。

5月MDI出口均价2193美元/吨，环比上涨13\%，同比上涨25\%。更值得关注的是进口价格的变化——5月中国MDI进口均价达到2979美元/吨，环比暴涨56\%，同比飙升115\%。

进出口价差的扩大是一个极其重要的信号。通常，中国作为全球最大的MDI生产国，出口价格低于进口价格是常态，但5月两者之间的价差扩大到786美元/吨，这是近年来罕见的水平。这说明全球高端MDI市场仍然供不应求，而中国供应商正在利用这一价差获取更高的出口利润。

回顾2024年，中国MDI出口均价从年初的1277美元/吨一路攀升至8月的2020美元/吨，随后有所回落。进入2025年，均价一度跌至1416美元/吨的低点，但自2026年3月以来连续三个月大幅反弹。这种V型反转的背后，是全球MDI产能投放节奏放缓和下游需求复苏的双重驱动。

3. 2025年的出口收缩并非需求问题，而是供给端的主动调整

2025年中国MDI出口量同比下降33\%，净出口从2024年的89.6万吨骤降至40万吨。乍看之下，这似乎是需求疲软的信号。但结合2026年以来的数据，JPM的分析指向了另一种解释：2025年的收缩更多是供给端主动调整的结果。

从历史数据看，2025年的出口量（80.5万吨）甚至低于2021年（101.4万吨）和2022年（98.8万吨）的水平，但2026年前五个月的出口量（39.6万吨）已经同比恢复15.6\%。如果这一趋势持续，2026年全年出口量有望回到100万吨以上。

值得注意的是，2025年进口量反而增长了31.4\%，达到40.5万吨。这意味着国内需求并未大幅萎缩，而是国内供应商在2025年主动减少了出口，以应对价格低迷和贸易摩擦。进入2026年，随着价格回升，出口动力重新被激活。

4. 报告尚未完全回答的关键问题：欧洲市场的承接能力能持续多久？



[... middle omitted ...]

均价2193美元/吨，虽然较2024年低点大幅回升，但仍远低于2021年9月2343美元/吨的历史高点。如果下游需求持续恢复，价格仍有上行空间；但如果欧洲经济出现超预期下行，价格可能再次承压。

第一，关注中国MDI出口均价与进口均价的价差。如果价差持续扩大，意味着中国供应商在全球定价体系中的地位在提升，这对国内龙头企业是明确的利好。

第二，关注对欧洲出口占比的变化。21\%的占比是否能够维持甚至提升，将直接决定中国MDI出口的长期增长空间。如果欧洲市场占比继续上升，说明中国供应商正在成功实现市场多元化。

第三，关注月度出口量的环比趋势。5月10.1万吨的出口量是2026年以来的最高水平，但能否在6月、7月保持这一量级，才是判断需求真实强度的关键。季节性因素和下游补库节奏都会影响月度数据，需要连续观察。

上述三个指标中，任何一个出现超预期变化，都可能改变当前的定价逻辑。JPM这份报告提供了详实的基础数据，但真正有价值的判断往往来自于对这些数据的持续追踪和交叉验证。

\subsection{[R023] JPM：中国锂贸易已发生结构性逆转，净进口时代开启}
{\small\color{refgray}归类：中国化工：传统品疲弱，氟化工与MDI的结构性机会}

中国锂产业正在经历一场被多数市场参与者低估的结构性转变。JPM最新发布的2026年5月中国锂进出口数据显示，中国已经从氢氧化锂的净出口国彻底转为净进口国，同时碳酸锂的进口依赖度正在以超出预期的速度加深。这不仅仅是月度波动，而是一个可能重塑全球锂定价体系与供应链格局的拐点。

2026年前五个月，中国氢氧化锂净进口量达到10.5千吨，而2025年同期还是净出口35.4千吨。碳酸锂净进口量同比增长54.5\%，达到151.4千吨。更值得关注的是，5月碳酸锂进口均价同比上涨102\%，氢氧化锂出口均价同比上涨64\%。价格与量的同步上升，指向的不是短期补库，而是供需关系的底层逻辑在变化。

这份报告的核心判断是：中国锂产业链的“两头在外”格局（上游资源依赖进口、下游产品出口全球）正在被打破，但方向与市场预期相反——中国正在从锂化学品出口国转变为净进口国。这将对全球锂矿商、加工企业以及下游电池制造商的定价权与竞争格局产生深远影响。

1. 氢氧化锂的贸易地位已发生不可逆转变，中国不再是全球最大供应方

JPM的数据清晰地展示了这一转变的轨迹。从2017年到2023年，中国一直是氢氧化锂的净出口国，净出口量从18.1千吨增长到126.2千吨，年复合增长率超过38\%。但2024年净出口量骤降至35.4千吨，2026年前五个月更是直接转为净进口10.5千吨。

月度数据更能说明变化的节奏。2026年1月至5月，中国氢氧化锂出口量分别为1.9千吨、2.6千吨、3.1千吨、5.5千吨和3.5千吨，远低于2023年同期的9.6千吨、10.3千吨、10.2千吨、10.1千吨和11.0千吨。出口量下降了约60\%-70\%。

与此同时，进口量在快速攀升。2026年前五个月，氢氧化锂进口量分别为6.8千吨、3.7千吨、6.1千吨、6.7千吨和3.9千吨，而2025年同期仅为1.0千吨、1.4千吨、2.0千吨、1.3千吨和0.8千吨。进口增长超过5倍。

2. 碳酸锂进口量价齐升，中国正从“加工枢纽”变为“消化黑洞”

碳酸锂的情况同样值得警惕。2026年5月，中国碳酸锂进口量达到37.6千吨，同比增长78\%，环比增长15\%。进口均价达到18,993美元/吨，同比上涨102\%，环比上涨15\%。

将时间拉长，趋势更加明显。2026年前五个月，碳酸锂净进口量为151.4千吨，同比增长54.5\%。而2025年全年净进口量也仅为237.7千吨。按照当前节奏，2026年全年净进口量可能突破360千吨，创下历史新高。

进口价格的上涨尤其值得关注。2025年5月，碳酸锂进口均价仅为9,392美元/吨，到2026年5月已经翻倍。这不仅仅是基数效应，而是反映了全球锂资源成本曲线的上移——澳洲锂辉石、南美盐湖的成本都在上升，而中国本土锂资源开采成本更高、环保约束更严。

JPM的数据揭示了一个有趣的现象：中国氢氧化锂出口价格与进口价格出现了背离。5月氢氧化锂出口均价为19,856美元/吨，环比上涨33\%，同比上涨64\%。但进口均价仅为18,000美元/吨，低于出口价格。

这种“出口价比进口价贵”的倒挂，在传统贸易模型中是不合理的。通常，出口价格应该低于进口价格，因为出口产品包含了加工溢价。但当前的情况表明，中国出口的氢氧化锂可能更多地流向高溢价市场（如日韩高端电池客户），而进口的氢氧化锂可能来自成本更低的海外产能。

碳酸锂的情况则相反。5月碳酸锂出口均价为18,993美元/吨，而进口均价为19,856美元/吨，进口价格高于出口价格。这符合中国作为碳酸锂净进口国的定位——进口的碳酸锂需要支付更高的边际成本。

4. 报告尚未完全回答的关键问题：中国本土锂资源开发为何没有跟上

JPM的报告提供了详实的数据，但有一个关键问题没有完全展开：既然进口依赖度在快速上升

[... middle omitted ...]

开始返销？如果返销量在增加，那么进口量的增长更多反映的是中国企业“资源全球化配置”的结果，而非单纯的供应短缺。这会影响对锂价的判断。

5. 投资者的观察框架：从“中国出口”转向“中国需求”的定价逻辑重构

基于JPM的数据，投资者需要重新建立对锂行业的三层观察框架：

第一层：中国贸易地位的转变是否可持续。氢氧化锂的净进口转变已经持续了五个月，碳酸锂的进口增速也在加速。如果这一趋势在2026年下半年延续，那么中国锂化学品价格将更多地由海外边际成本决定，而非国内供需。

第二层：下游需求结构的变化。高镍正极材料（NCM811等）对氢氧化锂的需求仍在增长，而磷酸铁锂对碳酸锂的需求也在扩大。如果中国电池企业的技术路线继续向高镍方向倾斜，氢氧化锂的进口依赖度可能会进一步加深。

第三层：全球锂资源供给的再平衡。中国从净出口国转为净进口国，意味着全球锂资源需要重新分配。澳洲和南美的锂矿商将获得更大的议价权，而中国的加工企业将面临成本上升的压力。这可能会推动新一轮的上下游整合。

\subsection{[R024] JPM：中国钛白粉出口的“量价齐升”并非故事的全部}
{\small\color{refgray}归类：中国化工：传统品疲弱，氟化工与MDI的结构性机会}

中国钛白粉出口正在经历一个被市场低估的结构性转变。JPM最新发布的2026年5月出口数据显示，中国钛白粉出口量同比增长13\%，出口均价环比上涨8.8\%、同比上涨7.0\%，达到每吨2218美元。表面上看，这是一组“量价齐升”的乐观信号。但深入拆解这份报告中的细分数据后，一个更复杂的图景浮现出来：真正驱动这一轮变化的，并非总量逻辑，而是产品结构、区域流向和定价权的深层分化。对于关注化工行业、尤其是涉及全球供应链重组的投资者而言，忽略这些细节，可能会错判中国钛白粉企业的真实竞争力和可持续性。

*这份JPM报告的真正价值，不在于告诉你“出口在增长”，而在于揭示了增长的质量正在发生根本性变化。**

1. 氯化法钛白粉正在成为出口增长的核心引擎，而非传统的硫酸法

JPM报告中最具区分度的数据点，是氯化法与硫酸法钛白粉出口的增速差异。2026年前五个月，中国氯化法钛白粉出口量同比增长39\%，而硫酸法仅增长6\%。这是一个数量级的差距。

氯化法钛白粉代表更高的技术壁垒、更稳定的品质和更高的附加值。长期以来，中国钛白粉出口以硫酸法为主，氯化法占比有限。但2026年的数据表明，中国企业在氯化法领域的产能释放和技术突破，正在加速改变出口结构。5月单月，氯化法出口量达34.4千吨，而2025年同期仅为24.9千吨。

KC评论： 氯化法出口增速是硫酸法的6.5倍，这不仅仅是产品组合的变化。它意味着中国钛白粉企业正在从“低成本、大路货”的定位，向中高端市场渗透。对于投资者而言，关注各家企业氯化法产能的释放节奏，远比关注总出口量更有意义。完整报告中的图表6提供了逐月比较，能帮助判断这一趋势的持续性。

这一结构性变化对全球钛白粉竞争格局的含义是深远的。传统上，氯化法钛白粉由科慕、特诺等国际巨头主导。中国企业的加速渗透，意味着全球高端市场正在出现新的供给压力。而硫酸法出口增速的放缓，则可能反映国内环保约束和产能出清正在发挥作用。

2. 出口目的地正在从欧盟向印度和东南亚转移，区域集中度下降

JPM报告展示了2025年与2026年出口目的地的显著变化。2026年前五个月，中国对欧盟27国的出口量为97千吨，同比增长29\%，但增速已低于印度。印度同期进口量达166千吨，同比增长32\%，取代欧盟成为中国钛白粉最大的出口目的地。

更值得关注的是，2025年全年中国对欧盟出口占比已降至9\%，而2024年这一比例为13\%。欧盟反倾销调查的影响正在持续显现。与此同时，土耳其、印度尼西亚、越南等市场的份额在提升。出口目的地从“高度依赖欧盟”向“多元化新兴市场”转变，这既分散了地缘政治风险，也意味着中国钛白粉企业需要适应不同市场的定价和渠道策略。

KC评论： 区域流向的变化直接关系到定价权。欧盟市场对品质和认证要求更高，但价格敏感度相对较低；印度和东南亚市场则更注重成本竞争力。JPM报告中的表4和表5给出了详细的国别数据，能帮助判断哪些企业在不同市场具有竞争优势。完整报告中的这些表格值得仔细研读。

这一趋势也暗示，中国钛白粉出口的“天花板”可能正在从地缘政治约束转向新兴市场需求容量。如果印度、东南亚的基础设施和房地产投资周期持续，中国钛白粉的出口增长仍有空间。但反过来，如果这些区域经济放缓，集中度下降反而可能放大出口波动。

5月出口均价2218美元/吨，环比上涨8.8\%，同比上涨7.0\%。这一数字看似亮眼，但JPM报告中的细分数据显示，价格回升的驱动力高度集中。

中国对欧盟27国的出口均价从1月的1945美元/吨持续攀升至5月的2431美元/吨，5个月累计涨幅达25\%。而同期对印度、东南亚等市场的均价涨幅远低于此。这意味着整体出口均价的回升，很大程度上是由欧盟市场的高价订单拉动的，而非全球性提价。

与此同时，中国钛白粉进口均价5月为29

[... middle omitted ...]

域结构，过度依赖欧盟市场的企业，其盈利稳定性可能低于市场预期。

2026年前五个月，中国钛白粉净出口量达857.2千吨，同比增长14.3\%，较2025年同期增加107千吨。这一数字创下历史同期新高。

但值得注意的是，5月单月出口量152.8千吨，虽然同比增长13\%，但环比4月的193.5千吨下降了21\%。这种“同比增、环比降”的模式，在硫酸法产品上表现得更为明显——硫酸法出口量5月为118.4千吨，环比4月的150.7千吨下降21.4\%。

环比下降可能反映季节性因素，也可能暗示海外市场在经历了一季度的集中补库后，采购节奏正在放缓。如果后续几个月出口量无法恢复环比增长，库存压力可能会在国内市场累积，进而对国内价格形成压制。

KC评论： 净出口量创新高是事实，但环比下降也是事实。关键在于判断这是短期波动还是趋势拐点。JPM报告中的表1提供了长达12年的月度出口数据，可以帮助读者判断5月环比下降是否处于历史正常波动区间。完整报告中的这些历史数据是判断趋势的宝贵工具。

\subsection{[R025] BARC：AI硬件股的估值游戏正在从光学转向服务器}
{\small\color{refgray}归类：AI/算力：从“光学狂欢”到“网络务实”，半导体拥挤度达历史极值}

当市场还在为AI算力需求是否见顶争论不休时，一份来自BARC的最新双周追踪报告揭示了一个更微妙的信号：AI硬件股的隐含估值正在经历结构性分化，而光学板块的“休息”可能只是表象。

这份报告的核心判断是：AI硬件股的隐含PE倍数已经从两个月前普遍的20-40倍区间，跃升至当前的40-60倍区间，但推动这一变化的驱动力正在从光学转向服务器与网络设备。市场对AI的定价逻辑，正在从“谁有AI概念”转向“谁的AI收入能持续兑现”。

这不是一个简单的板块轮动。它意味着，投资者需要重新审视手中AI硬件股的估值锚点——那些被市场赋予高AI溢价的股票，其支撑逻辑是否足够坚实？

1. 光学板块的估值回调并非基本面恶化，而是市场预期被过度透支

过去两周，BARC追踪的AI概念股平均下跌了中个位数，而光学板块是拖累主力。CIEN下跌25\%，FN下跌11\%，GLW也出现回调。直观判断是光学需求转弱，但BARC的分析指出，这更多是技术性回调而非基本面逆转。

CIEN在财报前股价已累计上涨107\%，远超同期标普500指数11\%和纳斯达克22\%的涨幅。财报本身并不差，但市场设定的预期过高。当实际数字无法匹配股价已经反映的“完美预期”时，回调便不可避免。FN的下行则更多是受CIEN的连带影响。

KC评论： 这提醒我们一个关键问题：AI硬件股的估值已经高度依赖“预期中的预期”。当股价提前反映了未来两到三个季度的利好，任何微小的miss都会引发剧烈修正。完整报告中的隐含AI PE时间序列图表，清晰地展示了这种预期透支的轨迹——从2月到5月，CIEN的隐含AI PE从58倍飙升至115倍，这个数字本身就意味着市场几乎没有任何容错空间。

对于关注光学板块的投资者而言，当前的回调是否提供了重新入场的机会，取决于能否区分“需求拐点”与“预期修正”。BARC并未给出明确结论，但暗示了关键变量：CIEN的订单积压和客户部署节奏。这些细节在完整报告中有更深入的拆解。

2. 服务器厂商正在成为AI估值的新引擎，但HPE的案例值得警惕

与光学板块形成鲜明对比的是，HPE在同期上涨了12\%，成为AI概念股中最大的赢家。其隐含AI PE倍数从系列追踪初期的个位数，迅速攀升至当前的58倍。驱动因素包括：强劲的财报、新的激进投资者入场，以及DELL财报带来的正面溢出效应。

HPE的案例揭示了当前市场的另一个特征：传统服务器厂商正在被重新定价。DELL的隐含AI PE维持在58倍，HPE也从底部快速追赶。BARC指出，这些公司的传统业务（on-premise）表现超预期，加上AI订单的持续增长，使得市场愿意给予更高的AI估值溢价。

但这里存在一个结构性矛盾：HPE的AI收入占比仅为16.6\%，远低于SMCI的86.9\%或FN的82.5\%。然而其隐含AI PE却与DELL持平，甚至高于部分AI收入占比更高的公司。这意味着，市场对HPE的AI估值溢价，很大程度上来自其传统业务“不拖后腿”+AI业务“有想象空间”的组合，而非AI业务本身的规模和增速。

KC评论： 这种估值逻辑的脆弱性在于，一旦传统业务出现波动或AI业务增速不及预期，估值修正的空间可能比光学板块更大。完整报告中有一张核心图表，将每家公司的AI收入占比、增长率与隐含AI PE做了并列对比。这张图的价值在于，它让投资者可以直观地看到哪些股票的估值有“水分”——即AI收入占比低但隐含AI PE高的公司，其估值对AI故事的依赖程度反而更高。

3. SMCI的融资事件揭示了一个被忽视的风险：AI订单的兑现能力比订单规模更重要

SMCI是本轮调整中跌幅最大的个股，两周内下跌40\%。直接触发因素是公司宣布了70亿美元的股权及股权挂钩融资计划，以支持AI订单交付。市场对此的反应是抛售，对稀释预期的担忧压倒了

[... middle omitted ...]

经不再是稀缺信息，市场开始关注“质”——即订单能否按时、按预期利润兑现。

KC评论： 这可能是这份报告中最值得深思的一个点。当整个AI产业链都在追逐“订单增长”时，SMCI的案例提醒我们：订单规模与收入确认之间存在巨大的不确定性。完整报告中对SMCI的详细分析，包括其历史交付记录和客户集中度，为判断这种不确定性提供了更完整的框架。对于持有或关注AI硬件股的投资者，理解“订单兑现率”这个指标，可能比追逐订单增速更有意义。

BARC在本期报告中首次系统性地分析了其覆盖范围内“非AI概念股”的估值情况。这些公司包括AAPL、GRMN、AXON、MSI等，其共同特征是AI收入占比极低或尚未被市场纳入AI叙事。

数据显示，这些股票自5月29日以来平均下跌了高个位数，而同期标普500仅下跌2\%。它们的平均CY27 PE为22倍，远低于AI概念股的平均27倍。更重要的是，其“核心PE”平均为31倍，高于AI概念股的16倍——这意味着非AI概念股的估值中，几乎没有AI溢价成分。

\subsection{[R026] Citi：端午出游数据背后，中国消费韧性比想象中更“挑”}
{\small\color{refgray}归类：中国消费与地产：K型分化加速，定价权成为核心考验}

当宏观数据与微观体感出现背离时，市场往往倾向于相信悲观的一面。但2026年端午假期的中国旅游数据，却讲述了一个更精细的故事。

根据中国文化和旅游部最新数据，端午三天假期（6月19-21日），国内旅游人次同比增长4.4\%，旅游收入同比增长4.0\%。在同期全国交通客运量同比下降0.9\%的背景下，这份成绩单显得尤为突出。

Citi团队在最新发布的报告中明确指出，这一数据“超出了我们的预期”。他们的核心判断是：中国休闲旅游需求依然坚韧，而出行结构正在发生深刻分化——非旅游目的的出行在减少，但休闲度假的意愿并未消退。

对于关注中国消费市场的决策者而言，这份报告的真正价值不在于确认“旅游复苏”，而在于揭示了一个更关键的变化：消费的韧性正在从“总量驱动”转向“场景驱动”，而能够识别并服务于特定场景的企业，正在获得定价权。

端午假期，全国交通客运总量同比下降0.9\%，其中占主导地位的自驾出行下降2.1\%。这一数字很容易被解读为“消费意愿不足”。但Citi分析指出，这种背离恰恰说明了问题的另一面。

雨天减少了非旅游目的的出行需求——通勤、探亲、短途办事等被压缩，但休闲旅游需求保持了韧性。换句话说，人们不是不出门了，而是“不出无意义的门”。当出行被赋予明确的“休闲度假”目的时，消费者的决策并未动摇。

这意味着什么？对于旅游、酒店、免税等行业的观察者而言，判断消费趋势的指标需要从“人流量”转向“人均消费意愿”和“有效出行目的”。宏观交通数据可能掩盖了结构性的消费韧性。

KC评论： Citi这个观察非常关键。市场习惯用“出行人数”来推断消费景气度，但这份报告提醒我们，出行结构比出行总量更能揭示真实需求。当你看到交通数据走弱时，不妨追问一句：是所有人都不走了，还是只有“非必要出行”在减少？完整报告中，Citi对比了不同出行方式的增速差异，并给出了更精细的拆分逻辑——这份图表对于理解中国消费的真实底色很有价值。

2. 酒店业正在经历“区域分化”，头部连锁的ADR定价权成为关键胜负手

Citi研报中一个值得关注的细节是，不同区域的酒店RevPAR表现出现了显著分化。以上海为例，端午期间酒店入住率同比增长了1.6个百分点。

这一分化背后，是酒店行业竞争逻辑的转变。在行业供给经历多年调整后，市场已从“数量扩张”进入“品质竞争”阶段。对于亚朵这类聚焦中高端市场、拥有强品牌溢价和管理能力的连锁酒店，其优势正在被重新定价。

Citi在报告中特别指出，亚朵在端午假期实现了正向RevPAR增长，入住率和平均房价均有改善。这正是Citi将其列为行业首选买入标的的核心逻辑——在RevPAR波动的环境下，亚朵的盈利韧性被市场低估，而其零售业务的持续增长又构成了第二增长曲线。

KC评论： 酒店业的区域分化，本质上是对“品牌溢价能力”的一次压力测试。当整体客流承压时，头部品牌反而能通过定价权维持甚至提升业绩。Citi报告中对亚朵的估值框架（14倍2026年EV/EBITDA，较行业平均溢价10\%）值得细读——它不只是数字，更是对“品牌护城河”的一种量化表达。完整报告中的估值模型和风险分析，可以帮助你理解这个溢价是否合理。

3. 出境游恢复正增长，但真正的增量来自“政策红利”驱动的入境游

端午期间，中国出入境日均流量同比增长12.9\%。其中中国大陆居民出境增长5.5\%，港澳台居民增长18.4\%，而外国人入境增长高达23.3\%。更值得注意的是，免签政策下的外国人入境人数同比增长15.2\%。

这组数据揭示了一个被市场低估的趋势：中国入境旅游市场正在经历一轮结构性增长，其动力并非来自自然恢复，而是来自签证便利化等政策红利的持续释放。

对于免税、酒店、高端零售等行业而言，这是一个重要的增量来源。海南离岛免税在端午期间销售额同比增长8.6\%，主

[... middle omitted ...]

6年春节和元旦期间人均消费大幅增长形成了鲜明对比。

Citi在报告中客观呈现了这一数据，但并未给出明确归因。这可能是天气因素导致的短期波动，也可能是消费结构变化的早期信号——当更多中低收入群体加入旅游大军时，人均消费自然会承压。

这里存在一个关键假设：如果人均消费持续承压，那么旅游行业的增长将更多依赖“量”而非“价”。这意味着，那些能够通过规模效应和成本控制维持利润率的企业将受益，而依赖高客单价和品牌溢价的企业可能面临压力。

报告没有给出这个问题的答案，但提出了一个值得持续跟踪的框架。对于投资者而言，未来几个假期的数据（特别是即将到来的国庆黄金周）将是验证这一趋势的关键窗口。

5. 对产业决策者的观察框架：从“总量”到“结构”的思维转换

综合Citi这份报告，我们可以提炼出一个观察中国消费市场的分析框架，它包含三个层次：

第一，区分“出行”与“消费”。交通客运量下降不代表消费意愿下降，关键在于区分“非必要出行”和“休闲度假出行”。后者才是消费的真实驱动力。

\subsection{[R027] Citi：端午节旅游数据背后的真正信号——OTA行业正在经历一场“消费降级式繁荣”}
{\small\color{refgray}归类：中国消费与地产：K型分化加速，定价权成为核心考验}

Citi：端午节旅游数据背后的真正信号——OTA行业正在经历一场“消费降级式繁荣”

端午节三天假期，中国国内旅游人次1.24亿，旅游收入445亿元，同比分别增长4.4\%和4.0\%。如果只看这两个数字，你可能会得出“消费复苏平稳”的结论。但Citi这份研报真正值得关注的判断，藏在一个不起眼的细节里：人均消费同比下降0.4\%。

这不是一个季度的问题。从2026年元旦到端午节，人均旅游消费同比增速分别为-1.0\%、-0.2\%、-0.2\%、-0.7\%、-0.4\%。连续五个假期，人均消费持续收缩。与此同时，出游人次却保持正增长。

这意味着什么？Citi的结论是：国内旅游市场整体表现“有韧性”，但韧性来自量的扩张，而非价的提升。对于OTA平台而言，这轮增长的红利，正在从“卖更贵的机票和酒店”转向“卖更多的低价短途产品”。

这并不是一个坏消息。Citi仍然维持对携程和同程旅行的买入评级。但这份报告真正有价值的，不是结论本身，而是它揭示的一个结构性变化：当人均消费持续下行，OTA平台的竞争逻辑正在被重新定义。

1. 人均消费连续五个假期下滑，这不再是天气或油价可以解释的短期波动

Citi报告提供了一个非常重要的数据序列。从2026年元旦到端午节，国内旅游人均消费分别为597元、1348元、455元、571元、359元，同比增速全部为负值。

Citi在报告中提到，端午节期间全国客运量同比下降0.9\%，铁路和航空分别增长2.9\%和0.6\%，但公路和水路分别下降1.2\%和1.8\%。研报将其归因于假期期间的暴雨天气和航空票价上涨。

但如果我们把时间拉长，连续五个假期人均消费为负，就不能简单用“天气”或“油价”来解释。更合理的推断是：中国消费者的旅游预算正在结构性收紧，人们仍然愿意出行，但更倾向于选择低成本方案。

KC评论： 人均消费下滑不是端午节特有的现象，而是一个延续了半年的趋势。Citi报告里这张表格（Figure 1）值得仔细看——它把2025年所有假期与2026年做了对比。2025年端午人均消费359元，同比增速是+0.2\%；2026年端午人均消费同样是359元，但同比增速变成了-0.4\%。同样的绝对值，不同的方向，说明基数效应和价格弹性都在发生变化。完整报告里还有2025年国庆的数据，可以作为对比基准。

2. 短途“微度假”和“小城”成为主旋律，OTA平台正在争夺同一个客群

Citi报告引用了三家OTA平台在端午期间的运营数据，指向同一个方向。

携程的数据显示，热门目的地包括大城市和“隐藏宝石”小城镇，女性用户和90后是主力，以跨省短途游为主。同程旅行则强调，短途“微度假”依然流行，江浙沪地区尤为突出，小城镇和“乡村漂”旅游趋势明显。飞猪的数据略有不同，人均预订价值和酒店入住时长显著增长，但驱动因素是对高品质旅行的偏好，以及西北和东北地区避暑和草原旅游需求的激增。

这三组数据看似有差异，但核心逻辑一致：消费者在缩短旅行半径，压缩单次支出，但出游频率没有下降。

对于OTA平台来说，这意味着两件事。第一，客单价承压是结构性的，而非周期性的。第二，用户争夺正在从“谁有更贵的酒店”转向“谁有更丰富的短途产品组合”。

Citi在报告中特别提到了同程旅行对腾讯平台的依赖。这恰恰是同程在“微度假”赛道上的核心优势——依托微信生态，触达下沉市场用户，提供低客单价、高频次的短途产品。而携程的优势在于其海外业务的结构性增长故事，Citi对其核心业务给出了20倍2026年预期市盈率，并给予20\%的溢价，理由是“结构性海外扩张”。

KC评论： Citi对携程的估值逻辑非常清晰——溢价来自海外，而非国内。国内短途游的低价化趋势，对携程来说反而是压力。真正受益于“微度假”浪潮的，是同程旅行这类深耕下沉市场、依赖社交裂变的

[... middle omitted ...]

平台意味着什么？

携程的海外业务一直被视为其估值溢价的核心支撑。但Citi的研报没有深入拆解：入境游的增长是否真的转化为携程的海外收入？还是说，这部分流量更多流向了国际OTA平台，或者被航空公司直销渠道截流？

报告提到，携程将在6月25日发布2026年第一季度业绩，Citi预计管理层对第二季度的指引将偏保守，原因是油价和调查的潜在影响。这里的“调查”具体指什么，Citi没有展开，但显然是一个需要关注的风险点。

*这恰恰是这份报告最值得追问的地方：跨境数据的亮眼表现，能否对冲国内人均消费下滑带来的压力？Citi的答案是“维持买入”，但逻辑链中有一个关键环节没有完全闭合。**

Citi对OTA板块的判断，建立在“国内旅游量增价稳”的假设之上。但连续五个假期人均消费为负，这个趋势何时能逆转？

报告没有给出明确的拐点判断。Citi的分析师更多是在描述现状，而非预测未来。这本身是严谨的——没有人能准确预测消费趋势的拐点。但对于投资者来说，这意味着需要关注以下几个变量：

\subsection{[R028] JPM：香港楼市最大的风险不是加息，而是恒指}
{\small\color{refgray}归类：中国消费与地产：K型分化加速，定价权成为核心考验}

香港房价年内已反弹超过10\%，达到JPM全年预测区间的上沿。但这家投行在最新报告中指出，下半年动能可能显著放缓。最值得关注的判断是：市场普遍担忧的资本外流管制和加息预期，可能都不是真正的风险。真正的下行风险，是香港股市的持续疲弱。

过去一个月，恒生指数从高点下跌约15\%。按照历史上房价滞后股市3-6个月的经验，这一轮股市回调对楼市的影响，可能才刚刚开始被计入。

JPM监测的多项高频指标出现了不一致的信号。这是判断市场是否正在转向的关键窗口。

看房预约量近期反而走强，过去三周超过570组，创下年内新高。这说明潜在购房者的兴趣并未消退。然而，实际的成交转化率在下降：一手新盘的首日售出率从此前的超过70\%降至64\%；35大屋苑的二手成交量连续两周低于60宗。

这种“看房多、成交少”的分化，指向一个典型的观望状态。购房者仍然关注市场，但在股市走弱和利率预期上升的背景下，许多人选择暂缓决策。

KC评论： 看房量是领先指标，成交是滞后指标。领先指标走强而滞后指标走弱，通常意味着市场正处于一个“临界点”之前。如果看房量随后也开始回落，那么观望情绪可能正在向实质性降温转化。完整报告中有详细的高频数据图表，可以帮助读者自行判断这一拐点的位置。

过去一个月，市场最关注的两个风险是：内地资本外流管制可能收紧，以及美联储加息预期升温。JPM的观点是，这两个因素方向性负面，但单独来看，可能不足以逆转当前的上行周期。

关于资本外流管制，报告指出，文件837并未明确提及房地产。内地居民通过每年5万美元的购汇额度购买香港房产，过去从未被允许，未来也不太可能被允许。真正的资金来源是离岸资产。报告估计，非香港常住的内地买家占交易量的5-10\%，按金额计算占10-15\%。这一比例虽然不低，但并非市场的绝对主导。

关于加息预期，JPM维持2026年不加息的判断。即便美联储加息，香港最优惠利率会跟随，但历史上加息周期（如2004-06年、2016-18年）并未必然导致楼市疲弱，前提是其他基本面因素保持健康。

KC评论： 市场往往高估政策风险的短期冲击，低估基本面韧性的长期支撑。JPM在完整报告中详细拆解了历史上加息与房价的关系，以及资本管制对内地买家的实际影响边界。这些分析对于判断当前风险溢价是否合理，非常有价值。

JPM认为，比起上述两个已知风险，更需要警惕的是恒生指数的持续弱势。

报告通过相关性分析发现，在所有房价驱动因素中，库存量的相关性最强，其次是恒生指数。恒指与香港房价的走势存在3-6个月的滞后关系。从1月高点至今，恒指已下跌约15\%。虽然尚未到金融危机级别，但如果这种弱势再持续3-6个月，对楼市情绪的压制将是实质性的。

这里有一个关键的例外需要留意： 年，恒指大幅上涨，但房价并未同步反弹，原因是当时库存高企。2011年，股市走弱，但房价仅轻微下跌。这说明，股市与房价的关系不是机械的，库存水平是重要的调节变量。

幸运的是，当前库存处于健康水平。一手市场未售库存约1.67万套，对应9个月去化周期；二手市场挂牌量对应7个月去化周期。即便假设年成交量下降20\%，去化周期也仍在可控范围内。

4. 报告尚未完全回答的关键问题：股市若持续走弱，房价调整幅度有多大？

这是整份报告最让人意犹未尽的地方。JPM明确指出了恒指走弱是最大的下行风险，但并未给出一个量化的情景分析：如果恒指再跌10\%、20\%，房价会相应调整多少？

报告提到了一个参照：当前香港地产股已从5月高点回调18\%，NAV折让50\%，处于历史均值以下1个标准差。但股价的调整与房价的调整，从来不是等比例的。

历史上，恒指与房价的弹性系数在不同周期中差异很大。1997-98年金融危机时，恒指跌幅超过60\%，房价跌幅超过60\%，两者几乎同步。2008年全球金融危机时，

[... middle omitted ...]

上1个标准差。JPM认为，当前估值已经基本定价了资本外流管制和加息这两个已知风险。

但问题在于，股市持续走弱的风险尚未被充分定价。如果恒指继续下行，地产股的NAV折让可能进一步扩大到历史极端水平。报告给出的参照是，新鸿基地产在NAV折让扩大到1个标准差以下时，潜在下行空间约至105港元（当前约113.6港元）。

选股策略上，报告偏向防御。偏好净现金或即将转为净现金的开发商，如长实集团和信和置业。新鸿基地产年内已跑赢恒指约28个百分点，短期可能面临获利回吐压力。高杠杆的新世界发展和恒基地产，在利率敏感度上更高，短期可能跑输。

报告列出了五个可能驱动地产股重新跑赢的条件：恒生指数的可持续反弹；加息预期的下降（利率按兵不动即可接受）；HIBOR的回落；内地买家减少背景下，销售和房价仍能保持强劲的证明；以及8-9月业绩期对派息率和租金增长的更强劲指引。

这里面，最核心的变量还是恒指。其他四个因素，要么是恒指反弹的结果，要么是在恒指疲弱背景下难以独立发挥作用的辅助条件。

\subsection{[R029] JEF：中国大模型正在用“成本-智能”剪刀差改写全球AI竞争格局}
{\small\color{refgray}归类：AI/算力：从“光学狂欢”到“网络务实”，半导体拥挤度达历史极值}

JEF：中国大模型正在用“成本-智能”剪刀差改写全球AI竞争格局

当市场还在争论AI泡沫何时破裂时，一份来自JEF的研报揭示了一个正在发生的结构性变化：中国大模型在智能水平上已接近美国同级产品，但API成本仅为后者的一个零头。这个“成本-智能”剪刀差的持续扩大，正在从根本上改变全球AI产业的价值分配逻辑。

这不是一个关于“追赶”的故事，而是一个关于“重新定义竞争规则”的故事。

这份长达50多页的研报通过追踪OpenRouter、Artificial Analysis等多个数据源的Token消耗、模型定价和智能指数，给出了一个清晰但尚未被市场充分定价的判断：中国AI产业链正从“成本优势”向“规模优势”切换，而字节跳动的火山引擎大会、腾讯的微信AI内测、以及阿里巴巴蔡崇信在VivaTech上释放的信号，都在指向同一个方向——2026年下半年可能是中国AI商业化的关键拐点。

1. 中国模型在Token消耗上已反超美国，但更关键的是“效率优势”而非“规模优势”

根据OpenRouter在6月15日至21日当周的数据，中国模型Token消耗量达到18.8万亿，而美国模型为5.8万亿。DeepSeek V4 Flash以4.94万亿的周消耗量位居榜首，小米MiMo-V2.5和MiniMax M3紧随其后。

这个数字本身并不令人意外——中国拥有更大的用户基数和更低的使用成本。真正值得关注的不是消耗量，而是这个消耗量背后的“效率逻辑”。

JEF的Silicon Data LLM Token Expenditure Index在6月14日至19日期间维持在1.64至1.68之间，较5月31日的2.04有明显下降。这个指数衡量的是Token消耗的“质量”——当用户向高成本模型迁移时指数上升，反之下降。指数的下降意味着用户正在大规模转向低成本模型，而中国模型恰恰是这个趋势的最大受益者。

KC评论： 这个指数下降的信号比Token总量增长更重要。它说明AI使用正在从“尝鲜性消费”转向“实用性消费”——用户不再盲目追求最强模型，而是选择“够用且便宜”的模型。这对中国AI公司是结构性利好，因为他们在这个象限的竞争力远超美国同行。完整报告中有超过50个数据点追踪这个趋势，包括按子行业拆分的用户行为变化，值得细看。

2026年斯坦福AI指数报告显示，美国顶级模型领先中国模型仅2.7\%。在Arena Elo评分中，Anthropic（1503）、xAI（1495）、Google（1494）、OpenAI（1481）位居前列，而阿里巴巴（1449）和DeepSeek（1424）已经进入第一梯队。

然而，在API定价上，差距是数量级的。以输入价格为例，阿里巴巴Qwen3.7-Max的推广价为每百万Token 6元人民币，而同类美国模型的定价远高于此。这种“智能接近、成本悬殊”的组合，正在重塑企业客户的采购决策逻辑。

JEF的研报明确指出，中国模型的成本优势源于架构创新：MoE（混合专家模型）、GQA（分组查询注意力）、稀疏注意力以及更高的MFU（模型算力利用率）。这些不是简单的“复制降本”，而是从底层架构出发的效率革命。

3. 火山引擎大会的三个关键信号：商业化正在从“讲故事”进入“算账”阶段

6月23-24日，火山引擎FORCE大会即将召开。JEF列出了五个关注点，其中三个最具商业含义：

第一，豆包大模型的日Token消耗量。3月26日披露的数字是120万亿，这个数字在三个月后的变化直接反映用户粘性和使用深度。

第二，商业化计划的更新，包括各种订阅套餐的定价。这是字节跳动从“烧钱获客”转向“回收成本”的关键信号。

第三，视频生成模型Seedance 2.0的ARR趋势。视频生成是AI应用中算力消耗最高、商业价值最大

[... middle omitted ...]

通、信息搜索、生产力、娱乐购物和工具生成等多个场景。

JEF将微信AI称为腾讯在A2A（Agent-to-Agent）领域的重要一步，并预计正式发布可能在2026年第四季度。

这个判断值得重视。微信拥有超过13亿月活用户，是中文互联网最大的“场景入口”。如果微信AI能够自然嵌入用户的日常沟通、信息获取和任务执行中，它可能成为A2A时代最大的“Agent分发平台”之一。

但这里有一个尚未被充分讨论的问题：微信AI的商业模式是什么？是订阅制、广告驱动，还是通过Agent调用抽成？JEF的研报没有给出明确答案，这部分需要从完整报告中的定价表和相关讨论中进一步推导。

5. 阿里巴巴的50万亿美元TAM：是“宏大叙事”还是“可操作的路径”？

蔡崇信在VivaTech 2026上抛出了一个惊人的数字：AI的TAM是50万亿美元，相当于全球GDP的一半。他认为中国在AI供应端投资不足，而阿里巴巴凭借其自由现金流和全栈能力（芯片、基础设施、Qwen模型、应用）处于有利位置。

\subsection{[R030] 摩根斯坦利：MS：中国消费股的真正考验不是需求，是定价权}
{\small\color{refgray}归类：中国消费与地产：K型分化加速，定价权成为核心考验}

当市场还在争论消费复苏的力度时，MS最新发布的这份中国消费者行业风险回报更新报告，给出了一个更冷静的判断：对于纸巾、卫生巾、护肤品这些看似刚需的品类，驱动股价的核心变量已经不是需求量的增长，而是企业在渠道碎片化时代能否守住定价权。

这份覆盖恒安国际、中顺洁柔、上海家化三家公司的研报，表面上是在调整盈利预测和目标价，但背后揭示了一个更深层的行业结构性变化——线上渠道和新零售的渗透率提升，正在系统性地侵蚀消费品的平均售价和利润空间。三家公司无一例外地面临同样的困境：销量增长尚可，但价格承压，盈利修复慢于预期。

报告中最值得关注的信号不是某个公司的评级变化，而是MS对三家公司2026/27年净利润预测的大幅下调——中顺洁柔被砍20\%/27\%，上海家化被砍35\%/39\%。这组数字背后，是投行对中国消费品公司盈利质量的一次重新定价。

过去几年，中国消费品市场的一个主流叙事是“集中度提升，龙头企业受益”。但MS的这份报告揭示了一个反直觉的现实：即使行业集中度在提升，龙头企业的定价权并未同步增强。

以恒安国际为例，这家纸巾龙头正在受益于中小厂商退出市场，市场份额持续扩大。但报告明确指出，销量增长的主要驱动力来自中小厂商退出后的份额承接，而非品牌溢价带来的自然增长。更关键的是，线上渠道和新零售渠道占比的持续上升，正在成为平均售价的持续拖累。

这意味着什么？当一个行业的增长主要靠“吃掉别人的份额”而非“让消费者为品牌多付钱”时，规模的扩张并不自动转化为利润的增长。恒安的纸巾业务预计2026年收入增长5\%，但毛利率从2025年的23\%微降至22.5\%——量在涨，价在跌，利润率原地踏步。

KC评论： 这里有一个容易被忽视的洞察——线上渠道对消费品公司来说，既是增长引擎，也是利润杀手。报告没有直接说，但我们可以推断：当一家公司的线上占比持续提升，而线上渠道的促销力度和价格透明度远高于线下，这家公司的长期毛利率天花板就会被压低。完整的研报里包含了更详细的渠道利润拆解，这张图表值得细看。

恒安国际是这三家公司中评级最温和的（等权重），支撑其股价的核心逻辑不是增长，而是每年每股1.4港元的派息承诺，对应约6\%的股息率。MS认为这提供了下行保护。

但问题在于：当一家公司的盈利本身在2026年预计下滑3\%，股息的安全性依赖于现金流而非利润增长。报告显示，恒安的经营利润率从2025年的10.6\%预计降至2026-28年的9.5-9.6\%区间，净利率从11\%降至 \%。盈利质量在缓慢但持续地下降。

对于一个以股息为核心卖点的股票，投资者需要回答一个关键问题：如果盈利继续承压，公司是否愿意维持派息？恒安的承诺是“每年1.4港元以上”，但报告没有给出这个承诺的可持续性压力测试。

KC评论： 恒安10倍2026年市盈率、6\%股息率，看起来不贵。但“不贵”和“有吸引力”是两回事。报告隐含的判断是：在没有盈利加速的证据之前，估值很难扩张。这意味着投资者买入恒安，本质上是在买入一个收益稳定的债券替代品，而非成长股。但债券替代品面临的风险是，如果盈利持续下滑，股息率的安全边际会收窄。完整的报告对恒安的DCF估值假设和敏感度分析有更详细的拆解，这些细节对判断安全边际至关重要。

中顺洁柔是这份报告中评级最悲观的（减持），MS给出的理由是：24倍2026年预期市盈率的估值，与有限的定价权、低净资产收益率和微薄的利润率改善空间之间，存在明显的不匹配。

报告的核心逻辑链条值得仔细推敲：中顺洁柔的收入增长同样依赖销量和份额承接，而非提价。线上渠道的促销力度和更激进的价格策略，持续稀释平均售价。毛利率虽然受益于产品高端化和效率提升，但纸浆成本的不确定性和渠道竞争限制了上行空间。更关键的是，盈利复苏虽然“可见”，但“温和”——运营杠杆被持续的品牌和电商投资

[... middle omitted ...]

柔经营费用率变化的图表，可以很清楚地看到这个趋势。如果你想理解为什么很多消费股“增收不增利”，这张图是一个很好的起点。

上海家化的情况更为极端。这家拥有六神、佰草集、玉泽等知名品牌的公司，正在经历一个“品牌投资驱动”的增长阶段。线上和抖音渠道占比的提升在战略上是正确的，但代价是高昂的内容投入、流量成本和品牌建设费用。

MS将上海家化2026/27年净利润预测大幅下调35\%/39\%，而收入预测实际上调了5\%——这意味着收入在改善，但利润改善的速度远低于预期。原因很明确：销售费用率将保持高位。

报告给出的核心矛盾是：上海家化正在做正确的事（重建品牌、拥抱线上），但这些正确的事在短期内会压制盈利。而市场给这只股票的估值是35倍以上2026年预期市盈率——对于一个盈利可见性有限、正在经历转型阵痛的公司，这样的估值很难找到支撑。

这里有一个值得深思的问题：当一家公司为了长期健康而牺牲短期利润时，市场应该给予多少耐心？MS的回答是：在盈利可见性改善之前，估值风险偏向下行。

\subsection{[R031] 摩根斯坦利：MS：PHEV关税担忧被过度定价，中国车企的缓冲比市场想象的厚}
{\small\color{refgray}归类：中国新能源车：6月订单透支7月，行业整合远未到终局}

摩根斯坦利：MS：PHEV关税担忧被过度定价，中国车企的缓冲比市场想象的厚

中国汽车板块近期的一次集体回调，表面看是欧盟可能将反补贴关税扩展至插电式混合动力车（PHEV）的传闻所致。但MS在6月22日发布的一份研报中给出了一个值得深思的判断：当前股价下跌已经price in了最坏情景，而中国车企的适应能力，很可能比市场预期的要强。

这份报告的核心主张并非“风险不存在”，而是“市场对风险的定价已经过头”。在港股恒指仅下跌0.5\%的背景下，主要中国车企股价跌幅达到3\%至5\%。MS认为，这种跌幅更多来自情绪和资金面的共振，而非基本面对PHEV关税风险的直接映射。

我们仔细拆解这份报告后认为，它提供的不是一个简单的“抄底”信号，而是一个重新审视中国车企全球化能力的分析框架。关税风险是真实的，但车企的应对工具箱——从渠道囤货到本地化生产——也比以往更丰富。

MS指出了一个被市场忽略的关键细节：股价跌幅最大的车企，并非那些PHEV出口欧盟占比最高的企业，而是南向资金持仓比例较高的公司。

这个观察非常敏锐。它意味着今天的抛售，更多是资金层面的“被动减仓”或“避险性卖出”，而不是投资者对特定公司PHEV敞口的精准定价。如果市场真的在理性定价关税风险，那么比亚迪、上汽、吉利这些PHEV出口占比较高的公司应该领跌；但报告数据显示，跌幅分布与南向资金持仓高度相关。

KC评论： 市场有时会“先开枪，再画靶子”。当负面新闻出现时，流动性好的股票往往最先被抛售，哪怕它的实际风险敞口并不大。这份报告提醒我们，区分“情绪驱动的下跌”和“基本面驱动的下跌”，在决策中至关重要。完整报告中还包含一张各车企PHEV出口欧盟占比的详细对比图，能帮助读者更直观地判断哪些公司真正暴露在风险之下。

这个结论的另一层含义是：如果后续没有实质性关税落地，这些被情绪错杀的股票存在修复空间。即使关税最终落地，当前股价也已经提前消化了部分利空。

2. PHEV关税若落地，对中国车企的冲击比BEV关税更复杂

要理解PHEV关税的潜在影响，必须先回顾2024年欧盟对中国纯电动车（BEV）加征关税后的市场反应。

当时，欧盟对BEV加征了7.8\%至35.3\%的额外关税（叠加10\%的基础关税），但PHEV被豁免。这一政策差异直接改变了中国车企的出口产品结构。根据报告引用的MarkLines数据，PHEV在中国对欧盟出口中的占比，已从2024年的6\%迅速攀升至2026年的24\%。

第一，PHEV已经成为中国车企对冲BEV关税风险的“逃生舱”。如果这个逃生舱也被关闭，中国车企将失去一个重要的战略缓冲。

第二，PHEV关税的冲击可能比BEV关税更复杂。报告指出，BEV关税对中国整体出口动能的抑制效果有限，只是改变了出口产品的组合。但PHEV的情况不同——它目前是中国车企在欧盟市场增长最快的细分领域，一旦加税，对利润率的侵蚀可能更直接。

不过，MS也提出了一个关键疑问：如果欧盟真的要对PHEV加征关税，它会采取与BEV相同的方式吗？ 报告没有给出明确答案，但暗示欧盟可能会调整策略，因为BEV关税的经验表明，简单加税并没有阻止中国车企的扩张。

KC评论： 这里有一个值得追问的问题：欧盟是否会区分“中国品牌PHEV”和“中国制造但属于欧洲品牌的PHEV”？如果关税设计更精细，对不同车企的影响差异可能非常大。完整报告对这一问题有更深入的假设情景分析。

3. 短期渠道囤货可能带来“抢装”行情，但窗口期只有1-2个月

报告给出了一个对短期交易者极具参考价值的判断：由于关税细则在1-2个月内不太可能确认，中国车企可能会选择加速向欧盟渠道发货（channel stuffing），就像2024年BEV关税落地前那样。

这意味着，未来几周，中国对欧盟的PHEV出口数据可

[... middle omitted ...]

把未来的需求前置，并不能解决关税一旦落地后的长期竞争力问题。

4. 车企的分化将加剧：谁的全球化能力更厚，谁就能扛过这一轮

报告对不同车企的PHEV敞口进行了量化分析。根据MarkLines数据，上汽、奇瑞和比亚迪的海外销量中，超过20\%来自在欧盟销售的PHEV，其次是吉利和零跑。而蔚来和小鹏的直接PHEV敞口有限，其股价下跌更多是受整体市场情绪拖累。

但MS的核心判断并非“PHEV敞口高的公司风险更大”，而是“那些已经在欧盟或海外布局了生产基地的车企，更有能力消化关税冲击”。报告明确提出，从BEV关税到整个新能源车（NEV）关税的升级，可能会加速中国车企在欧盟本地化生产，或从中国以外的工厂进行出口。

在这个框架下，MS维持了对小鹏（因其本地合作伙伴关系）、上汽（因估值重估机会）、比亚迪和吉利（因多生产基地布局）的增持评级。这份报告实际上在说：关税风险是一个筛选器，它会加速淘汰那些只有出口、没有本地化能力的车企，而头部玩家的竞争优势反而可能因为行业洗牌而强化。

\subsection{[R032] 摩根斯坦利：MS：中国生物医药的“全球叙事”正在换剧本}
{\small\color{refgray}归类：区域市场：亚洲太空经济被低估，印度制造业成本驱动型修复}

2026年6月17日至18日，MS在上海举办了首届中国生物医药研讨会。这场以“连接生物医药前沿：中国与全球创新”为主题的活动，聚集了超过50家中国生物制药/生物科技公司、20多家全球药企以及10余家PE/VC机构的代表。

这不是一场普通的行业聚会。它发生在全球对中国创新药资产定价逻辑正在发生根本性转变的关键节点。过去三年，市场对中国生物医药的关注高度集中在ADC、双抗和GLP-1这三个工程化能力强的赛道。但这份研报的核心判断是：全球对中国创新的认知，正在从“工程能力验证”进入“生物学差异化”的第二阶段。

换句话说，中国生物医药的价值叙事，不再只是“我们做得更快、更便宜”，而是“我们开始思考不同的问题”。

1. 创新1.0的“工程红利”已经定价完毕，但2.0的生物学风险尚未被充分理解

中国生物医药创新1.0的核心优势——工程能力、低成本开发、执行速度——已经成功让全球药企和投资人相信，中国是药物研发价值链上不可忽视的一环。MS研报明确指出，这一阶段的成果已经基本被市场消化和定价。

但真正让这次研讨会值得关注的是，全球药企和PE/VC代表开始认可中国生物医药从“快速跟进”向“差异化生物学和药物形式”的转变。First-in-class（首创）创新正在涌现，尽管仍然处于早期阶段。

KC评论： 这听起来像一句正确的废话，但它的投资含义很具体。如果市场已经给“中国速度”打了分，那么未来超额收益的来源，将取决于谁能真正做“全球新”的生物学。但报告也坦诚地指出，首创创新意味着更高的生物学风险——投资人可能还没有完全意识到这种风险的量级。这恰恰是完整报告中最值得细读的部分：它列出了哪些疾病领域和技术平台被PE/VC认为最有潜力，以及哪些风险是当前市场定价尚未充分反映的。

PE/VC投资人小组重点关注的五大机会方向包括：慢性病的超长效给药平台、肝外递送（如siRNA和saRNA）、生物药/多肽的口服替代方案、体内基因编辑以及细胞疗法。这些方向覆盖了肿瘤、中枢神经系统、免疫和炎症等多个治疗领域。

2. BD交易模式正在从“一次性卖资产”升级为“共建价值”，但谈判桌上有新的规则

传统的外许可模式仍然是主流，尤其适用于临床阶段、开发路径清晰的资产。但研讨会传递出的更重要的信号是：随着全球药企对中国发现平台的信心逐步建立，NewCo、共同开发（co-co）和战略合作等模式正在快速崛起。

这意味着中国生物医药公司不再只是“资产出售方”，而是有机会成为“价值共享方”。恒瑞医药就是一个典型案例：其GLP-1资产通过NewCo模式“Kailera”在18个月内完成了两轮私募融资和纳斯达克IPO，恒瑞除了获得里程碑付款和特许权使用费外，还保留了长期股权收益。

但交易模式的升级也带来了更高的“入场门槛”。全球药企在研讨会上明确列出了常见的“交易杀手”：薄弱的科学依据、不成熟的数据、知识产权风险、糟糕的全球可转化性、CMC（化学、制造和控制）的不确定性以及复杂的跨境结构。

KC评论： 这其实是一个“反常识”的信号。很多人以为中国药企出海的最大障碍是地缘政治，但这次研讨会上，全球药企代表认为资产层面的许可交易相对安全——前提是公司结构和知识产权已经做好了尽职调查准备。真正被频繁提及的deal-breaker，反而是科学和临床设计层面的基本功。换句话说，地缘政治风险是“系统变量”，而科学严谨性才是“个体变量”。完整报告里详细拆解了全球药企在谈判中真正看重的条款权重排序，以及中国公司如何在临床设计阶段就为BD做好准备——这部分对任何有出海意图的生物科技公司创始人来说，几乎是必读的操作指南。

对中国公司而言，最好的准备方式是以BD思维来设计临床试验：采用高标准对照组、确保数据可以迁移到美国或全球监管包、尽可能包含多样化患者群体、保持

[... middle omitted ...]

优势在于庞大的工程人才库、高效的CRO基础设施和快速的湿实验执行能力，这些与AIDD形成了完整的药物发现闭环。

但关键瓶颈同样明确：专有数据访问权（包括失败实验的数据集）和生物学知识。单纯拥有模型能力而没有资产验证的公司，将面临越来越激烈的竞争。

真正能够建立更高竞争壁垒的AIDD公司，是那些能够整合AI、自动化、转化生物学和可信的内部资产，并形成清晰商业化模式的企业。

4. 恒瑞的全球化路径提供了一个可复制的范式，但核心在于“端到端能力”

恒瑞医药管理层在研讨会上展示了其全球化路径，这被MS视为中国生物制药全球化的关键案例。恒瑞的路径核心在于：凭借端到端的开发和商业化能力，通过灵活多样的BD模式快速拓展全球版图。

除了前面提到的GLP-1 NewCo案例，恒瑞还与GSK和百时美施贵宝达成了两项广泛的战略合作，覆盖了早期I期、IND前和发现阶段的一系列资产。这种共同开发（co-co）安排可以利用双方的专业知识来解决靶点、开发或商业化相关的瓶颈，加速进入临床。

\subsection{[R033] 摩根斯坦利：MS：6月订单透支7月，新能源车市的分水岭不在销量，在“催熟”与“真转化”}
{\small\color{refgray}归类：中国新能源车：6月订单透支7月，行业整合远未到终局}

摩根斯坦利：MS：6月订单透支7月，新能源车市的分水岭不在销量，在“催熟”与“真转化”

6月最后一周，中国新能源车市的周订单数据出现了一个值得警惕的信号。根据MS最新发布的周度渠道调研，6月15-21日当周，主要品牌订单总量环比大幅攀升，但推动力高度依赖两个因素：季度末促销冲刺和新车型集中交付转化。MS分析师Tim Hsiao团队在报告中明确警告，这种订单前置正在“inflating bookings ahead of a likely July air pocket”——即6月的繁荣可能正在透支7月的需求，形成典型的“订单空气包”。

这不是一个简单的季节性波动问题。它指向的是中国新能源车市场正在经历的结构性分化：真正依靠产品力实现“自然转化”的品牌，与依赖促销和订单前置“催熟”的品牌，正在走向截然不同的命运。这份报告的核心判断是，新车型驱动的转化是可持续的，而促销驱动的订单前置将在7月暴露真实需求底色。

1. 季度末促销正在制造一个“虚假繁荣”的6月，7月订单大概率显著回落

MS的调研数据显示，6月15-21日当周，比亚迪订单达到 万辆，环比暴增61\%，同比也有30\%的增长。这一数字表面上看极为强劲，但报告明确指出，支撑这一增长的核心因素是唐系列新车型的预订单转化，以及季度末更为激进的促销活动。

“Pull orders forward”是这份报告的关键词。季度末促销的本质，是将原本可能在7月或更晚时间发生的购买决策，通过折扣和限时优惠提前锁定。这种做法在财务报表上会美化6月的交付和订单数据，但代价是7月的需求池被提前抽干。MS用“air pocket”来描述7月的预期状况，这个词在航空术语中指气流导致的突然失重，在汽车行业语境下，意味着订单量可能从6月的高点骤降。

KC评论： 季度末促销对车企来说像一剂兴奋剂，短期数据好看，但药效过后是更深的疲态。对于投资者而言，6月的同比环比增长数据需要谨慎解读——真正重要的是7月的订单能否稳住。这也是为什么MS将行业评级维持为“In-Line”，而非上调至“Attractive”。

2. 比亚迪和零跑证明：新车型才是订单转化的“硬通货”，促销只是辅助

在促销驱动的市场环境中，仍有品牌展现出不同的增长逻辑。比亚迪当周订单的61\%环比增长，核心驱动力来自“Great Tang”新车型的预订单转化。零跑汽车同样表现突出，订单达到 万辆，环比增长43\%，同比增长93\%，其增长动力来自C系列车型的改款。

这两个案例的共同点在于：新车型或改款车型带来的订单转化，是“增量”而非“前置”。消费者因为产品本身的吸引力而做出购买决定，而非因为限时优惠而被迫提前消费。这种转化具有更强的可持续性，因为它不依赖于时间窗口的挤压。

报告也给出了反例。理想汽车当周订单仅6.5-6.7万辆，环比增长16\%，但同比下降28\%。MS明确指出，L系列车型的销售受到抑制，原因是消费者在等待当天发布的2026款L8。这意味着，即使理想汽车有新车发布计划，但在发布前的“空窗期”，老车型的需求被自然冻结。这不是促销能解决的问题，而是产品周期管理的必然结果。

KC评论： 新车型转化与促销转化的本质区别在于，前者是“创造需求”，后者是“转移需求”。投资者在评估车企订单数据时，应该区分这两类增长的占比。MS的报告虽然没有给出每个品牌的具体拆分，但通过对比比亚迪和零跑vs其他品牌的增长质量，已经给出了清晰的判断框架。

3. 蔚来、小鹏、小米：分化加剧，谁的“第二曲线”能撑起估值？

蔚来当周订单 万辆，环比下降10\%，但同比暴增190\%。摩根的评论值得注意：Onvo品牌在L60发布后的销量出现下滑，而ES9的交付可能在6月超过7000辆。这揭示了一个关键问题——蔚来的多品牌战略正在经历“此消彼长”的

[... middle omitted ...]

入局者，小米仍然保持了不错的增长势头，但环比的个位数增长暗示，SU7的订单爆发期可能正在过去，进入平稳爬坡阶段。对于小米而言，真正的考验在于产能爬坡速度和后续车型的发布节奏。

特斯拉中国当周订单 万辆，环比增长5\%，同比下滑13\%。在MS的覆盖中，特斯拉由分析师Andrew Percoco负责，这意味着这份报告对特斯拉的观察更多是作为市场参照物。特斯拉订单的“largely flat”状态，反映了两个事实：第一，特斯拉在中国市场的品牌号召力仍然存在，订单没有出现崩塌式下滑；第二，在缺乏新车型和重大价格调整的情况下，特斯拉的增长动能明显落后于比亚迪和中国本土新势力。

KC评论： 特斯拉的订单平稳但缺乏弹性，这在中国新能源车市场快速迭代的背景下，是一个危险信号。当比亚迪、零跑、华为系等品牌不断推出新车型和改款时，特斯拉依靠Model 3和Model Y的“老本”能撑多久？MS没有直接给出答案，但这恰恰是完整报告中值得深入阅读的部分——特斯拉的估值逻辑是否需要重新定价？

\subsection{[R034] 摩根斯坦利：半导体设备的下一个瓶颈不是芯片，是连接}
{\small\color{refgray}归类：AI/算力：从“光学狂欢”到“网络务实”，半导体拥挤度达历史极值}

这份报告最值得看的判断不是DRAM涨价，也不是设备出货量超预期。摩根斯坦利在2026年6月的日本半导体设备月度报告中，提出了一个真正具有结构意义的观点：当Agentic AI从概念走向部署，整个计算架构的瓶颈正在从芯片算力转向内存带宽，并进一步转向设备间的互联能力。这意味着，过去三年市场习惯的“买设备就是买AI基建”的逻辑，正在被重新定义。

市场目前对日本半导体设备板块的定价，基本建立在“AI拉动先进制程、存储扩产”这条主线上。但这份报告揭示了一条更隐蔽、也更值得追踪的线索：Agentic AI的工作负载模式，与Chatbot时代有本质区别。后者是单次请求-响应，前者需要AI代理在数十到数百个循环中反复调用上下文、检索知识、执行推理。这直接导致对DRAM和NAND的需求量级发生指数级跃迁，而非线性增长。

报告给出的数字值得反复推敲：仅NVIDIA Vera CPU的推出，就可能使2025年基础上的总DRAM bit需求增加约16\%。这不是一个增量，而是一个跳变。更关键的是，报告指出，连接性将成为下一个瓶颈——当芯片和内存的密度足够高时，信号在铜缆上能传输的距离本身就成了物理限制。

KC评论： 这意味着，市场目前对半导体设备公司的估值，可能还没有充分反映“互联设备”这一细分赛道的结构性溢价。报告没有直接给出买入建议，但它提供的框架暗示：那些在测试、互联、封装环节有技术壁垒的公司，其增长曲线可能比纯前道设备公司更陡峭。完整报告中有详细的互联设备收入拆分和竞争格局分析，值得细读。

1. Agentic AI正在改变内存需求的“量纲”，而非“斜率”

市场对AI拉动存储需求的认知，基本停留在“HBM供不应求、DDR5涨价”这个层面。但摩根斯坦利在报告中用一个简单的架构图，点明了本质差异：Chatbot模式下，一次推理只需要一次内存访问；而Agent模式下，GPU需要在一个循环中反复读取工作内存，CPU则需要承载更长的上下文计算。

这带来的直接后果是，内存需求的驱动因素从“每台服务器的容量”变成了“每个代理任务的循环次数”。后者是一个乘数效应，而非加法效应。报告估算，仅Vera CPU的部署，就能使2025年基础上的DRAM bit需求增长约16\%。这还只是一个CPU型号的拉动。

更值得注意的是，报告在NAND部分也强调了类似逻辑：Agentic AI导致更长的上下文、更复杂的数据保留需求，使得AI存储架构从“GPU溢出-内存-闪存-HDD”的四层金字塔，变成每一层的容量和带宽都需要同步扩张。这不是简单的“存储跟着计算走”，而是“存储决定了计算能否有效展开”。

KC评论： 如果这个判断成立，那么市场对存储设备公司的估值模型可能需要从“周期股”切换到“成长股”。报告中的DRAM和NAND价格趋势图、bit出货量同比拆分，是验证这一判断的关键数据。完整报告里有详细的WSTS数据拆解和摩根斯坦利的预测对比。

2. 测试设备正在从“后道”变成“瓶颈”，Advantest被上调为Top Pick的逻辑在此

摩根斯坦利在5月26日将Advantest上调为Top Pick，这一动作在报告中被放在显眼位置。但真正值得关注的不是评级本身，而是背后的逻辑链条：Agentic AI对内存和计算芯片的测试需求，不是线性增长，而是指数增长。

报告给出的Advantest收入和利润预测，已经清晰展示了这一趋势。Computing/Communication测试业务收入从2026财年的7290亿日元，预计增长到2027财年的10700亿日元，再到2028财年的15550亿日元。这背后隐含的测试机台出货量，从2026财年的2700台增长到2027财年的3835台，再到2028财年的5411台。

KC评论： 这不是一个简单的“AI

[... middle omitted ...]

判断的依据非常硬核：信号在铜缆上能传输的距离，随着数据速率的提升而急剧缩短。当数据传输速率达到200Gbps以上时，铜缆的有效传输距离可能只有几米。这意味着，在数据中心内部，设备之间的物理布局、互联架构、光模块和电模块的配合，正在成为制约系统性能的关键。

这一判断对半导体设备投资的意义在于：那些在互联相关设备——如测试接口、探针卡、封装设备、基板制造设备——有技术积累的公司，其增长逻辑将从“跟着制程走”变成“跟着架构走”。后者是一个更长期、更稳定的需求驱动。

报告没有直接点名哪些公司受益最大，但可以从其覆盖的股票中看出线索：Disco（切割/研磨设备）、Tokyo Seimitsu（探针台/划片机）、Towa（封装设备）等，都在互联链路上占据关键位置。

KC评论： 这是报告中最值得深入挖掘的部分。如果互联成为瓶颈，那么设备公司的估值逻辑将从“CAPEX周期”切换到“架构升级周期”。完整报告中有关于互联设备市场规模的测算和竞争格局分析，这些数据在公开渠道很难找到。

\subsection{[R035] NOM：端午消费数据透露的信号，比表面数字更值得警惕}
{\small\color{refgray}归类：中国消费与地产：K型分化加速，定价权成为核心考验}

消费数据正在传递一个清晰的信号：复苏的斜率在放缓，而市场对此的定价可能还不够充分。

NOM在端午假期后发布的一份消费相关研报中，给出了一个简洁但有力的判断：端午假期旅游出行和消费数据表现“平淡”（lukewarm），与五一和春节假期相比，增长动能明显减弱。这份报告的价值不在于它描述了“端午消费不太好”这个现象，而在于它通过对比不同假期的数据、不同出行方式的分化，以及商业区人流与销售额的背离，揭示了一个更值得关注的结构性变化——消费者的支出意愿正在变得更加审慎，而企业需要为此调整预期。

对于产业决策者和投资者而言，当前的关键问题不是“消费会不会突然崩溃”，而是“当消费增速从‘补偿性反弹’回归‘常态化低速增长’，哪些公司还能保持盈利韧性，哪些公司会被暴露在估值风险中”。

1. 端午数据不是孤立事件，而是消费动能系统性回落的第三个信号

NOM报告中最值得关注的一组对比，不是端午数据本身，而是它相对于五一和春节的减速幅度。

端午假期日均旅游出行量约2.16亿人次，同比增速接近零。相比之下，2026年五一假期日均出行量同比增长3.5\%，春节假期同比增长8.2\%。从8\%到3.5\%再到接近0\%，这不是单次假期的波动，而是一个清晰的减速序列。

NOM将端午数据疲弱归因于三个因素：华南不利天气、出行燃料成本上升、以及更广泛的消费领域持续疲软。天气和燃料成本是短期扰动，但“消费领域持续疲软”才是结构性因素。如果只看端午数据，很容易将增速放缓归因于天气，但结合五一和春节的趋势来看，天气只是加速了已经存在的放缓过程。

KC评论： 消费数据的减速序列比单次数据更有说服力。8\%到3.5\%到接近0\%的增速递减，意味着即使没有天气等一次性因素，消费的内生增长动能也在减弱。投资者需要问的问题是：这种减速是短期调整，还是消费进入新常态的前兆？完整报告中包含了对不同出行方式的进一步拆解，以及NOM对消费板块的选股框架，这些细节对于判断上述问题至关重要。

端午假期铁路出行量同比增长2.9\%，但增速低于五一假期的4.6\%；自驾出行量同比下降2.1\%，而五一期间自驾是同比增长2.6\%的；航空出行量同比增长0.6\%，相比五一期间的下降5.7\%有所恢复。

自驾出行的转负是最值得关注的信号。自驾通常被视为消费弹性较高的出行方式——它不需要提前购票，决策门槛低，对家庭预算的敏感度也更高。当消费者开始减少自驾出行，转而选择铁路或干脆不出行，这背后反映的可能是对油费、过路费等可变成本的敏感度上升。

航空出行的恢复则可能更多受供给侧因素驱动——五一期间航空出行下降5.7\%可能部分由于机票价格偏高，端午期间航空公司加大促销力度推动了需求恢复。但0.6\%的增速仍然微弱，说明价格弹性并没有完全释放被压抑的需求。

KC评论： 出行方式的分化本质上是一个“消费降级但未消失”的故事。消费者仍然在出行，但在选择更便宜的选项。这种“精打细算”的消费行为，对于航空、酒店、旅游等行业的定价能力意味着什么？完整报告中对商业区人流与销售额的背离分析，会进一步印证这一判断。

3. 商业区人流增长快于销售额增长，这是消费客单价下行的直接证据

NOM报告引用了商务部的数据：端午假期全国78个重点商业区的人流量同比增长4.0\%，但同口径销售额仅增长3.5\%。五一假期这两个数字分别是5.0\%和5.3\%。

人流增速高于销售额增速，意味着客单价在下降。消费者仍然在逛商场、逛商业区，但购买意愿和购买力在减弱。这是一个典型的“量增价减”格局，对于依赖客单价提升来驱动收入增长的零售企业来说，这是一个明确的警示信号。

更值得关注的是，五一假期的人流增速和销售额增速之间的差距（5.0\% vs 5.3\%）是销售额跑赢人流的，说明当时消费者的购买力仍然较强。而端午期间这个关系逆转，销售额增

[... middle omitted ...]

的多品牌矩阵和成本控制能力，使其在行业增速放缓时仍有可能维持利润率稳定。

Laopu Gold的逻辑则更具特殊性。NOM认为该股票的“风险回报比具有吸引力”，尤其是在黄金价格在地缘政治环境趋于宽松的情况下企稳的假设下。Laopu Gold作为一家黄金珠宝零售商，其估值与金价高度相关。NOM对Laopu的推荐，本质上是在押注金价不会大幅下跌，同时公司自身的品牌溢价和产品结构能够支撑其高于行业平均的估值倍数。

这两只股票的推荐逻辑共同指向一个核心判断：在消费增速放缓的环境下，市场将奖励那些有“定价权”或“利润率确定性”的公司，而惩罚那些依赖行业beta增长的公司。

KC评论： NOM的选股逻辑实际上提供了一个消费投资的观察框架：在消费增速系统性放缓的背景下，投资者应该寻找那些“即使行业增速为零，也能通过份额提升或利润率改善实现盈利增长”的公司。完整报告中包含了对这两只股票的估值方法论和风险因素的详细拆解，这些细节对于理解NOM的推荐是否适用于自己的投资组合至关重要。

\subsection{[R036] Citi：铝价缺口比市场想象的更深，中国铝业和中国宏桥的买入窗口已打开}
{\small\color{refgray}归类：能源与大宗：铝缺口被低估，黄金与2013年本质不同，航运逻辑切换}

Citi：铝价缺口比市场想象的更深，中国铝业和中国宏桥的买入窗口已打开

市场对电解铝行业的担忧主要集中在两个数字上：48百万吨和霍尔木兹海峡。前者是中国国家统计局4月年化产量引发的产能上限恐慌，后者是中东供应恢复可能带来的价格冲击。但Citi在6月22日发布的这份研报中，给出了一个与市场直觉相反的核心判断——这两个担忧都被夸大了。中国产能并未触碰天花板，海外供应恢复的速度将慢于需求修复的速度。结论是：铝价和行业利润将维持高位更久，当前股价回调是增持中国铝业和中国宏桥的窗口。

这份报告没有停留在简单的供需平衡表推演，而是通过拆解数据来源的可靠性、海外产能恢复的实际节奏、以及价格弹性对利润的杠杆效应，构建了一个从宏观商品判断到个股估值锚定的完整逻辑链。以下是核心要点的梳理与解读。

投资者最直接的焦虑来自中国国家统计局4月公布的电解铝产量数据。按该月产量年化计算，中国运行产能已逼近48百万吨，这被解读为产能天花板政策可能被突破的信号。

Citi在报告中直接否定了这一推论。该机构指出，国家统计局的月度产量数据本身波动极大，行业参与者更少使用。以5月数据为例，年化产量已回落至45.8百万吨。更关键的是，自从2025年3月以来，统计局数据计算出的月均运行产能多次出现环比下降，这与行业实际感受——利润率持续攀升带动产能利用率走高——完全背离。

KC评论： 这段话的核心价值不在于指出统计局数据有误差，而在于点明一个投资逻辑：当市场对一个“坏数据”做出过度反应时，往往提供了逆向买入的机会。Citi通过交叉验证行业数据源（Mysteel、SMM与国家统计局），确认产能政策没有变化，这就把“产能突破天花板”这个系统性风险从定价中剔除了。完整报告中还包含了对产能政策细节的进一步拆解，这里没有完全展开。

中东铝库存外流是另一个压制铝价情绪的因素。市场担心，一旦海峡复航，积压的中东铝锭将大量涌入市场，造成短期供应过剩。

Citi引用了Mysteel的估算：中东地区可运出的铝库存不足40万吨。这个数字对于全球年消费量接近7000万吨的市场而言，属于一次性冲击，而非趋势性逆转。

更大的变量在印尼。市场对印尼新增产能的担忧一直存在，但Citi指出，电力供应是印尼铝产能建设的关键瓶颈。发电机组建设需要时间，产能爬坡同样需要周期。根据Mysteel的预测，2026年全年海外新增铝产能约185万吨，实际贡献产量约120万吨。但与此同时，中东产能暂停和欧洲产能的缓慢恢复，将导致2026年海外铝产量同比减少160万吨。一进一出，2026年海外实际供给是收缩的。

Citi大宗商品团队对铝价的判断是明确的：预计2026年三季度平均铝价约为每吨4000美元。核心逻辑是霍尔木兹海峡复航后，需求恢复的速度将快于供应恢复的速度。

这一判断建立在两个关键假设之上。第一，全球铝市场在2026年将出现约200万吨的初级缺口，2027年缺口收窄至约27万吨。第二，即便考虑海外产能恢复，2027年海外产量同比增幅也只有约300万吨，不足以填补累积的库存缺口。

KC评论： 200万吨和27万吨这两个数字之间的变化，暗示了一个重要的二阶效应：2026年的缺口可能驱动价格冲高，而2027年缺口的大幅收窄则意味着价格顶部可能出现在2026年。报告没有直接讨论铝价见顶的时间点，但读者可以从这份供需推演中自己推导。这部分的时间序列假设和情景分析，在完整报告中有更细致的拆解。

4. 铝价每变动10\%，中国铝业和中国宏桥的利润弹性高达35\%和28\%

Citi在这份报告中给出了一张非常直观的敏感性分析表。以2026年预测为基础，假设氧化铝价格固定在每吨2843元，铝价每上涨10\%，中国铝业的净利润将增加35\%，中国宏桥的净利润将增加28\%。反之，如果铝价下跌10\%，两

[... middle omitted ...]

以及政策风险何时会重启

Citi的看多逻辑建立在两个核心假设之上：中国产能政策不变，以及海外供应恢复慢于需求。但报告本身也列出了几个值得警惕的风险因素。

第一，如果铝价持续超预期上涨，中国政府是否会松动供给侧改革政策？报告明确将“政府放松限产政策”列为中国铝业和中国宏桥的共同风险。第二，成本端的不确定性。氧化铝价格的波动、电力成本的上升，都可能侵蚀利润弹性。第三，印尼产能的实际落地节奏。Mysteel的预测是基于当前建设进度，但电力基础设施的延误可能改变供需平衡表。

KC评论： Citi的估值锚定在2026年，但铝行业的周期性决定了利润高位不可能无限期维持。投资者需要追问的是：2027年之后，全球铝市场是否会从短缺转为过剩？报告给出的2027年缺口只有27万吨，已经非常接近平衡。如果海外产能恢复顺利，或者中国需求增速放缓，利润高点可能比市场预期的更早到来。这些问题的答案，取决于对海外产能爬坡曲线和中国下游需求结构的更深入分析，而这正是完整报告可以提供的增量信息。

\subsection{[R037] GS：日本央行内部已有人要求加速加息，关键看通胀是否会超2\%}
{\small\color{refgray}归类：宏观与利率：美联储“沉默”重塑全球定价，日本央行内部分歧加大 / 日本市场：结构性牛市第二幕，四大驱动力同步强化}

GS：日本央行内部已有人要求加速加息，关键看通胀是否会超2\%

日本央行6月货币政策会议的意见摘要，表面上只是常规披露。但GS日本经济学家团队在最新研报中捕捉到了一个值得深读的信号：有两位委员明确主张加速加息，认为政策利率应尽快接近中性利率水平。

这并非市场主流叙事。当前全球投资者的注意力仍集中在日本央行何时、以何种节奏缩减购债规模，以及日元汇率是否会继续承压。但GS指出，真正需要跟踪的，是“基础通胀超过2\%的风险评估”——这个判断将决定少数派的加速加息主张是否会扩散为央行共识。

换句话说，日本央行下一次加息的时间窗口，可能比市场预期更早，也可能比市场预期更晚。而决定权，掌握在通胀数据手里。

6月日本央行的政策决定本身并不令人意外：加息至1\%，并从下一财年起停止缩减日本国债购买规模。市场对此已有充分预期。

一位委员明确反对加息，认为提高政策利率会抑制企业固定投资，进而同时压低通胀、产出和就业。这位委员的立场代表了对经济复苏可持续性的担忧。

但更多的委员支持加息，理由是两个：经济下行风险已经减弱；基础CPI存在向上偏离2\%的风险。

GS的解读是：支持加息的理由中，“通胀超调风险”已经取代“经济过热风险”，成为央行决策的核心变量。这意味着，日本央行的加息逻辑正在从“预防性”转向“应对性”。

KC评论： 这不是日本央行第一次讨论加息，但这是第一次在意见摘要中出现“基础通胀可能超过2\%”的系统性担忧。对投资者来说，这意味着日本央行的加息决策将越来越依赖实际通胀数据，而非经济增速。谁在跟踪日本的核心CPI分项，谁就能预判下一次加息时点。

两位委员分别提出了明确主张：一位认为“有必要尽快将政策利率向中性利率靠拢”；另一位更直接地给出了数字——中性利率大约在2\%左右，并认为“以几个月为间隔考虑加息是合理的”。

这两个表述叠加在一起，意味着在目前1\%的政策利率水平上，这两位委员希望看到大约100个基点的加息空间。

当然，目前这只是两位委员的个人意见。GS在报告中特别强调，这些意见目前仍是少数派。但问题在于：这个少数派是否会扩大？

GS的判断是：这取决于未来对“基础通胀超过2\%风险”的评估。如果未来几个月的核心通胀数据持续高于预期，更多委员可能转向支持加速加息。

KC评论： 2\%的中性利率目标并不是GS的预测，而是央行内部委员的表述。但这一表述本身就很有价值——它给出了一个可以跟踪的“政策目标区间”。如果市场开始相信日本央行确实想把利率提到2\%，那么日本国债收益率曲线的定价逻辑就要重写。完整报告中，GS对中性利率的测算方法和假设条件值得细读。

6月会议的另一项重要决定是：从下一财年起停止缩减日本国债购买规模。

这一决定并非没有争议。一位委员明确反对停止缩减，认为这会被市场解读为“财政融资”或“压低长期利率的尝试”。这是一个非常尖锐的批评——如果市场真的这么理解，日元可能进一步承压，日本国债的信用溢价也可能上升。

但大多数委员选择了暂停。理由是：日本国债市场的功能虽然有所改善，但仍需考虑市场稳定。

GS没有在摘要中给出自己的判断，但隐含的逻辑很清晰：日本央行在“正常化”和“市场稳定”之间，选择了后者。这也意味着，日本央行短期内不会主动引发国债市场动荡，但长期缩减的方向并未改变。

KC评论： 停止缩减购债，不等于日本央行重回宽松。它更像是一个“缓兵之计”——先稳住市场，再找机会继续缩减。对投资者来说，这意味日本国债的供给压力短期内不会加剧，但长期风险并未消除。完整报告中，GS对日本央行购债操作的技术细节和未来路径有更详细的拆解。

GS这篇研报最大的价值，不是给出了一个明确的预测，而是提出了一个需要持续跟踪的问题：两位委员的加速加息主张，何时会扩散到更多委员？

这个问题目前没有答案。但GS给出了跟

[... middle omitted ...]

是个别委员的表述。如果日本央行的正式研究或会议纪要中出现对中性利率的系统性讨论，说明加速加息正在进入政策议程。

第三，日元汇率的反应。如果市场开始定价更快的加息路径，日元可能提前升值。这会反过来影响日本的进口价格和通胀预期，形成一个自我强化的循环。

这三个信号都不需要复杂的模型，但需要持续跟踪。这也是为什么我们每天会在社群中更新国际投行的最新研报摘要和关键数据图表——帮助读者在信息过载中抓住真正重要的变量。

这份GS研报的价值，不在于它给出了一个确定的结论，而在于它帮助读者识别了一个正在形成的政策分歧。分歧本身不是答案，但它是所有重要变化的前奏。

如果你希望持续跟踪日本央行的政策动向，以及全球主要投行对日本市场的最新判断，欢迎加入我们的社群。每天，我们会由AI agent与人工review相结合，生成国际投行中文摘要与KC评论，日均更新约10至40页，并整理当天最新的数据图表合集。这些内容既方便你直接喂给AI模型进行深度分析，也方便你人工快速把握全球市场动态。

\subsection{[R038] Bernstein：半导体唯一的游戏还在，但拥挤度已到历史极值}
{\small\color{refgray}归类：AI/算力：从“光学狂欢”到“网络务实”，半导体拥挤度达历史极值 / 风险偏好：AI主题拥挤但需求无虞，可持续投资“喘口气”}

Bernstein：半导体唯一的游戏还在，但拥挤度已到历史极值

AI主题已经膨胀到几乎拖拽整个半导体板块前进的地步——除了一个反常的例外：AI算力公司本身。这是Bernstein最新一期半导体周期速览报告给出的核心判断。报告作者Stacy Rasgon团队用一份详尽的“Tearsheet”展示了当前周期的位置，结论是：AI需求毫无放缓迹象，但几乎所有跟踪指标都在亮黄灯。估值、拥挤度、库存天数，每一个都在警示风险，而市场的反应却是继续加注。

这份报告最值得注意的判断并非“AI还会涨”，而是：当前这轮行情的驱动力已经切换，从基本面改善转向了“只有这一个游戏可玩”的叙事惯性。 SOX指数年初至今上涨107\%，是标普500涨幅的11倍。但Forward EPS同期也上涨了75\%。这意味着估值扩张并非全部，盈利增长提供了支撑。但问题在于，当几乎所有非AI公司都在强行构建AI叙事时，市场已经进入了一个“只要沾上AI就能涨”的状态。

对于产业决策者和投资者而言，真正需要回答的问题是：这个状态还能持续多久？Bernstein的答案是——可能比大多数人预期的更久，但途中会有剧烈颠簸。

KC评论： 报告的核心洞察是“唯一可玩的游戏”，而不是“唯一的增长”。这暗示了资金流向的集中度风险。当所有人都在追逐同一主题时，一旦出现任何边际变化，调整的幅度也会是同等的。完整报告中的Tearsheet图表（Exhibit 1）用红绿灯系统标出了7个关键指标的状态，其中6个已经从三个月前的“中性”转为“负面”。这些图表值得细看。

1. 真正的“输家”是AI算力公司本身，但它们的估值却出奇便宜

这是一个反直觉的现象。Bernstein指出，当市场注意力从GPU转向内存、半导体设备、光通信、模拟/功率器件，再到最近的CPU时，AI算力公司本身反而成了“输家”——它们的股价相对于板块其他成分表现落后。但报告认为，这种落后恰恰创造了机会。

“瓶颈”概念驱动的投资逻辑有一个根本缺陷：如果NVDA和AVGO不涨，所有围绕它们构建的瓶颈叙事最终都会失效。报告明确表示，两家公司目前的估值都“不合理地便宜”。NVDA的Blackwell/Rubin产品线到2027年的累计收入预期达到1万亿美元，而当前市场对NVDA的盈利预期仍然偏低。AVGO的1000亿美元AI收入目标，在报告看来甚至偏保守。

关键数据点：NVDA当前Forward PE仅22.7倍（2026年），AVGO为33.8倍（2026年）。对于一个年增长仍保持在三位数区间的行业，这样的估值并不算疯狂。

KC评论： 报告没有展开的是“估值便宜”的前提——当前盈利预期是否足够保守。如果AI需求出现任何低于预期的信号，这些“便宜”的估值会迅速变得昂贵。完整报告中有详细的盈利模型拆解，包括每个季度的收入驱动因素，这些是判断预期是否安全的关键。

Bernstein对半导体设备板块的整体判断是“继续持有”，但承认估值已经大幅上升。报告特别强调WFE（晶圆制造设备）支出仍有上行空间，尤其是DRAM、先进逻辑和先进封装领域。AI需求最终会转化为更多的芯片、更多的晶圆和更多的设备。

在三家覆盖的公司中，Bernstein倾向于Applied Materials（AMAT），理由是DRAM敞口和相对估值优势。AMAT、KLAC和LRCX均被给予“跑赢”评级。值得注意的是，报告提到了KLAC的10:1拆股后的模型修正，但并未因此改变评级。

一个值得关注的细节：报告认为中国市场的替换风险在KLAC身上更低。这暗示了地缘政治因素仍然是设备股定价的重要变量，但不同公司的暴露程度存在差异。

Bernstein对模拟芯片板块的态度是“观望”。TXN和ADI的营收已经连续多个季度实现双位数增长，这意味着行业不仅

[... middle omitted ...]

实。NXPI虽然更便宜，但汽车敞口过大，且渠道库存正在上升。

KC评论： 模拟芯片的估值问题本质上是“复苏预期已经充分定价”。30-40倍的PE意味着市场已经将未来2-3年的增长折现到当前价格。完整报告中有详细的估值比较图表（Exhibits 55-56），展示了模拟板块相对于标普500的溢价水平。这些图表可以帮助判断当前估值是否还有安全边际。

Bernstein的Tearsheet中，最令人不安的指标是“拥挤度”。报告显示，半导体板块相对于整个市场的拥挤度已经远高于历史区间上限。三个月前这个指标还是“中性”，现在已经转为“负面”。

历史经验表明，拥挤度达到极端水平后，板块往往会出现剧烈调整。但报告没有给出一个明确的触发条件。这是当前周期最不确定的部分——当所有资金都在同一个方向时，任何边际变化都可能引发连锁反应。

报告提到了“泡沫”讨论正在增加，但认为当前还远未到真正的泡沫阶段，因为盈利增长在支撑股价。然而，如果盈利增长出现任何放缓，估值扩张的支柱就会动摇。

\subsection{[R039] Bernstein：车企造人形机器人，真正的护城河不在机器人本身}
{\small\color{refgray}归类：风险偏好：AI主题拥挤但需求无虞，可持续投资“喘口气”}

Bernstein：车企造人形机器人，真正的护城河不在机器人本身

汽车制造商正在集体涌入人形机器人赛道。从特斯拉到现代，从小鹏到小米，2026年北京车展上，多家车企的人形机器人成为展台主角。这不是跨界玩票，也不是资本市场的故事新编。Bernstein最新研报给出了一个核心判断：车企进入人形机器人领域，不是多元化，而是主业的自然延伸——它们的竞争优势恰恰来自造车本身积累的硬件、软件和规模化能力。

这份报告的价值不在于罗列谁家机器人走得更好看，而在于它拆解了一个关键问题：当所有车企都在做同一件事，谁真正拥有结构性优势，谁只是在跟风？

理解这一轮车企进军机器人，首先要跳出“机器人是新兴行业”的思维定式。Bernstein的研报从底层逻辑切入：一辆电动车和一个双足机器人，共享的零部件远比外界想象的多。

电机、减速器、传感器——这些构成机器人关节的核心部件，与新能源汽车的电驱系统高度同源。更关键的是制造端。车企现有的生产线、供应链管理能力、质量控制系统，可以直接迁移到机器人的量产过程中。这不仅仅是成本优势，更是时间优势：当纯机器人公司还在为一条产线融资时，车企已经在用现成的工厂做原型迭代。

KC评论： 这份报告最值得细读的是它对“复用”的量化分析。读者可以重点关注研报中关于零部件重叠度和制造成本对比的图表，它们直观展示了为什么车企的起跑线比专业机器人公司靠前得多。完整报告中还有特斯拉和小鹏的零部件拆解对比，值得深入看。

从投资角度看，这意味着：评估一家车企的机器人业务前景，不能只看它机器人本身的技术参数，更要看它能把多少造车能力“平移”过去。那些垂直整合程度高、核心零部件自研比例大的车企，天然拥有更深的护城河。

硬件复用是显性优势，但Bernstein的研报更强调一个容易被忽视的层面：软件和算法的部分重叠。

自动驾驶和人形机器人的底层技术架构正在趋同。VLA模型（视觉-语言-动作模型）和世界模型，是两者共同的技术基石。特斯拉的Optimus和FSD共享同一套感知和决策框架，小鹏的IRON机器人与其自动驾驶系统在AI模型层面也有大量交叉。

但这里有一个关键区别：不是所有车企的自动驾驶能力都足够强。Bernstein在报告中对各家车企的自动驾驶水平做了分级，这直接决定了它们在机器人软件栈上的起点。

KC评论： 报告没有展开的一个二阶问题是：如果一家车企的自动驾驶本身还处于追赶阶段，它所谓的“软件复用”到底能复用多少？这个问题在完整报告的技术对比部分有更细致的拆解，包括各家的模型架构差异和训练数据来源。建议读者拿到完整报告后重点看“Software Overlap”那一节的图表。

这意味着，机器人业务的长期竞争力，很大程度上取决于车企在自动驾驶领域的真实积累。那些在自动驾驶上已经建立技术壁垒的公司，在机器人赛道上的优势会随时间放大。

Bernstein的研报提出了一个务实的观点：人形机器人的第一个大规模应用场景不会是家庭，而是工厂。

这恰恰是车企的天然主场。特斯拉已经在自己的工厂内部署Optimus执行零件搬运任务，宝马的Spartanburg工厂用Figure机器人完成了超过3万辆车的生产辅助工作，现代旗下的波士顿动力正在将Atlas机器人引入重载制造环节。

第一，内部部署意味着车企可以在真实生产环境中快速迭代，发现的问题可以直接反馈给研发团队，形成“部署-反馈-改进”的闭环。这比任何实验室测试都有效。

第二，工厂场景的ROI更容易计算。Bernstein的研报给出了一个关键判断：当人形机器人的投资回收期低于5年时，在汽车制造中的部署就具有经济吸引力。目前报告估算的回收期是7-8年，但预计未来3-5年内随着机器人成本下降和劳动力成本上升，这个阈值将被突破。

KC评论： 这里的核心变量是“回收期

[... middle omitted ...]

间维度。报告对小鹏和特斯拉给出了相对积极的评价，但同时强调“长期时间线”的不确定性。小鹏的IRON机器人计划在2026年底量产，特斯拉的目标是2026年有限商业化——这些时间表能否兑现，取决于灵巧手等关键瓶颈的突破速度。

其次是商业模式。车企进入机器人领域，最终是要卖机器人还是卖服务？特斯拉和现代倾向于硬件销售，丰田选择了RaaS（机器人即服务）模式，小鹏则在探索消费级市场。不同的商业模式对应完全不同的估值逻辑和竞争壁垒。

第三是竞争格局的演化。目前车企的优势很大程度上来自“跨界复用”，但当专业机器人公司（如Figure、Agility、宇树科技）也开始规模化量产时，这种优势还能持续多久？报告提到了这个风险，但没有给出明确的判断。

KC评论： 这些未解问题正是完整报告最有价值的部分。Bernstein在报告中提供了详细的竞争格局图谱和各家公司的技术路线对比，包括灵巧手、运动控制、AI模型等核心模块的进展。对于希望深入理解这个赛道的读者，这些细节比结论本身更重要。

\subsection{[R040] JEF：AI服务器PCB的“Kyber延迟”比想象中更值得警惕}
{\small\color{refgray}归类：AI/算力：从“光学狂欢”到“网络务实”，半导体拥挤度达历史极值}

JEF：AI服务器PCB的“Kyber延迟”比想象中更值得警惕

如果2027年你预期AI服务器会全面切换到一种全新的背板架构——Kyber，那现在可能需要重新校准预期。JEF的最新供应链调研指向一个概率极高的结论：这个架构大概率要推迟到2028年，甚至可能被取消。

这个判断的直接后果是：2027年AI PCB和CCL的全球市场规模，将比JEF原先的预测分别减少5\%和8\%。如果Kyber最终彻底消失，2028年的冲击会扩大到11\%和16\%。

表面上看，5\%和8\%的调整似乎不算灾难。但真正值得关注的不是数字本身，而是这件事揭示的产业链权力结构变化：谁在技术迭代中真正受益，谁在承担迭代失败的风险。这份报告给出的答案，对PCB和CCL领域的投资布局有直接含义。

JEF的报告指出，Kyber延迟的核心原因在于“正交背板PCB”的复杂程度超出了预期。原本计划用这种PCB替代铜缆线束来实现机柜内部的高速互联，但技术挑战至今仍未解决。供应链早在5月就察觉到端倪，到最近几周，2027年看不到Kyber已经成为一个“高度可能的事件”。

这意味着2027年的Rubin Ultra将延续Oberon架构——也就是当前的NVL72方案。供应链正在努力让Kyber在2028年以简化版（从4-canister缩减为2-canister）落地，但JEF明确表示“可见度很低”。

KC评论： 技术推迟本身不是新闻，但JEF的调研揭示了一个更关键的结构性问题：当下一代架构无法按时落地，现有的Oberon架构生命周期就被迫延长。这对铜缆连接器厂商是利好，因为他们被PCB替代的威胁暂时缓解了。但对PCB厂商来说，这意味着他们预期的“量价齐升”周期被推迟了，而竞争压力不会消失。

在整体TAM下修的前提下，JEF给出了一个看似矛盾但逻辑清晰的判断：PCB供应链整体承压，但上游材料商会继续跑赢下游板厂。

理由很直接：上游正在经历行业性的供给紧张，尤其是玻纤布和CCL。这些材料供应商在当前阶段更容易转嫁成本上涨，甚至从每次涨价中获取额外利润。相比之下，PCB制造商转嫁能力明显弱于上游。

这个判断背后的逻辑是：在技术迭代节奏被打乱的阶段，产业链的利润分配会向供给刚性更强的环节倾斜。上游材料的扩产周期长、进入壁垒高，而PCB板厂的产能相对灵活。当需求增速放缓但绝对值仍在增长时，这种结构性优势会被放大。

KC评论： 这个结论对投资者意味着，即使在AI PCB整体TAM下修的情况下，上游材料股和下游板厂的股价表现可能出现明显分化。报告没有展开的是，这种分化是否已经反映在当前估值中，以及如果Kyber在2028年最终落地，这种优势是否会反转。这些正是完整报告中值得深挖的内容。

JEF给出了两个情景的数字：2027年AI PCB和CCL TAM分别下调5\%和8\%，但如果Kyber推迟到2028年或取消，2028年的下调幅度会扩大到11\%和16\%。

关键在于，2028年的冲击不是简单的线性外推。如果Kyber被取消，意味着整个行业在 年期间将停留在Oberon架构上，而Oberon架构对PCB的规格要求和用量都低于Kyber。更低的单机用量叠加更长的产品周期，会导致PCB厂商的营收和利润增速出现双杀。

同时，JEF也强调，材料和规格的研发并没有停止。交换板和中板可能在2027年迁移到M9/10/PTFE材料，CoWoP仍然瞄准最早2027年渗透。这意味着PCB的单机价值量仍在上升，只是上升的斜率被拉平了。

4. 报告没有完全回答的关键问题：Oberon延长带来的竞争格局重塑

JEF的报告中有一个判断值得反复推敲：铜缆厂商是这一延迟的受益者，因为被PCB替代的威胁被推迟了。这隐含了一个更深层的问题——如果Oberon架构的生命周期延长2-3年，铜

[... middle omitted ...]

该如何给PCB和铜缆相关公司定价。完整报告中的图表和原始数据可以帮助读者自己推演不同情景下的概率分布。这也是为什么我们建议读者获取原始报告进行深度阅读。

基于JEF的这份报告，我们建议读者建立一个三层的观察框架，用来跟踪这一事件的演变：

第一层是技术验证信号。关注2026年下半年供应链是否会传出Kyber简化版（2-canister）的工程验证进展。如果到2026年底仍没有明确的验证结果，2028年落地概率将进一步下降。

第二层是材料迁移信号。报告提到M9/10/PTFE和CoWoP的研发仍在进行，这些材料的渗透节奏是判断PCB价值量能否独立于架构升级而增长的指标。如果2027年这些材料如期进入交换板和中板，即使Kyber推迟，PCB厂商仍有机会通过单机价值量提升来对冲量上的损失。

第三层是竞争格局信号。关注铜缆厂商在Oberon延长期间的产能扩张和客户绑定情况。如果铜缆厂商在此期间大幅提升产能并锁定长期订单，PCB厂商在下一轮迭代中的议价能力将被进一步削弱。

\subsection{[R041] 摩根斯坦利：日本电子化学品正在经历一场“材料级”的供应链重构}
{\small\color{refgray}归类：日本市场：结构性牛市第二幕，四大驱动力同步强化}

摩根斯坦利：日本电子化学品正在经历一场“材料级”的供应链重构

市场对日本电子化学品的关注，长期以来集中在光刻胶、高纯度化学品等“明星材料”上。但摩根斯坦利在2026年6月的最新月度报告中，给出了一个值得警惕的判断：这一轮产业链变化，真正的结构性机会不在最热门的单品，而在那些被市场低估的“中间层”材料——覆铜板（CCL）和封装材料。

这份覆盖日本电子化学品板块的报告，维持行业“In-Line”评级，但个股推荐出现了明确的分化。报告最核心的信号是：FC-BGA（倒装芯片球栅阵列）用覆铜板正经历“非常强劲”的需求周期，而中国市场的光刻胶需求是另一大驱动力。但这两条线索背后，隐藏着截然不同的竞争逻辑。

*KC评论：** 大多数投资者习惯用“半导体材料”这个笼统概念来覆盖日本公司。但摩根斯坦利这份报告提醒我们，电子化学品内部已经出现严重分化——材料的结构性短缺和周期性过剩并存。读完整份报告，才能真正理解为什么Resonac被列为“Overweight”，而Tokyo Ohka Kogyo却是“Equal-weight”。

摩根斯坦利报告明确指出，“FC-BGA CCLs”的表现“非常强劲”。这不是一个模糊的行业观察，而是对特定细分赛道的确认。

覆铜板是PCB的核心基材，而FC-BGA封装基板用的CCL，技术要求远高于普通PCB。随着AI芯片、高性能计算对封装密度和信号完整性的要求持续提升，FC-BGA基板的需求在过去18个月持续走高。但真正值得关注的是，摩根斯坦利判断Resonac很可能启动产能扩张，并进一步提价。

这意味着什么？在电子产业链中，材料供应商通常处于“被动响应”位置——晶圆厂或封装厂给出规格，材料商按单生产。但Resonac的提价信号表明，FC-BGA CCL的供需格局已经紧张到让供应商获得了定价权。这不是周期性的“补库存”，而是结构性需求增长带来的议价能力转移。

2. Sumitomo Bakelite的封装材料正在进入台积电与OSAT供应链，这是比业绩本身更重要的信号

报告提到，Sumitomo Bakelite当前的业务“强劲”，两个关键进展值得关注：一是颗粒状封装材料向Intel的出货量正在增加；二是液态封装材料已开始向台积电和OSAT（外包封装测试厂）提供样品。

*KC评论：** 样品供应（sample supply）阶段，意味着产品已经通过了初步验证，进入客户认证的后期。对于一家日本材料公司来说，进入台积电供应链是“质变”信号。这份报告没有展开的是，液态封装材料在先进封装中的用量正在快速提升，而Sumitomo Bakelite在这个领域的技术积累，可能比市场预期的更深。完整报告中有详细的客户认证进度和产能规划，是判断未来12-18个月业绩拐点的关键。

3. 中国市场的光刻胶需求是真实驱动力，但Tokyo Ohka Kogyo的增速可能受限

报告指出，“中国市场正在驱动光刻胶需求”。这与过去两年中国本土晶圆厂扩产、国产化替代加速的趋势一致。但摩根斯坦利对Tokyo Ohka Kogyo给出了“Equal-weight”评级，目标价9,400日元，较当前股价有约15\%的下行空间。

理由很直接：Tokyo Ohka Kogyo的增长率低于JSR，这反映了其MOR（金属氧化物抗蚀剂）产能较低。 MOR是EUV光刻工艺中的关键材料，随着3nm及以下制程的量产，MOR需求正在爆发。Tokyo Ohka Kogyo在这个领域的产能布局落后于竞争对手，这意味着它可能无法充分受益于这一轮技术升级带来的增量市场。

*KC评论：** 这是一个典型的“赛道对了，但公司不一定对”的案例。中国光刻胶需求是真实且持续的，但如果一家公司不能在EUV相关材料上建立产能优势，其增长天花板就会比同行低。报

[... middle omitted ...]

，而是专注于技术面和供需面。对于投资者来说，这是一个需要自己补充判断的关键变量。完整报告中包含了各公司在中国市场的收入占比和客户分布数据，这些细节是评估风险的基础。

5. 估值信号：Resonac的ROE改善路径是核心叙事，但市场可能低估了执行难度

从估值表看，Resonac的ROE从2025年的4.3\%预期提升至2027年的16.1\%，这是一个相当激进的改善路径。公司当前P/E（2026e）为32.6倍，P/B为4.2倍，在电子化学品板块中属于偏高水平。

摩根斯坦利对Resonac的“Overweight”评级，本质上是在押注其重组和产能扩张的成功。但这里有一个尚未验证的假设：Resonac能否在扩产的同时维持提价能力？

FC-BGA CCL的强劲需求是提价的基础，但一旦产能扩张步伐超过需求增长，定价权就会迅速消失。报告提到了“进一步提价”的可能性，但没有给出具体的幅度和时间表。这是完整报告中值得细读的部分——产能扩张计划的具体节奏，以及客户对提价的接受度。

\subsection{[R042] 摩根斯坦利：MS：AI投资正从“光学狂欢”转向“网络务实”}
{\small\color{refgray}归类：AI/算力：从“光学狂欢”到“网络务实”，半导体拥挤度达历史极值}

市场对AI基础设施的叙事正在经历一次微妙的但关键的切换。当所有人还在为光模块的订单爆发和激光器产能焦虑时，MS的最新行业反馈揭示了一个正在发生的结构性转向：投资者的关注重心，正从光通信的“需求发现”阶段，移向网络设备的“推理落地”阶段。

这份报告的核心判断并非否定光学的长期价值，而是指出一个短期但重要的节奏变化——过去一个月，网络设备板块上涨6.7\%，而光学板块下跌10\%。这不是简单的板块轮动，而是市场在用一个更务实的框架重新审视AI基础设施的投资逻辑。

光学的需求故事并未终结。MS明确指出，需求条件仍然强劲，这不是投资者当前争论的焦点。真正的辩论已经上移了一个层次：在增量产能陆续释放、技术路线尚未收敛的背景下，这些订单的利润质量到底能维持多久？

报告梳理了四个关键变量：CPO/NPO技术的落地时间与形态、激光器新玩家的产能扩张、训练与推理之间资本配置的切换对规模光学定价的影响，以及光纤现货市场的价格信号。这些因素叠加在一起，使得光学的短期可见度下降，波动性上升。

KC评论： 这其实是AI硬件投资从“炒预期”进入“验利润”的必然阶段。当所有人都知道订单会增长时，股价的边际驱动力就从“有没有订单”变成了“订单能赚多少钱”。这份报告最值得深读的部分，是它如何拆解不同光学公司的利润结构差异——不是所有光学股都站在同一块地基上。

2. 一个被市场过度解读的产能新闻：Source Photonics的扩张更像“扰动”而非“颠覆”

上周最引发市场关注的数据点，是苏州东山精密公告其子公司Source Photonics计划投资12亿美元建设光芯片及模块产能。消息一出，Lumentum和Coherent股价承压，市场担忧激光器市场将面临新一轮供给冲击。

但MS给出了一个冷静的校准：Source Photonics在激光器市场的份额仍然较小，行业反馈显示，其在竞争性讨论中并未形成实质性威胁。目前市场由Lumentum、Coherent、住友电工和博通主导。更重要的是，Source Photonics的产能扩张更可能影响的是CW激光器市场，而非利润率更高的EML市场。

这意味着，市场对这一消息的定价可能存在过度反应。但报告也承认，在缺乏现货市场数据的情况下，这种情绪驱动的波动在短期内很难完全消除。

与光学的争论不休形成鲜明对比的是，MS在网络设备领域找到了更坚实的共识。报告特别点名Arista Networks和Cisco，认为这两家公司在未来一个月内将迎来更多正面催化剂，直至财报季。

核心驱动力来自一个判断：随着AI从训练阶段向推理阶段演进，网络设备的角色将从“管道”升级为“平台”。推理场景对网络延迟、负载均衡和边缘算力的要求，使得高端交换机和智能网络方案的价值量显著提升。这不是一个简单的量增逻辑，而是一个价增逻辑——用户愿意为更好的网络性能支付溢价。

KC评论： 如果你把AI推理看作一个“服务交付”问题，而不是“模型训练”问题，网络设备的位置就从成本中心变成了体验中心。这是这份报告最有趣的洞察之一。完整报告里对Arista和Cisco的估值差异、产品组合和客户粘性有更详细的拆解，这些细节才是判断哪家公司能真正吃到这波红利的关键。

4. 光学板块中唯一“干净”的故事：Corning的逻辑为何不同

在光学板块的普遍噪音中，MS认为Corning是未来一个月“最干净的故事”。背后的逻辑值得推敲：Corning的商业模式以长期协议为主，而非现货定价。这意味着，尽管亚洲光纤现货市场的价格波动剧烈，但Corning的利润表不会受到直接冲击。

报告引用了Fujikura上调全年业绩指引这一事件——其营业利润预期上调47\%，理由是需求更强、定价更好、氢气短缺影响弱于预期。虽然MS认为亚洲现货价格的解读对Cornin

[... middle omitted ...]

重要的未解问题：如果推理成为主导需求，对光学市场的总量和结构意味着什么？

报告指出，在训练阶段，规模光学（scale-across）的需求强劲且定价相对宽松。但进入推理阶段，由于推理场景更强调成本效率和规模部署，光学组件的定价竞争可能会加剧。这意味着，即使总需求继续增长，利润率的结构可能发生变化——这可能是当前市场对光学板块定价分歧的根源。

此外，CPO技术的落地时间表仍是一个巨大的未知数。如果CPO在2027年或2028年实现规模商用，现有的可插拔光学模块市场格局将被重塑。报告认为，规模升级中额外的光学内容比其形式更重要，但这个判断本身就需要持续验证。

KC评论： 这个问题的答案，将决定光学板块未来两年的投资框架。如果推理场景下的光学需求只是“量增价跌”，那么当前的高估值需要重新审视；如果光学在推理中找到了新的价值锚点（比如更高性能的激光器或更复杂的封装方案），那么现在的波动可能就是机会。完整的报告里提供了更细致的需求拆分和敏感性分析，这些才是自己判断的基础。

\subsection{[R043] JPM：AI颠覆论被高估，数字银行龙头才是真正的护城河}
{\small\color{refgray}归类：未被主题摘要引用}

市场对AI颠覆软件行业的焦虑从未消退。当ChatGPT掀起生成式AI浪潮后，每一家SaaS公司都被迫回答一个问题：你的产品会被AI一键替代吗？在数字银行这个细分领域，JPM最新发布的研报给出了一个反共识的判断——AI颠覆风险被明显高估了。

这份覆盖Q2 Holdings、Alkami和nCino三家公司的深度报告，核心结论清晰而锐利：那些深耕银行核心工作流、数据整合与合规要求的垂直SaaS厂商，其护城河远比市场想象的深厚。AI新进入者无法绕过银行监管、历史数据访问权限和复杂系统集成这三道墙。

更关键的是，JPM不仅给出了判断，还明确指出了两条投资主线：Q2 Holdings作为有机增长的优质标的，Alkami作为高确定性的并购候选。而nCino虽然估值有吸引力，但竞争格局和AI叙事的清晰度仍逊一筹。

1. 银行SaaS的护城河不在于技术，而在于合规与集成复杂度

市场普遍担心AI原生公司会像颠覆内容创作一样颠覆银行软件。但这份研报揭示了一个被忽视的事实：银行软件的核心价值不在代码本身，而在系统集成、合规认证和历史数据积累。

JPM分析师在访问nCino总部后明确指出，客户对底层模型并不敏感，他们真正关心的是谁承担治理与监管责任。nCino联合创始人的观察更为直接——编码只是整个解决方案的一小部分，设计、测试、治理和持续支持才是真正的大头。

这意味着，即使AI能够生成代码，也无法自动完成与银行核心系统的对接、通过监管审计、处理遗留数据的兼容性问题。这些需要数年积累的行业知识和信任关系，才是数字银行SaaS最深的护城河。

KC评论： 很多投资者把AI时代的护城河等同于模型能力，但在银行这个行业，合规和信任才是真正的准入壁垒。完整报告中包含了对三家公司在合规深度上的详细对比，以及各家公司与银行核心系统的集成层级图，这些细节能帮你更精确地判断哪家公司的护城河最深。

在JPM的评级框架中，Q2 Holdings是最清晰的有机增长故事。公司所处的数字银行市场已经形成双寡头格局，Q2和Alkami共同主导，竞争烈度远低于市场想象。

几个关键数据点值得关注：Q2拥有约2500万数字银行账户的可触达市场，但只有约20\%迁移到了现代化平台。这意味着80\%的渗透空间，且这些客户多为区域银行和信用合作社，它们升级的紧迫性远高于大型银行。合同周期长达5-7年，每年仅400-500个续约节点，收入可见性极强。

管理层给出的FY27订阅增长指引为 \%，但JPM认为这仅仅是一个下限。更重要的是，公司预计未来五年EBITDA利润率将从FY25的24\%扩张至约35\%，且已实现超过90\%的自由现金流转化率。无负债、高转化、保守的并购态度——这些特征让Q2在基本面驱动的市场中具备溢价基础。

JPM给予FY27 15倍EV/EBITDA估值，对应目标价60美元，较当前股价有近40\%的上行空间。

KC评论： 15倍EV/EBITDA对于一家拥有双寡头地位、利润率持续扩张、且无负债的公司来说，看起来并不贵。但完整报告中包含了详细的利润率拆解模型和客户迁移时间表，这些才是判断估值是否合理的关键依据。

如果说Q2是稳健的有机增长标的，那么Alkami则是JPM眼中最具确定性的并购候选。激进投资者Jana Partners已经公开施压公司重启出售流程，这为股价提供了额外的催化剂。

Alkami的独特价值在于其云原生、多租户架构，这使其在部署速度和客户体验上具备优势。公司的客户留存率持续高于98\%，客户流失的主要原因不是竞争，而是客户自身的并购活动。这种粘性在软件行业中极为罕见。

谁会是潜在买家？JPM认为私募股权是最可能的收购方。原因在于：FIS、Fiserv、Jack Henry等战略买家拥有大量遗留系统，收购Alkami会面临生态

[... middle omitted ...]

估值有吸引力，但竞争格局和AI叙事仍不够清晰

nCino是三家覆盖公司中唯一被给予中性评级的。尽管JPM在访问其总部后对AI防御能力的态度有所改善，但相比Q2和Alkami，nCino面临两个核心问题。

一是终端市场不够集中。nCino主要服务商业银行业务，这个市场的碎片化程度远高于零售数字银行。客户分布在不同的银行规模、地域和业务重点中，这使得规模效应的释放更加困难。

二是AI叙事不如同行清晰。虽然nCino已经在Banking Advisor、Digital Partners等产品中展示了AI应用，但JPM认为市场上存在更多明确的AI受益标的，nCino的差异化仍需要时间证明。

不过，nCino的估值确实有吸引力。FY27 9倍EV/EBITDA，在垂直软件公司中处于低位。公司正处于定价模式转型期，同时削减研发支出并加大销售投入，如果这些策略成功，订阅收入有望在几个季度后加速增长。此外，与埃森哲、德勤等系统集成商的合作加深，也为利润率改善提供了另一条路径。

\subsection{[R044] 摩根斯坦利：新兴市场债券的“油价逻辑”尚未兑现}
{\small\color{refgray}归类：宏观与利率：美联储“沉默”重塑全球定价，日本央行内部分歧加大}

这份报告的核心判断简洁但有力：油价已经大幅回落，但新兴市场本地货币债券收益率并未同步下行。这中间存在一个10-15个基点的“便宜度”缺口，而摩根斯坦利认为这个缺口将在未来几周被填补。

这不是一份简单的“看多”报告。它更像是一份战术指南：在油价下跌、美联储暂停、风险偏好尚可的三重背景下，新兴市场债券的定价出现了结构性错位，而这种错位恰恰是机会所在。

报告发布的时间点是2026年6月22日，正值美国-伊朗冲突高峰过后、油价从高点显著回落之际。市场对新兴市场的情绪已经从三月的恐慌中恢复，但债券收益率的下行幅度远不及油价跌幅。摩根斯坦利的策略师团队——由James Lord领衔——认为，这种滞后不是基本面问题，而是由三个临时性因素造成的：美国利率走高、美元强势、以及市场对新兴市场央行加息的过度定价。

报告最核心的观察来自一个简单的相关性分析。在冲突爆发后的前两个半月，新兴市场本地货币债券指数收益率与油价走势高度同步。但过去一个月，两者出现了明显的背离：油价继续下行，而收益率却停滞不前。

用更宽泛的能源价格指数（Bloomberg Energy Index）来衡量，这种背离幅度约为10-15个基点。换句话说，按照历史关系，新兴市场债券收益率应该比现在的水平低10-15个基点。

KC评论： 10-15个基点听起来不大，但对于一个规模数万亿美元的资产类别而言，这意味着数十亿美元的定价偏差。更重要的是，这种偏差不是由基本面恶化引起的，而是由三个“临时性因素”造成的——这意味着它更可能被修复，而不是持久存在。

为什么收益率没有跟随油价下跌？摩根斯坦利给出了三个原因：美国国债收益率走高、美元走强、以及市场对新兴市场央行加息的定价。但关键在于，这三个因素都在边际上出现逆转信号。

摩根斯坦利的经济学家团队几乎在所有新兴市场国家都持有比市场更鸽派的利率预期。报告展示了一张对比图：市场定价显示新兴市场央行未来12个月将加息，而摩根斯坦利的基线预测是降息。

这种分歧的根源在于对通胀前景的判断不同。市场仍然在为高油价时期的通胀尾部风险定价，但摩根斯坦利认为，随着油价持续回落，通胀压力将显著缓解，央行将获得降息空间。

具体来看，南非市场定价隐含约75个基点的加息预期，而摩根斯坦利认为最多还有一次25个基点的加息。这种分歧为利率交易创造了机会：报告推荐做陡南非收益率曲线（1y1y-2y2y steepener），利用市场定价与摩根斯坦利基线之间的差距。

KC评论： 这里有一个值得追问的问题：摩根斯坦利的通胀模型是否过于乐观？如果油价回落的速度慢于预期，或者核心通胀的粘性超出想象，那么央行的鸽派空间就会被压缩。报告没有完全展开这个风险，但它在“Required reading”部分推荐了一篇关于南非宏观策略的文章，那里可能有更详细的风险分析。

新兴市场货币是这份报告的另一条主线。报告构建了一个新兴市场货币对G3（美元、欧元、日元）以及更广泛发达国家货币篮子（G3加瑞士法郎、澳元、加元）的总回报指数。这个指数在6月继续创出新高。

过去一个月、一个季度和年初至今，新兴市场货币对所有融资货币都实现了正回报——除了澳元（年初至今）和美元（本月至今，微幅负回报）。这个表现是在美联储6月会议偏鹰派的背景下取得的，这本身就说明新兴市场货币的上涨动力并非仅仅来自美元弱势。

报告认为，这种韧性的背后是三个结构性因素：风险偏好持续强劲、全球投资者对新兴市场本地资产的配置仍然偏低、以及新兴市场的基本面和政策环境在改善。

KC评论： 这里有一个有趣的悖论。报告说投资者对新兴市场的配置仍然偏低，但期权市场数据显示投资者已经建立了大量的HUF多头头寸，并且近期在加仓ZAR、CNH、IDR和INR。这意味着“低配”可能正在被快速修正，而一旦修正完成，

[... middle omitted ...]

——既没有极端流出，也没有极端流入。报告暗示，如果出现极端流出（目前还没有），那将是一个很强的风险回报入场点。

这份报告的一个关键假设是“美联储今年将按兵不动”。但报告本身也承认，美国经济数据强劲、FOMC偏鹰派，这导致了“美国例外论”的重新讨论。

问题在于：如果美国经济持续超预期，美联储被迫再次加息，或者至少维持高利率更长时间，那么新兴市场债券的“补跌”逻辑就会受到挑战。报告在讨论US收益率走高时提到了这一点，但没有深入分析这种情景的概率和影响。

对于读者来说，这是一个需要自己判断的关键变量。摩根斯坦利的基线是美联储暂停，但如果这个假设被打破，报告中推荐的几乎所有交易——做多新兴市场货币、做多新兴市场债券——都需要重新评估。

报告展示了一张年初至今的新兴市场总回报地图。拉丁美洲是唯一一个多个国家同时实现正回报的地区，其中哥伦比亚由于政策变革预期成为整个资产类别中表现最好的市场。匈牙利在中东欧地区领跑，南非也有所贡献。而亚洲地区普遍落后，只有中国提供了正回报。

\section{Disclaimer}
{\scriptsize\color{refgray}
This community is a paid learning and information-sharing community created for general educational purposes only. The membership fee is charged solely for access to community discussions, educational content, organization, and operational costs. It is not an advisory fee, management fee, performance fee, brokerage fee, or compensation for personalized financial, investment, legal, tax, or professional advice.\par
I participate and share information solely as an individual and for educational discussion among community members. I am not acting as a financial adviser, investment adviser, broker, dealer, portfolio manager, legal adviser, tax adviser, accountant, fiduciary, or any other licensed professional adviser. Nothing shared in this community constitutes, and should not be interpreted as, financial advice, investment advice, legal advice, tax advice, a recommendation, solicitation, offer, endorsement, instruction, or guarantee to buy, sell, hold, short, trade, or invest in any security, cryptocurrency, fund, derivative, strategy, or other financial product.\par
All content shared in this community, including messages, discussions, links, files, charts, examples, market commentary, watchlists, AI-generated or AI-polished summaries, and third-party materials, is general, impersonal, and educational in nature. It does not take into account any individual's financial situation, investment objectives, risk tolerance, investment horizon, liquidity needs, tax status, legal restrictions, or suitability requirements. No content should be relied upon as suitable for any specific person, account, portfolio, or financial decision.\par
Information shared in this community may be incomplete, inaccurate, outdated, speculative, AI-generated, AI-edited, or based on third-party internet sources that have not been independently verified. I make no representation or warranty as to the accuracy, completeness, timeliness, reliability, or suitability of any information shared.\par
Financial markets involve substantial risk, including the possible loss of principal. Past performance, historical data, hypothetical examples, simulated results, backtests, personal opinions, or educational examples do not guarantee future results. Each member is solely responsible for conducting their own research, exercising independent judgment, and making their own decisions. Before making any financial, investment, legal, or tax decision, members should consult qualified and licensed professionals.\par
Participation in this community, payment of any membership fee, or communication with me or other members does not create any adviser-client, fiduciary, professional, contractual, agency, partnership, employment, or other advisory relationship. By joining or participating in this community, you acknowledge and agree that you are solely responsible for your own decisions, actions, transactions, profits, losses, risks, and consequences, and that I shall not be liable for any direct, indirect, incidental, consequential, financial, legal, tax, or other loss or damage arising from or related to any content, discussion, information, or material shared in this community.\par
}
\end{document}
